\input{preamble}

% OK, start here.
%
\begin{document}

\title{Cohomologie des espaces algébriques}

\maketitle

\phantomsection
\label{section-phantom}

\tableofcontents




\section{Introduction}
\label{section-introduction}

\noindent
Dans ce chapitre, nous traitons de la cohomologie des espaces algébriques.
Bien que nous démontrions certains résultats sur la cohomologie des faisceaux abéliens,
nous nous intéressons principalement à la cohomologie des faisceaux quasi-cohérents, c'est-à-dire
que nous démontrons les analogues des résultats du chapitre ``Cohomologie des schémas''.
Certains résultats de ce chapitre se trouvent dans \cite{Kn}.

\medskip\noindent
Un ingrédient important qui manque dans ce chapitre est le
{\it principe d'induction}, c'est-à-dire l'analogue, pour les espaces algébriques quasi-compacts
et quasi-séparés, du
chapitre Cohomologie des schémas, lemme \ref{coherent-lemma-induction-principle}.
Ce principe est formulé précisément et démontré en détail dans
Catégories dérivées des espaces, section \ref{spaces-perfect-section-induction}.
Au lieu du principe d'induction, nous utilisons dans ce chapitre le
complexe de {\v C}ech alterné, voir
section \ref{section-alternating-cech}.
Il est conçu pour démontrer des énoncés d'annulation tels que
la proposition \ref{proposition-vanishing},
mais, dans certains cas, le principe d'induction est un outil plus puissant
et peut-être plus « standard ». Nous encourageons le lecteur
à consulter le principe d'induction
après avoir lu une partie du contenu de cette section.



\section{Conventions}
\label{section-conventions}

\noindent
L'hypothèse générale est que tous les schémas sont contenus dans
un grand site fppf $\Sch_{fppf}$. Tous les anneaux $A$ considérés
ont la propriété que $\Spec(A)$ est (isomorphe à) un
objet de ce grand site.

\medskip\noindent
Soit $S$ un schéma et soit $X$ un espace algébrique sur $S$.
Dans ce chapitre et le suivant, nous écrirons $X \times_S X$
pour le produit de $X$ par lui-même (dans la catégorie des espaces algébriques
sur $S$), plutôt que $X \times X$.








\section{Images directes supérieures}
\label{section-higher-direct-image}

\noindent
Soit $S$ un schéma. Soit $X$ un espace algébrique représentable sur $S$.
Soit $\mathcal{F}$ un module quasi-cohérent sur $X$ (voir
Propriétés des espaces, section \ref{spaces-properties-section-quasi-coherent}).
D'après
Descente, proposition \ref{descent-proposition-same-cohomology-quasi-coherent},
les groupes de cohomologie $H^i(X, \mathcal{F})$ coïncident avec le groupe de
cohomologie usuel calculé dans la topologie de Zariski du
module quasi-cohérent sur le schéma représentant $X$.

\medskip\noindent
Plus généralement, soit $f : X \to Y$ un morphisme quasi-compact et quasi-séparé
d'espaces algébriques représentables $X$ et $Y$. Soit
$\mathcal{F}$ un module quasi-cohérent sur $X$. D'après
Descente, lemme \ref{descent-lemma-higher-direct-images-small-etale},
le faisceau $R^if_*\mathcal{F}$ coïncide avec l'image
directe supérieure usuelle calculée pour la topologie de Zariski
du module quasi-cohérent sur le schéma représentant $X$
qui se déduit vers le schéma représentant $Y$.

\medskip\noindent
Plus généralement encore, supposons que $f : X \to Y$ soit un
morphisme représentable, quasi-compact et
quasi-séparé d'espaces algébriques sur $S$. Soit $V$ un schéma
et soit $V \to Y$ un morphisme étale surjectif. Posons $U = V \times_Y X$
et soit $f' : U \to V$ le changement de base de $f$. Pour tout
Pour tout $\mathcal{O}_X$-module quasi-cohérent $\mathcal{F}$, nous avons
\begin{equation}
\label{equation-representable-higher-direct-image}
R^if'_*(\mathcal{F}|_U) = (R^if_*\mathcal{F})|_V,
\end{equation}
voir
Propriétés des espaces,
lemme \ref{spaces-properties-lemma-pushforward-etale-base-change-modules}.
Comme $f' : U \to V$ est un morphisme quasi-compact et quasi-séparé
de schémas, la remarque du paragraphe précédent permet de
calculer $R^if'_*(\mathcal{F}|_U)$ en considérant $\mathcal{F}|_U$ comme un
faisceau quasi-cohérent sur le schéma $U$, et $f'$ comme un morphisme de schémas.
Nous utiliserons souvent ce fait sans autre mention.

\medskip\noindent
Nous démontrons ensuite que les images directes supérieures de faisceaux quasi-cohérents sont
quasi-cohérentes pour tout morphisme quasi-compact et quasi-séparé d'
espaces algébriques. Dans la démonstration, nous utilisons une astuce ; une démonstration « meilleure » utiliserait
un complexe de {\v C}ech relatif, comme cela est expliqué dans
Faisceaux sur les champs, sections \ref{stacks-sheaves-section-cech} et
\ref{stacks-sheaves-section-sheaf-cech-complex} ff.

\begin{lemma}
\label{lemma-higher-direct-image}
Soit $S$ un schéma. Soit $f : X \to Y$ un morphisme d'espaces algébriques
sur $S$. Si $f$ est quasi-compact et quasi-séparé, alors $R^if_*$
transforme les $\mathcal{O}_X$-modules quasi-cohérents en
$\mathcal{O}_Y$-modules quasi-cohérents.
\end{lemma}

\begin{proof}
Soit $V \to Y$ un morphisme étale où $V$ est un schéma affine. Posons
$U = V \times_Y X$ et notons $f' : U \to V$ le morphisme induit.
Soit $\mathcal{F}$ un $\mathcal{O}_X$-module quasi-cohérent. D'après
Propriétés des espaces, lemme
\ref{spaces-properties-lemma-pushforward-etale-base-change-modules}
nous avons
$R^if'_*(\mathcal{F}|_U) = (R^if_*\mathcal{F})|_V$.
Puisque la propriété d'être un module quasi-cohérent est locale pour la
topologie étale sur $Y$ (voir
Propriétés des espaces, lemme
\ref{spaces-properties-lemma-characterize-quasi-coherent})
nous pouvons remplacer $Y$ par $V$, c'est-à-dire supposer que $Y$ est un schéma affine.

\medskip\noindent
Supposons $Y$ affine. Comme $f$ est quasi-compact, $X$ est
quasi-compact. Nous pouvons donc choisir un schéma affine $U$ et un
morphisme étale surjectif $g : U \to X$, voir
Propriétés des espaces,
lemme \ref{spaces-properties-lemma-quasi-compact-affine-cover}.
Figure
$$
\xymatrix{
U \ar[r]_g \ar[rd]_{f \circ g} & X \ar[d]^f \\
& Y
}
$$
Le morphisme $g : U \to X$ est représentable, séparé
et quasi-compact puisque $X$ est quasi-séparé. Le lemme
s'applique donc à $g$ (par la discussion précédente).
Il s'applique également à $f \circ g : U \to Y$ (car il s'agit d'un morphisme
de schémas affines).

\medskip\noindent
Dans la situation décrite au paragraphe précédent, nous montrerons par
récurrence sur $n$ que $IH_n$ : pour tout faisceau quasi-cohérent $\mathcal{F}$
sur $X$, les faisceaux $R^if\mathcal{F}$
sont quasi-cohérents pour $i \leq n$.
Le cas $n = 0$ résulte de
Morphismes des espaces, lemme \ref{spaces-morphisms-lemma-pushforward}.
Supposons $IH_n$. Dans le reste de la démonstration, nous montrons que $IH_{n + 1}$ est valable.

\medskip\noindent
Soit $\mathcal{H}$ un $\mathcal{O}_U$-module quasi-cohérent.
Considérons la suite spectrale de Leray
$$
E_2^{p, q} = R^pf_* R^qg_* \mathcal{H}
\Rightarrow
R^{p + q}(f \circ g)_*\mathcal{H}
$$
Cohomologie sur les sites, lemme \ref{sites-cohomology-lemma-relative-Leray}.
Comme $R^qg_*\mathcal{H}$ est quasi-cohérent par $IH_n$, les faisceaux
$R^pf_*R^qg_*\mathcal{H}$ sont quasi-cohérents pour $p \leq n$.
Les faisceaux $R^{p + q}(f \circ g)_*\mathcal{H}$ sont tous
quasi-cohérents (en fait nuls pour $p + q > 0$, mais cela ne nous est pas nécessaire).
En degrés $\leq n + 1$, le seul module dont nous n'avons pas encore
établi la quasi-cohérence est $E_2^{n + 1, 0} = R^{n + 1}f_*g_*\mathcal{H}$.
De plus, les différentielles
$d_r^{n + 1, 0} : E_r^{n + 1, 0} \to E_r^{n + 1 + r, 1 - r}$
sont nulles puisque le terme cible est nul. Comme $\QCoh(\mathcal{O}_X)$
est une sous-catégorie de Serre faible de $\textit{Mod}(\mathcal{O}_X)$
(Propriétés des espaces, lemme
\ref{spaces-properties-lemma-properties-quasi-coherent})
il s'ensuit que $R^{n + 1}f_*g_*\mathcal{H}$
est quasi-cohérent (les détails sont omis).

\medskip\noindent
Soit $\mathcal{F}$ un $\mathcal{O}_X$-module quasi-cohérent.
Posons $\mathcal{H} = g^*\mathcal{F}$. Le morphisme d'adjonction
$\mathcal{F} \to g_*g^*\mathcal{F} = g_*\mathcal{H}$ est injectif
car $U \to X$ est étale surjectif. Considérons la suite exacte
$$
0 \to \mathcal{F} \to g_*\mathcal{H} \to \mathcal{G} \to 0
$$
où $\mathcal{G}$ est le conoyau de la première flèche et, en particulier,
quasi-cohérent. En appliquant la suite exacte longue de cohomologie, on obtient
$$
R^nf_*g_*\mathcal{H} \to
R^nf_*\mathcal{G} \to
R^{n + 1}f_*\mathcal{F} \to
R^{n + 1}f_*g_*\mathcal{H} \to
R^{n + 1}f_*\mathcal{G}
$$
Le conoyau de la première flèche est quasi-cohérent et
nous avons vu ci-dessus que $R^{n + 1}f_*g_*\mathcal{H}$ est quasi-cohérent.
Ainsi, $R^{n + 1}f_*\mathcal{F}$ possède une filtration à $2$ étapes,
dont le premier quotient est quasi-cohérent et le second un sous-module d'un
faisceau quasi-cohérent. Comme $\mathcal{F}$ est un module quasi-cohérent arbitraire
$\mathcal{O}_X$-module, le même résultat vaut pour $\mathcal{G}$.
Nous pouvons donc choisir une suite exacte
$0 \to \mathcal{A} \to R^{n + 1}f_*\mathcal{G} \to \mathcal{B}$
où $\mathcal{A}$ et $\mathcal{B}$ sont des $\mathcal{O}_Y$-modules quasi-cohérents.
Alors le noyau $\mathcal{K}$ de
$R^{n + 1}f_*g_*\mathcal{H} \to R^{n + 1}f_*\mathcal{G}
\to \mathcal{B}$ est quasi-cohérent ; on obtient donc un morphisme
$\mathcal{K} \to \mathcal{A}$ dont le noyau $\mathcal{K}'$ est
également quasi-cohérent. Ainsi $R^{n + 1}f_*\mathcal{F}$ s'insère dans une suite
exacte
$$
R^nf_*g_*\mathcal{H} \to
R^nf_*\mathcal{G} \to
R^{n + 1}f_*\mathcal{F} \to \mathcal{K}' \to 0
$$
où tous les modules sont quasi-cohérents, sauf éventuellement $R^{n + 1}f_*\mathcal{F}$.
Nous concluons que $R^{n + 1}f_*\mathcal{F}$ est quasi-cohérent, c'est-à-dire que
$IH_{n + 1}$ est vérifiée, comme souhaité.
\end{proof}

\begin{lemma}
\label{lemma-quasi-coherence-higher-direct-images-application}
Soit $S$ un schéma. Soit $f : X \to Y$ un morphisme quasi-séparé et quasi-compact
d'espaces algébriques sur $S$. Pour tout
$\mathcal{O}_X$-module quasi-cohérent $\mathcal{F}$ et tout objet affine $V$ de
$Y_\etale$, on a
$$
H^q(V \times_Y X, \mathcal{F}) = H^0(V, R^qf_*\mathcal{F})
$$
pour tout $q \in \mathbf{Z}$.
\end{lemma}

\begin{proof}
Comme la formation de $Rf_*$ commute à la localisation étale,
(Propriétés des espaces, lemme
\ref{spaces-properties-lemma-pushforward-etale-base-change-modules})
nous pouvons remplacer $Y$ par $V$ et supposer que $Y = V$ est affine.
Considérons la suite spectrale de Leray
$E_2^{p, q} = H^p(Y, R^qf_*\mathcal{F})$
convergeant vers $H^{p + q}(X, \mathcal{F})$, voir
Cohomologie sur les sites, lemme \ref{sites-cohomology-lemma-Leray}.
Par le lemme \ref{lemma-higher-direct-image},
les faisceaux $R^qf_*\mathcal{F}$ sont quasi-cohérents. D'après
Cohomologie des schémas, lemme
\ref{coherent-lemma-quasi-coherent-affine-cohomology-zero}
on a $E_2^{p, q} = 0$ lorsque $p > 0$.
La suite spectrale dégénère donc en $E_2$, ce qui conclut.
\end{proof}




\section{Morphismes finis}
\label{section-finite-morphisms}

\noindent
Voici quelques résultats valables pour tous les faisceaux abéliens
(en particulier pour les modules quasi-cohérents).
Nous {\bf avertissons} le lecteur que ces lemmes ne sont pas valables pour
les morphismes finis de schémas munis de la topologie de Zariski.

\begin{lemma}
\label{lemma-finite-higher-direct-image-zero}
Soit $S$ un schéma. Soit $f : X \to Y$ un morphisme intégral (par exemple fini)
d'espaces algébriques. Alors
$f_* : \textit{Ab}(X_\etale) \to \textit{Ab}(Y_\etale)$
est un foncteur exact et $R^pf_* = 0$ pour $p > 0$.
\end{lemma}

\begin{proof}
D'après Propriétés des espaces, lemme
\ref{spaces-properties-lemma-pushforward-etale-base-change}
nous pouvons calculer les images directes supérieures sur un recouvrement étale de $Y$.
Nous pouvons donc supposer que $Y$ est un schéma.
Cela implique que $X$ est un schéma (Morphismes des espaces, lemme
\ref{spaces-morphisms-lemma-integral-local}).
Dans ce cas, nous pouvons appliquer
Cohomologie étale, lemme \ref{etale-cohomology-lemma-what-integral}.
Dans le cas fini, le lecteur pourra consulter la
proposition moins technique de Cohomologie étale,
\ref{etale-cohomology-proposition-finite-higher-direct-image-zero}.
\end{proof}

\begin{lemma}
\label{lemma-stalk-push-finite}
Soit $S$ un schéma. Soit $f : X \to Y$ un morphisme fini d'espaces algébriques
sur $S$. Soit $\overline{y}$ un point géométrique de $Y$ dont les
relèvements notés $\overline{x}_1, \ldots, \overline{x}_n$ dans $X$. Alors
$$
(f_*\mathcal{F})_{\overline{y}} =
\prod\nolimits_{i = 1, \ldots, n}
\mathcal{F}_{\overline{x}_i}
$$
pour tout faisceau $\mathcal{F}$ sur $X_\etale$.
\end{lemma}

\begin{proof}
Choisissons un voisinage étale $(V, \overline{v})$ de $\overline{y}$.
Alors la fibre $(f_*\mathcal{F})_{\overline{y}}$
est la fibre de $f_*\mathcal{F}|_V$ en $\overline{v}$.
D'après Propriétés des espaces,
lemme \ref{spaces-properties-lemma-pushforward-etale-base-change}
nous pouvons remplacer $Y$ par $V$ et $X$ par $X \times_Y V$.
Alors $Z \to X$ est un morphisme fini de schémas et le résultat est
la proposition de Cohomologie étale
\ref{etale-cohomology-proposition-finite-higher-direct-image-zero}.
\end{proof}

\begin{lemma}
\label{lemma-finite-rings}
Soit $S$ un schéma. Soit $\pi : X \to Y$ un morphisme fini d'espaces
algébriques sur $S$. Soit $\mathcal{A}$ un faisceau d'anneaux sur $X_\etale$.
Soit $\mathcal{B}$ un faisceau d'anneaux sur $Y_\etale$.
Soit $\varphi : \mathcal{B} \to \pi_*\mathcal{A}$
un homomorphisme de faisceaux d'anneaux, de sorte que l'on obtient un
morphisme de topos annelés
$$
f = (\pi, \varphi) :
(\Sh(X_\etale), \mathcal{A})
\longrightarrow
(\Sh(Y_\etale), \mathcal{B}).
$$
Pour un faisceau de $\mathcal{A}$-modules $\mathcal{F}$ et un
faisceau de $\mathcal{B}$-modules $\mathcal{G}$, l'application canonique
$$
\mathcal{G} \otimes_\mathcal{B} f_*\mathcal{F}
\longrightarrow
f_*(f^*\mathcal{G} \otimes_\mathcal{A} \mathcal{F}).
$$
est un isomorphisme.
\end{lemma}

\begin{proof}
L'application est adjointe à l'application
$$
f^*\mathcal{G} \otimes_\mathcal{A}
f^* f_*\mathcal{F} =
f^*(\mathcal{G} \otimes_\mathcal{B} f_*\mathcal{F})
\longrightarrow
f^*\mathcal{G} \otimes_\mathcal{A} \mathcal{F}
$$
provenant de $\text{id} : f^*\mathcal{G} \to f^*\mathcal{G}$
et du morphisme d'adjonction $f^* f_*\mathcal{F} \to \mathcal{F}$.
Pour voir que cette application est un isomorphisme, nous pouvons vérifier sur les fibres
(Propriétés des espaces, théorème
\ref{spaces-properties-theorem-exactness-stalks}).
Soit $\overline{y}$ un point géométrique de $Y$ et
soient $\overline{x}_1, \ldots, \overline{x}_n$ les points géométriques
de $X$ situés au-dessus de $\overline{y}$.
En calculant l'effet de nos applications sur les fibres, nous voyons qu'il faut
démontrer
$$
\mathcal{G}_{\overline{y}}
\otimes_{\mathcal{B}_{\overline{y}}}
\left(
\bigoplus\nolimits_{i = 1, \ldots, n} \mathcal{F}_{\overline{x}_i}
\right) =
\bigoplus\nolimits_{i = 1, \ldots, n}
(\mathcal{G}_{\overline{y}}
\otimes_{\mathcal{B}_{\overline{x}}}
\mathcal{A}_{\overline{x}_i}) \otimes_{\mathcal{A}_{\overline{x}_i}}
\mathcal{F}_{\overline{x}_i}
$$
ce qui est vrai. Nous avons utilisé ici que
le produit tensoriel commute au passage aux fibres, ainsi que le
comportement des fibres par image inverse,
Propriétés des espaces, lemme \ref{spaces-properties-lemma-stalk-pullback}, et
le comportement des fibres sous poussée le long d'une immersion fermée, voir
lemme \ref{lemma-stalk-push-finite}.
\end{proof}

\noindent
Nous terminons cette section par une formule de projection très générale
pour les morphismes finis.

\begin{lemma}
\label{lemma-projection-formula-finite}
Avec $S$, $X$, $Y$, $\pi$, $\mathcal{A}$, $\mathcal{B}$, $\varphi$ et $f$
comme dans le lemme \ref{lemma-finite-rings}, on a
$$
K \otimes_\mathcal{B}^\mathbf{L} Rf_*M =
Rf_*(Lf^*K \otimes_\mathcal{A}^\mathbf{L} M)
$$
dans $D(\mathcal{B})$ pour tout $K \in D(\mathcal{B})$ et
$M \in D(\mathcal{A})$.
\end{lemma}

\begin{proof}
Comme $f_*$ est exact (lemme \ref{lemma-finite-higher-direct-image-zero}),
le foncteur $Rf_*$ se calcule en appliquant $f_*$ à un complexe représentant quelconque.
Choisissons un complexe $\mathcal{K}^\bullet$ de $\mathcal{B}$-modules
représentant $K$, K-plat à termes plats, voir
Cohomologie sur les sites, lemme \ref{sites-cohomology-lemma-K-flat-resolution}.
Alors $f^*\mathcal{K}^\bullet$ est K-plat à termes plats, voir
Cohomologie sur les sites, lemme \ref{sites-cohomology-lemma-pullback-K-flat}.
Choisissons un complexe quelconque $\mathcal{M}^\bullet$ de $\mathcal{A}$-modules
représentant $M$. Alors
il faut montrer
$$
\text{Tot}(\mathcal{K}^\bullet \otimes_\mathcal{B} f_*\mathcal{M}^\bullet)
=
f_*\text{Tot}(f^*\mathcal{K}^\bullet \otimes_\mathcal{A} \mathcal{M}^\bullet)
$$
car, par nos choix, ces complexes représentent respectivement les membres de droite et de gauche
de la formule du lemme.
Comme $f_*$ commute aux sommes directes
(par exemple d'après la description des fibres dans
le lemme \ref{lemma-stalk-push-finite}),
il suffit de vérifier les égalités
$$
\mathcal{K}^n \otimes_\mathcal{B} f_*\mathcal{M}^m
=
f_*(f^*\mathcal{K}^n \otimes_\mathcal{A} \mathcal{M}^m)
$$
qui résultent du lemme \ref{lemma-finite-rings}.
\end{proof}









\section{Colimites et cohomologie}
\label{section-colimits}

\noindent
Le lemme suivant s'applique en particulier aux diagrammes de
faisceaux quasi-cohérents.

\begin{lemma}
\label{lemma-colimits}
Soit $S$ un schéma. Soit $X$ un espace algébrique sur $S$.
Si $X$ est quasi-compact et quasi-séparé, alors
$$
\colim_i H^p(X, \mathcal{F}_i)
\longrightarrow
H^p(X, \colim_i \mathcal{F}_i)
$$
est un isomorphisme
pour tout diagramme filtrant de faisceaux abéliens sur $X_\etale$.
\end{lemma}

\begin{proof}
C'est une conséquence du lemme
Cohomologie sur les sites,
\ref{sites-cohomology-lemma-colim-works-over-collection}.
À savoir, soit $\mathcal{B} \subset \Ob(X_{spaces, \etale})$
l'ensemble des espaces quasi-compacts et quasi-séparés étales au-dessus de $X$.
Si $U \in \mathcal{B}$, comme $U$ est quasi-compact,
la famille des recouvrements finis $\{U_i \to U\}$ avec $U_i \in \mathcal{B}$
est cofinale dans l'ensemble des recouvrements de $U$ dans $X_{spaces, \etale}$. Par
le lemme de Morphismes des espaces
\ref{spaces-morphisms-lemma-quasi-compact-quasi-separated-permanence}
l'ensemble $\mathcal{B}$ vérifie toutes les hypothèses du
lemme de Cohomologie sur les sites
\ref{sites-cohomology-lemma-colim-works-over-collection}.
Comme $X \in \mathcal{B}$, le résultat s'ensuit.
\end{proof}

\begin{lemma}
\label{lemma-colimit-cohomology}
\begin{slogan}
Les images directes supérieures des morphismes qcqs commutent aux colimites filtrantes
de faisceaux.
\end{slogan}
Soit $S$ un schéma. Soit $f : X \to Y$ un morphisme quasi-compact et quasi-séparé
d'espaces algébriques sur $S$. Soit $\mathcal{F} = \colim \mathcal{F}_i$
une colimite filtrante de faisceaux abéliens sur $X_\etale$.
Alors, pour tout $p \geq 0$, on a
$$
R^pf_*\mathcal{F} = \colim R^pf_*\mathcal{F}_i.
$$
\end{lemma}

\begin{proof}
Nous utiliserons le fait que le morphisme de topos $f_{small} : X_{small} \to Y_{small}$
provient du morphisme de sites
$f_{spaces, \etale} : X_{spaces, \etale} \to Y_{spaces, \etale}$
associé au foncteur continu $V \longmapsto X \times_Y V$,
voir Propriétés des espaces, lemme
\ref{spaces-properties-lemma-functoriality-etale-site}.
Nous appliquerons le lemme
\ref{sites-cohomology-lemma-higher-direct-image-colimit}
à ce morphisme de sites. Comme tout objet de $Y_{spaces, \etale}$
admet un recouvrement par des objets affines, il suffit de montrer que
pour $V$ affine et étale au-dessus de $Y$, on a
$H^p(X \times_Y V, \mathcal{F}) = \colim H^p(X \times_Y V, \mathcal{F}_i)$.
Comme $V$ est affine, l'espace algébrique $X \times_Y V$ est
quasi-compact et quasi-séparé.
Nous pouvons donc appliquer le lemme \ref{lemma-colimits}
pour conclure.
\end{proof}

\noindent
Le lemme suivant montre que les modules de présentation finie se comportent
comme attendu sur les espaces algébriques quasi-compacts et quasi-séparés.

\begin{lemma}
\label{lemma-finite-presentation-quasi-compact-colimit}
Soit $S$ un schéma. Soit $X$ un espace algébrique quasi-compact et quasi-séparé
sur $S$. Soit $I$ un ensemble dirigé et
soit $(\mathcal{F}_i, \varphi_{ii'})$ un système indexé par $I$
de $\mathcal{O}_X$-modules. Soit $\mathcal{G}$ un
$\mathcal{O}_X$-module de présentation finie. Alors
$$
\colim_i \Hom_X(\mathcal{G}, \mathcal{F}_i)
=
\Hom_X(\mathcal{G}, \colim_i \mathcal{F}_i).
$$
En particulier, $\Hom_X(\mathcal{G}, -)$ commute aux colimites filtrantes
dans $\QCoh(\mathcal{O}_X)$.
\end{lemma}

\begin{proof}
L'égalité affichée est un cas particulier du lemme de
Modules sur les sites, lemme \ref{sites-modules-lemma-finite-presentation-quasi-compact-colimit}.
Pour l'appliquer, il faut vérifier les hypothèses du
lemme de Sites \ref{sites-lemma-directed-colimits-global-sections}, partie (4),
pour le site $X_\etale$.
Pour cela, nous vérifierons les hypothèses
(2)(a), (2)(b) et (2)(c) de
la remarque de Sites \ref{sites-remark-stronger-conditions}.
À savoir, soit $\mathcal{B} \subset \Ob(X_\etale)$ l'ensemble des objets affines.
Alors
\begin{enumerate}
\item Puisque $X$ est quasi-compact, il existe un $U \in \mathcal{B}$
tel que $U \to X$ soit surjectif
(Propriétés des espaces, lemme
\ref{spaces-properties-lemma-quasi-compact-affine-cover}),
d'où $h_U^\# \to *$ est surjectif.
\item Pour $U \in \mathcal{B}$, tout recouvrement étale $\{U_i \to U\}_{i \in I}$
de $U$ admet un raffinement par un recouvrement étale fini
$\{U_j \to U\}_{j = 1, \ldots, m}$ avec $U_j \in \mathcal{B}$
(Topologies, lemme \ref{topologies-lemma-etale-affine}).
\item Pour $U, U' \in \Ob(X_\etale)$, on a
$h_U^\# \times h_{U'}^\# = h_{U \times_X U'}^\#$.
Si $U, U' \in \mathcal{B}$, alors $U \times_X U'$ est quasi-compact
car $X$ est quasi-séparé ; voir par exemple le lemme de Morphismes des espaces
\ref{spaces-morphisms-lemma-quasi-compact-quasi-separated-permanence}
Par conséquent, il existe un morphisme étale surjectif
$U'' \to U \times_X U'$ avec $U'' \in \mathcal{B}$
Propriétés des espaces, lemme
\ref{spaces-properties-lemma-quasi-compact-affine-cover}).
Autrement dit, il existe des morphismes $U'' \to U$ et $U'' \to U'$
tels que l'application $h_{U''}^\# \to h_U^\# \times h_{u'}^\#$ soit
surjective.
\end{enumerate}
Pour l'énoncé final, observons que le foncteur d'inclusion
$\QCoh(\mathcal{O}_X) \to \textit{Mod}(\mathcal{O}_X)$
commute aux colimites et que les modules de présentation finie
sont quasi-cohérents. Voir
Propriétés des Espaces, lemme
\ref{spaces-properties-lemma-properties-quasi-coherent}.
\end{proof}





\section{Le complexe de {\v C}ech alterné}
\label{section-alternating-cech}

\noindent
Soit $S$ un schéma. Soit $f : U \to X$ un morphisme étale d'espaces
algébriques sur $S$. Le foncteur
$$
j : U_{spaces, \etale} \longrightarrow X_{spaces, \etale},\quad
V/U \longmapsto V/X
$$
induit une équivalence entre $U_{spaces, \etale}$ et la catégorie localisée
$X_{spaces, \etale}/U$, voir
Propriétés des espaces, section \ref{spaces-properties-section-localize}.
Il existe donc des foncteurs
$$
f_! : \textit{Ab}(U_\etale) \longrightarrow
\textit{Ab}(X_\etale),\quad
f_! : \textit{Mod}(\mathcal{O}_U) \longrightarrow \textit{Mod}(\mathcal{O}_X),
$$
qui sont adjoints à gauche des foncteurs
$$
f^{-1} : \textit{Ab}(X_\etale) \longrightarrow
\textit{Ab}(U_\etale),\quad
f^* : \textit{Mod}(\mathcal{O}_X) \longrightarrow \textit{Mod}(\mathcal{O}_U)
$$
voir
Modules sur les sites, section \ref{sites-modules-section-localize}.
Avertissement : a priori, ce foncteur n'a
rien à voir avec la cohomologie à supports compacts !
Dans la référence ci-dessus, nous avons appelé ce foncteur « extension par zéro ».
Notons que les deux versions de $f_!$ coïncident puisque $f^* = f^{-1}$ pour
les faisceaux de $\mathcal{O}_X$-modules.

\medskip\noindent
Comme nous utiliserons cette construction ci-dessous, rappelons quelques-unes de ses
propriétés. Soit $\mathcal{G}$ un faisceau abélien sur $U_\etale$ ;
le faisceau $f_!$ est le faisceautisé du préfaisceau
$$
V/X \longmapsto
f_!\mathcal{G}(V) =
\bigoplus\nolimits_{\varphi \in \Mor_X(V, U)}
\mathcal{G}(V \xrightarrow{\varphi} U),
$$
voir
Modules sur les sites, lemme \ref{sites-modules-lemma-extension-by-zero}.
De plus, si $\mathcal{G}$ est un $\mathcal{O}_U$-module, alors $f_!\mathcal{G}$
est le faisceautisé du même préfaisceau de groupes abéliens,
muni de façon évidente d'une structure de $\mathcal{O}_X$-module
(voir loc. cit.). Soit $\overline{x} : \Spec(k) \to X$
un point géométrique. Il existe alors une identification canonique
$$
(f_!\mathcal{G})_{\overline{x}} =
\bigoplus\nolimits_{\overline{u}} \mathcal{G}_{\overline{u}}
$$
où la somme porte sur tous les $\overline{u} : \Spec(k) \to U$ tels que
$f \circ \overline{u} = \overline{x}$, voir
Modules sur les sites, lemme \ref{sites-modules-lemma-stalk-j-shriek}
et
Propriétés des espaces, lemme
\ref{spaces-properties-lemma-points-small-etale-site}.
Dans la suite, nous étudions le faisceau $f_!\underline{\mathbf{Z}}$.
Ici, $\underline{\mathbf{Z}}$ désigne le faisceau constant sur
$X_\etale$ ou $U_\etale$.

\begin{lemma}
\label{lemma-product-is-tensor-product}
Soit $S$ un schéma. Soient $f_i : U_i \to X$ des morphismes étales
d'espaces algébriques sur $S$. Alors il existe des isomorphismes
$$
f_{1, !}\underline{\mathbf{Z}} \otimes_{\mathbf{Z}}
f_{2, !}\underline{\mathbf{Z}}
\longrightarrow
f_{12, !}\underline{\mathbf{Z}}
$$
où $f_{12} : U_1 \times_X U_2 \to X$ est le morphisme structural
et
$$
(f_1 \amalg f_2)_! \underline{\mathbf{Z}}
\longrightarrow
f_{1, !}\underline{\mathbf{Z}} \oplus
f_{2, !}\underline{\mathbf{Z}}
$$
\end{lemma}

\begin{proof}
Une fois l'application définie, ce sera un isomorphisme d'après la description
des fibres ci-dessus. Pour définir l'application, il suffit de travailler au niveau
des préfaisceaux. Nous devons donc définir une application
$$
\left(\bigoplus\nolimits_{\varphi_1 \in \Mor_X(V, U_1)} \mathbf{Z}\right)
\otimes_{\mathbf{Z}}
\left(\bigoplus\nolimits_{\varphi_2 \in \Mor_X(V, U_2)} \mathbf{Z}\right)
\longrightarrow
\bigoplus\nolimits_{\varphi \in \Mor_X(V, U_1 \times_X U_2)}
\mathbf{Z}
$$
Nous envoyons l'élément $1_{\varphi_1} \otimes 1_{\varphi_2}$ sur l'élément
$1_{\varphi_1 \times \varphi_2}$, avec les notations évidentes. Nous omettons la preuve
de la seconde égalité.
\end{proof}

\noindent
Une autre propriété importante est l'application trace
$$
\text{Tr}_f : f_!\underline{\mathbf{Z}} \longrightarrow \underline{\mathbf{Z}}.
$$
L'application trace est adjointe à
l'application $\mathbf{Z} \to f^{-1}\underline{\mathbf{Z}}$ (qui est un isomorphisme).
Si $\overline{x}$ est comme ci-dessus, l'application $\text{Tr}_f$ sur la fibre en $\overline{x}$
est l'application
$$
(\text{Tr}_f)_{\overline{x}} :
(f_!\underline{\mathbf{Z}})_{\overline{x}} =
\bigoplus\nolimits_{\overline{u}} \mathbf{Z}
\longrightarrow
\mathbf{Z} = \underline{\mathbf{Z}}_{\overline{x}}
$$
qui fait la somme des entiers donnés. Cela résulte de l'adjonction à l'application
$1 : \mathbf{Z} \to f^{-1}\underline{\mathbf{Z}}$. En particulier, si
$f$ est étale et surjectif, alors $\text{Tr}_f$ est surjectif.

\medskip\noindent
Supposons que $f : U \to X$ soit un morphisme étale surjectif
d'espaces algébriques. Considérons le {\it complexe de Koszul}
associé à l'application trace définie ci-dessus
$$
\ldots \to \wedge^3f_!\underline{\mathbf{Z}} \to
\wedge^2f_!\underline{\mathbf{Z}} \to f_!\underline{\mathbf{Z}} \to
\underline{\mathbf{Z}} \to 0
$$
Ici, les puissances extérieures sont prises sur le faisceau d'anneaux $\underline{\mathbf{Z}}$.
Les applications sont définies par la règle
$$
e_1 \wedge \ldots \wedge e_n \longmapsto
\sum\nolimits_{i = 1, \ldots, n} (-1)^{i + 1}
\text{Tr}_f(e_i)
e_1 \wedge \ldots \wedge \widehat{e_i} \wedge \ldots \wedge e_n
$$
où $e_1, \ldots, e_n$ sont des sections locales de $f_!\underline{\mathbf{Z}}$.
Soit $\overline{x}$ un point géométrique de $X$ et posons
$M_{\overline{x}} = (f_!\underline{\mathbf{Z}})_{\overline{x}} =
\bigoplus_{\overline{u}} \mathbf{Z}$. La fibre du complexe ci-dessus en
$\overline{x}$ est le complexe
$$
\ldots \to \wedge^3 M_{\overline{x}} \to \wedge^2 M_{\overline{x}}
\to M_{\overline{x}} \to \mathbf{Z} \to 0
$$
qui est exact puisque $M_{\overline{x}} \to \mathbf{Z}$ est surjectif, voir
Compléments d'algèbre, lemme \ref{more-algebra-lemma-homotopy-koszul-abstract}.
Ainsi, si $K^\bullet = K^\bullet(f)$ désigne le complexe défini par
$K^i = \wedge^{i + 1}f_!\underline{\mathbf{Z}}$, on obtient un
quasi-isomorphisme
\begin{equation}
\label{equation-quasi-isomorphism}
K^\bullet \longrightarrow \underline{\mathbf{Z}}[0]
\end{equation}
Nous utilisons le complexe $K^\bullet$ pour définir ce que nous appelons
le complexe de {\v C}ech alterné associé à $f : U \to X$.

\begin{definition}
\label{definition-alternating-cech-complex}
Soit $S$ un schéma. Soit $f : U \to X$ un morphisme étale surjectif
d'espaces algébriques sur $S$. Soit $\mathcal{F}$ un objet de
$\textit{Ab}(X_\etale)$. Le
{\it complexe de {\v C}ech alterné}\footnote{Cette notation peut être non standard}
$\check{\mathcal{C}}^\bullet_{alt}(f, \mathcal{F})$
associé à $\mathcal{F}$ et $f$ est le complexe
$$
\Hom(K^0, \mathcal{F}) \to \Hom(K^1, \mathcal{F}) \to
\Hom(K^2, \mathcal{F}) \to \ldots
$$
dont les groupes de morphismes sont calculés dans $\textit{Ab}(X_\etale)$.
\end{definition}

\noindent
Le lecteur peut vérifier que si $U = \coprod U_i$ et si $f|_{U_i} : U_i \to X$
est l'immersion ouverte d'un sous-espace, alors
$\check{\mathcal{C}}_{alt}^\bullet(f, \mathcal{F})$ coïncide avec le complexe
introduit dans
Cohomologie, section \ref{cohomology-section-alternating-cech}
pour le recouvrement de Zariski $X = \bigcup U_i$ et la restriction
de $\mathcal{F}$ au site de Zariski de $X$. Mais surtout,
il importe cependant de relier la cohomologie du
{\v C}ech alterné à la cohomologie.

\begin{lemma}
\label{lemma-alternating-cech-to-cohomology}
Soit $S$ un schéma. Soit $f : U \to X$ un morphisme étale surjectif
d'espaces algébriques sur $S$. Soit $\mathcal{F}$ un objet de
$\textit{Ab}(X_\etale)$. Il existe une application canonique
$$
\check{\mathcal{C}}^\bullet_{alt}(f, \mathcal{F})
\longrightarrow
R\Gamma(X, \mathcal{F})
$$
dans $D(\textit{Ab})$. De plus, il existe une suite spectrale dont la page $E_1$ est
$$
E_1^{p, q} =
\Ext_{\textit{Ab}(X_\etale)}^q(K^p, \mathcal{F})
$$
convergeant vers $H^{p + q}(X, \mathcal{F})$, où
$K^p = \wedge^{p + 1}f_!\underline{\mathbf{Z}}$.
\end{lemma}

\begin{proof}
Rappelons le quasi-isomorphisme
$K^\bullet \to \underline{\mathbf{Z}}[0]$, voir
(\ref{equation-quasi-isomorphism}).
Choisissons une résolution injective $\mathcal{F} \to \mathcal{I}^\bullet$
dans $\textit{Ab}(X_\etale)$. Considérons le bicomplexe
$\Hom(K^\bullet, \mathcal{I}^\bullet)$ dont les termes sont
$\Hom(K^p, \mathcal{I}^q)$. La différentielle
$d_1^{p, q} : A^{p, q} \to A^{p + 1, q}$
est induite par la différentielle $K^{p + 1} \to K^p$
et la différentielle $d_2^{p, q} : A^{p, q} \to A^{p, q + 1}$ est induite par
la différentielle
$\mathcal{I}^q \to \mathcal{I}^{q + 1}$.
Notons $\text{Tot}(\Hom(K^\bullet, \mathcal{I}^\bullet))$
le complexe total associé, voir
Homologie, section \ref{homology-section-double-complexes}.
Nous utiliserons les deux suites spectrales
$({}'E_r, {}'d_r)$ et $({}''E_r, {}''d_r)$
associées à ce bicomplexe, voir
Homologie, section \ref{homology-section-double-complex}.

\medskip\noindent
Comme $K^\bullet$ est une résolution de $\underline{\mathbf{Z}}$,
on voit que les complexes
$$
\Hom(K^\bullet, \mathcal{I}^q) :
\Hom(K^0, \mathcal{I}^q) \to
\Hom(K^1, \mathcal{I}^q) \to
\Hom(K^2, \mathcal{I}^q) \to \ldots
$$
sont acycliques en degrés strictement positifs et ont pour $H^0$
$\Gamma(X, \mathcal{I}^q)$. Par conséquent, d'après
Homologie, lemme \ref{homology-lemma-double-complex-gives-resolution}
l'application naturelle
$$
\mathcal{I}^\bullet(X) \longrightarrow
\text{Tot}(\Hom(K^\bullet, \mathcal{I}^\bullet))
$$
est un quasi-isomorphisme de complexes de groupes abéliens.
En particulier, on conclut que
$H^n(\text{Tot}(\Hom(K^\bullet, \mathcal{I}^\bullet))) = H^n(X, \mathcal{F})$.

\medskip\noindent
L'application $\check{\mathcal{C}}^\bullet_{alt}(f, \mathcal{F}) \to
R\Gamma(X, \mathcal{F})$ du lemme est la composée de
$\check{\mathcal{C}}^\bullet_{alt}(f, \mathcal{F}) \to
\text{Tot}(\Hom(K^\bullet, \mathcal{I}^\bullet))$
avec l'inverse du quasi-isomorphisme affiché.

\medskip\noindent
Considérons enfin la suite spectrale $({}'E_r, {}'d_r)$.
Nous avons
$$
E_1^{p, q} = q\text{-ième cohomologie de }
\Hom(K^p, \mathcal{I}^0) \to
\Hom(K^p, \mathcal{I}^1) \to
\Hom(K^p, \mathcal{I}^2) \to \ldots
$$
Ceci démontre le lemme.
\end{proof}

\noindent
Il résulte du lemme qu'il est important de comprendre
les groupes d'extensions $\Ext_{\textit{Ab}(X_\etale)}(K^p, \mathcal{F})$,
c'est-à-dire les foncteurs dérivés à droite de
$\mathcal{F} \mapsto \Hom(K^p, \mathcal{F})$.

\begin{lemma}
\label{lemma-compute}
Soit $S$ un schéma. Soit $f : U \to X$ un morphisme étale, surjectif et séparé
d'espaces algébriques sur $S$. Pour $p \geq 0$, posons
$$
W_p = U \times_X \ldots \times_X U \setminus \text{toutes les diagonales}
$$
où le produit fibré comporte $p + 1$ facteurs.
Le groupe $S_{p + 1}$ agit librement sur $W_p$ au-dessus de $X$ et
$$
\Hom(K^p, \mathcal{F}) = S_{p + 1}\text{-anti-invariants de }
\mathcal{F}(W_p)
$$
fonctoriellement en $\mathcal{F}$, où
$K^p = \wedge^{p + 1}f_!\underline{\mathbf{Z}}$.
\end{lemma}

\begin{proof}
Comme $U \to X$ est séparé, la diagonale $U \to U \times_X U$ est une
immersion fermée. Comme $U \to X$ est étale, cette diagonale
$U \to U \times_X U$ est une immersion ouverte, voir
Morphismes des espaces, lemmes
\ref{spaces-morphisms-lemma-etale-unramified} et
\ref{spaces-morphisms-lemma-diagonal-unramified-morphism}.
Ainsi $W_p$ est un sous-espace à la fois ouvert et fermé de
$U^{p + 1} = U \times_X \ldots \times_X U$. L'action de $S_{p + 1}$
sur $W_p$ est libre, car nous avons supprimé les points fixes de l'action.
D'après
lemme \ref{lemma-product-is-tensor-product}
on obtient
$$
(f_!\underline{\mathbf{Z}})^{\otimes p + 1} =
f^{p + 1}_!\underline{\mathbf{Z}} = (W_p \to X)_!\underline{\mathbf{Z}}
\oplus Rest
$$
où $f^{p + 1} : U^{p + 1} \to X$ est le morphisme structural.
En examinant les fibres au-dessus d'un point géométrique $\overline{x}$ de $X$,
on voit que
$$
\left(
\bigoplus\nolimits_{\overline{u} \mapsto \overline{x}} \mathbf{Z}
\right)^{\otimes p + 1}
\longrightarrow
(W_p \to X)_!\underline{\mathbf{Z}}_{\overline{x}}
$$
est le quotient dont le noyau est engendré par tous les tenseurs
$1_{\overline{u}_0} \otimes \ldots \otimes 1_{\overline{u}_p}$
où $\overline{u}_i = \overline{u}_j$ pour certains $i \not = j$.
Ainsi, l'application quotient
$$
(f_!\underline{\mathbf{Z}})^{\otimes p + 1}
\longrightarrow
\wedge^{p + 1}f_!\underline{\mathbf{Z}}
$$
se factorise par $(W_p \to X)_!\underline{\mathbf{Z}}$ ; on obtient donc
$$
(f_!\underline{\mathbf{Z}})^{\otimes p + 1}
\longrightarrow
(W_p \to X)_!\underline{\mathbf{Z}}
\longrightarrow
\wedge^{p + 1}f_!\underline{\mathbf{Z}}
$$
Cela montre déjà que $\Hom(K^p, \mathcal{F})$ est (fonctoriellement) un
sous-groupe de
$$
\Hom((W_p \to X)_!\underline{\mathbf{Z}}, \mathcal{F}) = \mathcal{F}(W_p)
$$
Pour l'identifier aux éléments anti-invariants sous $S_{p + 1}$, il faut montrer que
la surjection $(W_p \to X)_!\underline{\mathbf{Z}}
\to \wedge^{p + 1}f_!\underline{\mathbf{Z}}$ est le quotient anti-invariant maximal
$S_{p + 1}$-quotient anti-invariant. Autrement dit, il faut montrer que
$\wedge^{p + 1}f_!\underline{\mathbf{Z}}$ est le quotient de
$(W_p \to X)_!\underline{\mathbf{Z}}$ par le sous-faisceau engendré par
les sections locales $s - \text{signature}(\sigma)\sigma(s)$, où $s$ est
une section locale de $(W_p \to X)_!\underline{\mathbf{Z}}$.
Cela se vérifie sur les fibres, où c'est clair.
\end{proof}

\begin{lemma}
\label{lemma-twist}
Soit $S$ un schéma. Soit $W$ un espace algébrique sur $S$.
Soit $G$ un groupe fini agissant librement sur $W$.
Soit $U = W/G$, voir
Propriétés des espaces, lemme \ref{spaces-properties-lemma-quotient}.
Soit $\chi : G \to \{+1, -1\}$ un caractère.
Alors il existe un faisceau de $\mathbf{Z}$-modules localement libre de rang 1
$\underline{\mathbf{Z}}(\chi)$ sur $U_\etale$ tel que, pour tout
faisceau abélien $\mathcal{F}$ sur $U_\etale$, on ait
$$
H^0(W, \mathcal{F}|_W)^\chi =
H^0(U, \mathcal{F} \otimes_{\mathbf{Z}} \underline{\mathbf{Z}}(\chi))
$$
\end{lemma}

\begin{proof}
Le morphisme quotient $q : W \to U$ est un $G$-torseur ; autrement dit, il existe
un morphisme étale surjectif $U' \to U$ tel que
$W \times_U U' = \coprod_{g \in G} U'$ comme espaces munis de l'action de $G$ au-dessus de $U'$.
(On peut notamment prendre $U' = W$.) Ainsi, $q_*\underline{\mathbf{Z}}$ est un
module de $\mathbf{Z}$ localement libre de rang fini muni d'une action de $G$.
Pour tout point géométrique $\overline{u}$ de $U$, on obtient des $G$-équivariants
isomorphismes
$$
(q_*\underline{\mathbf{Z}})_{\overline{u}}
= \bigoplus\nolimits_{\overline{w} \mapsto \overline{u}} \mathbf{Z}
= \bigoplus\nolimits_{g \in G} \mathbf{Z} = \mathbf{Z}[G]
$$
où le second signe $=$ utilise un point géométrique
$\overline{w}_0$ au-dessus de $\overline{u}$ et
envoie le facteur correspondant à $g \in G$ sur le facteur
correspondant à $g(\overline{w}_0)$. Nous avons
$$
H^0(W, \mathcal{F}|_W) =
H^0(U, \mathcal{F} \otimes_\mathbf{Z} q_*\underline{\mathbf{Z}})
$$
car
$q_*\mathcal{F}|_W = \mathcal{F} \otimes_\mathbf{Z} q_*\underline{\mathbf{Z}}$
comme on peut le vérifier après restriction à $U'$.
$$
\underline{\mathbf{Z}}(\chi) =
(q_*\underline{\mathbf{Z}})^\chi \subset
q_*\underline{\mathbf{Z}}
$$
qui est le sous-faisceau des sections se transformant selon $\chi$.
Pour tout point géométrique $\overline{u}$ de $U$, on a
$$
\underline{\mathbf{Z}}(\chi)_{\overline{u}} =
\mathbf{Z} \cdot \sum\nolimits_g \chi(g) g
\subset
\mathbf{Z}[G] = (q_*\underline{\mathbf{Z}})_{\overline{u}}
$$
Il s'ensuit que $\underline{\mathbf{Z}}(\chi)$ est localement libre de
rang 1 (plus précisément, cela doit être vérifié après restriction à $U'$).
Notons que, pour tout $\mathbf{Z}$-module $M$, les $\chi$-semi-invariants
de $M[G]$ sont les éléments de la forme $m \cdot \sum\nolimits_g \chi(g) g$.
Ainsi, pour tout faisceau abélien $\mathcal{F}$ sur $U$, on a
$$
\left(\mathcal{F} \otimes_\mathbf{Z} q_*\underline{\mathbf{Z}}\right)^\chi
=
\mathcal{F} \otimes_\mathbf{Z} \underline{\mathbf{Z}}(\chi)
$$
car l'égalité se vérifie sur toutes les fibres. Le résultat du lemme s'obtient en
prenant les sections globales.
\end{proof}

\noindent
Nous pouvons maintenant tout rassembler et obtenir le résultat suivant
particulièrement agréable.

\begin{lemma}
\label{lemma-alternating-spectral-sequence}
Soit $S$ un schéma. Soit $f : U \to X$ un morphisme d'espaces algébriques
sur $S$, surjectif, \'etale et séparé. Pour $p \geq 0$, posons
$$
W_p = U \times_X \ldots \times_X U \setminus \text{toutes les diagonales}
$$
(avec $p + 1$ facteurs), comme dans le lemme \ref{lemma-compute}.
Soit $\chi_p : S_{p + 1} \to \{+1, -1\}$ le caractère signe.
Soient $U_p = W_p/S_{p + 1}$ et $\underline{\mathbf{Z}}(\chi_p)$ les objets définis comme dans le
lemme \ref{lemma-twist}.
Alors la suite spectrale du
lemme \ref{lemma-alternating-cech-to-cohomology}
a pour page $E_1$
$$
E_1^{p, q} =
H^q(U_p, \mathcal{F}|_{U_p} \otimes_\mathbf{Z} \underline{\mathbf{Z}}(\chi_p))
$$
et aboutit à $H^{p + q}(X, \mathcal{F})$.
\end{lemma}

\begin{proof}
Remarquons que, puisque l'action de $S_{p + 1}$ sur $W_p$ est au-dessus de $X$, on
obtient bien un morphisme $U_p \to X$. Puisque $W_p \to X$ est \'etale et que
$W_p \to U_p$ est \'etale surjectif, il s'ensuit
que $U_p \to X$ est lui aussi \'etale ; voir
Morphismes d'espaces, lemme \ref{spaces-morphisms-lemma-etale-local}.
Par conséquent, un objet injectif de
$\textit{Ab}(X_\etale)$ devient par restriction un objet injectif de
$\textit{Ab}(U_{p, \etale})$ ; voir
Cohomologie sur les sites, lemme \ref{sites-cohomology-lemma-cohomology-of-open}.
De plus, le foncteur
$\mathcal{G} \mapsto
\mathcal{G} \otimes_\mathbf{Z} \underline{\mathbf{Z}}(\chi_p))$
est une auto-équivalence de $\textit{Ab}(U_p)$, donc transforme les objets injectifs
en objets injectifs et est exact (car
$\underline{\mathbf{Z}}(\chi_p)$ est un
$\underline{\mathbf{Z}}$-module inversible). Ainsi, étant donnée une résolution injective
$\mathcal{F} \to \mathcal{I}^\bullet$ dans $\textit{Ab}(X_\etale)$,
le complexe
$$
\Gamma(U_p,
\mathcal{I}^0|_{U_p} \otimes_\mathbf{Z} \underline{\mathbf{Z}}(\chi_p))
\to
\Gamma(U_p,
\mathcal{I}^1|_{U_p} \otimes_\mathbf{Z} \underline{\mathbf{Z}}(\chi_p))
\to
\Gamma(U_p,
\mathcal{I}^2|_{U_p} \otimes_\mathbf{Z} \underline{\mathbf{Z}}(\chi_p))
\to \ldots
$$
calcule
$H^*(U_p,
\mathcal{F}|_{U_p} \otimes_\mathbf{Z} \underline{\mathbf{Z}}(\chi_p))$.
D'autre part, d'après
le lemme \ref{lemma-twist},
il est égal au complexe des éléments anti-invariants sous $S_{p + 1}$ de
$$
\Gamma(W_p, \mathcal{I}^0) \to
\Gamma(W_p, \mathcal{I}^1) \to
\Gamma(W_p, \mathcal{I}^2) \to \ldots
$$
qui, d'après
le lemme \ref{lemma-compute},
est égal au complexe
$$
\Hom(K^p, \mathcal{I}^0) \to
\Hom(K^p, \mathcal{I}^1) \to
\Hom(K^p, \mathcal{I}^2) \to \ldots
$$
qui calcule
$\Ext^*_{\textit{Ab}(X_\etale)}(K^p, \mathcal{F})$.
En combinant ce qui précède, on conclut.
\end{proof}





\section{Annulation en degrés supérieurs pour les faisceaux quasi-cohérents}
\label{section-higher-vanishing}

\noindent
Dans cette section, nous montrons que, pour tout espace algébrique quasi-compact et
quasi-séparé $X$, il existe un entier
$n = n(X)$ tel que la cohomologie de tout faisceau quasi-cohérent
sur $X$ est nulle en degrés strictement supérieurs à $n$.

\begin{lemma}
\label{lemma-quasi-coherent-twist}
Soient $S$, $W$, $G$, $U$ et $\chi$ comme dans le
lemme \ref{lemma-twist}.
Si $\mathcal{F}$ est un $\mathcal{O}_U$-module quasi-cohérent,
alors $\mathcal{F} \otimes_{\mathbf{Z}} \underline{\mathbf{Z}}(\chi)$ l'est aussi.
\end{lemma}

\begin{proof}
La structure de $\mathcal{O}_U$-module est claire. Pour vérifier que
$\mathcal{F} \otimes_{\mathbf{Z}} \underline{\mathbf{Z}}(\chi)$
est quasi-cohérent, il suffit de le vérifier localement pour la topologie \'etale.
Comme $\underline{\mathbf{Z}}(\chi)$
est un $\underline{\mathbf{Z}}$-module localement libre de rang fini, on en déduit le lemme.
\end{proof}

\noindent
La proposition suivante est intéressante même lorsque $X$ est un schéma.
Elle constitue la généralisation naturelle du
lemme \ref{coherent-lemma-vanishing-nr-affines} du chapitre Cohomologie des schémas.
Avant de l'énoncer, remarquons que, pour tout morphisme \'etale
$f : U \to X$ d'un schéma affine vers un espace algébrique
quasi-séparé $X$, les fibres de $f$ sont universellement bornées ; en particulier,
il existe un entier $d$ tel que toutes les fibres de $|U| \to |X|$
soient de cardinal au plus $d$ ; c'est l'implication
$(\eta) \Rightarrow (\delta)$ du lemme
\ref{decent-spaces-lemma-bounded-fibres} du chapitre Espaces décents.

\begin{proposition}
\label{proposition-vanishing}
Soit $S$ un schéma. Soit $X$ un espace algébrique sur $S$.
Supposons $X$ quasi-compact et séparé.
Soit $U$ un schéma affine, et soit
$f : U \to X$ un morphisme \'etale surjectif.
Soit $d$ un majorant du cardinal des fibres de
$|U| \to |X|$. Alors, pour tout $\mathcal{O}_X$-module quasi-cohérent $\mathcal{F}$,
on a $H^q(X, \mathcal{F}) = 0$ pour tout $q \geq d$.
\end{proposition}

\begin{proof}
Nous utiliserons la suite spectrale du
lemme \ref{lemma-alternating-spectral-sequence}.
Le lemme s'applique, car $f$ est séparé puisque $U$ est séparé ; voir
Morphismes d'espaces, lemme
\ref{spaces-morphisms-lemma-compose-after-separated}.
Puisque $X$ est séparé, le schéma $U \times_X \ldots \times_X U$ est un sous-schéma fermé
de
$U \times_{\Spec(\mathbf{Z})} \ldots \times_{\Spec(\mathbf{Z})} U$
et est donc affine. Ainsi, $W_p$ est affine. Par conséquent, $U_p = W_p/S_{p + 1}$ est un
schéma affine d'après
Groupoïdes, proposition \ref{groupoids-proposition-finite-flat-equivalence}.
La discussion de la
section \ref{section-higher-direct-image}
montre que la cohomologie des faisceaux quasi-cohérents sur $W_p$ (considéré comme espace
algébrique) coïncide avec la cohomologie du faisceau quasi-cohérent correspondant
sur le schéma affine sous-jacent ; elle est donc nulle en degrés positifs d'après
Cohomologie des schémas, lemme
\ref{coherent-lemma-quasi-coherent-affine-cohomology-zero}.
D'après
le lemme \ref{lemma-quasi-coherent-twist},
les faisceaux
$\mathcal{F}|_{U_p} \otimes_\mathbf{Z} \underline{\mathbf{Z}}(\chi_p)$
sont quasi-cohérents. Ainsi,
$H^q(W_p,
\mathcal{F}|_{U_p} \otimes_\mathbf{Z} \underline{\mathbf{Z}}(\chi_p))$
est nul pour $q > 0$. Par définition de l'entier $d$, on voit que
$W_p = \emptyset$ pour $p \geq d$. Par conséquent,
$H^0(W_p,
\mathcal{F}|_{U_p} \otimes_\mathbf{Z} \underline{\mathbf{Z}}(\chi_p))$
est également nul pour $p \geq d$.
Cela démontre la proposition.
\end{proof}

\noindent
Dans le lemme suivant, nous établissons qu'un espace algébrique quasi-compact et
quasi-séparé est de dimension cohomologique finie
pour les modules quasi-cohérents. Nous ne précisons la borne que parce que
nous l'utiliserons plus loin pour démontrer un résultat analogue pour les images directes
supérieures.

\begin{lemma}
\label{lemma-vanishing-quasi-separated}
Soit $S$ un schéma. Soit $X$ un espace algébrique sur $S$.
Supposons $X$ quasi-compact et quasi-séparé.
On peut alors choisir
\begin{enumerate}
\item un schéma affine $U$,
\item un morphisme \'etale surjectif $f : U \to X$,
\item un entier $d$ majorant les degrés des fibres de $U \to X$,
\item pour tout $p = 0, 1, \ldots, d$, un morphisme \'etale surjectif
$V_p \to U_p$, où $V_p$ est un schéma affine et $U_p$ est défini comme dans le
lemme \ref{lemma-alternating-spectral-sequence}, et
\item un entier $d_p$ majorant le degré des fibres de $V_p \to U_p$.
\end{enumerate}
De plus, chaque fois que l'on dispose de (1) -- (5), pour tout
$\mathcal{O}_X$-module quasi-cohérent $\mathcal{F}$, on a $H^q(X, \mathcal{F}) = 0$ pour tout
$q \geq \max(d_p + p)$.
\end{lemma}

\begin{proof}
Puisque $X$ est quasi-compact, on peut trouver un morphisme \'etale surjectif
$U \to X$ avec $U$ affine ; voir
Propriétés d'espaces, lemme
\ref{spaces-properties-lemma-quasi-compact-affine-cover}.
D'après
Espaces décents, lemme \ref{decent-spaces-lemma-bounded-fibres},
les fibres de $f$ sont universellement bornées ; on peut donc trouver $d$.
On a $U_p = W_p/S_{p + 1}$ et $W_p \subset U \times_X \ldots \times_X U$
est ouvert et fermé. Puisque $X$ est quasi-séparé, les schémas
$W_p$ sont quasi-compacts, donc $U_p$ est quasi-compact.
Puisque $U$ est séparé, les schémas $W_p$ sont séparés ; ainsi,
$U_p$ est séparé d'après (la version absolue du)
Espaces, lemme \ref{spaces-lemma-quotient-finite-separated}.
D'après
Propriétés d'espaces, lemme
\ref{spaces-properties-lemma-quasi-compact-affine-cover},
on peut trouver les morphismes $V_p \to W_p$.
D'après
Espaces décents, lemme \ref{decent-spaces-lemma-bounded-fibres},
on peut trouver les entiers $d_p$.

\medskip\noindent
On utilise alors la suite spectrale
$$
E_1^{p, q} =
H^q(U_p, \mathcal{F}|_{U_p} \otimes_\mathbf{Z} \underline{\mathbf{Z}}(\chi_p))
\Rightarrow
H^{p + q}(X, \mathcal{F})
$$
voir
le lemme \ref{lemma-alternating-spectral-sequence}.
Par définition de l'entier $d$, on voit que $U_p = 0$ pour $p \geq d$.
D'après la proposition \ref{proposition-vanishing}
et
le lemme \ref{lemma-quasi-coherent-twist},
on voit que
$H^q(U_p,
\mathcal{F}|_{U_p} \otimes_\mathbf{Z} \underline{\mathbf{Z}}(\chi_p))$
est nul pour $q \geq d_p$ lorsque $p = 0, \ldots, d$.
On en déduit le lemme.
\end{proof}









\section{Annulation des images directes supérieures}
\label{section-vanishing-higher-direct-images}

\noindent
On applique les résultats de la
section \ref{section-higher-vanishing}
afin d'obtenir l'annulation des images directes supérieures des faisceaux quasi-cohérents
pour les morphismes quasi-compacts et quasi-séparés. Ce résultat est utile, car
il permet, dans certaines situations, de raisonner par récurrence descendante sur le degré
cohomologique.

\begin{lemma}
\label{lemma-vanishing-higher-direct-images}
Soit $S$ un schéma. Soit $f : X \to Y$ un
morphisme d'espaces algébriques sur $S$.
Supposons que
\begin{enumerate}
\item $f$ soit quasi-compact et quasi-séparé, et
\item $Y$ soit quasi-compact.
\end{enumerate}
Alors il existe un entier $n(X \to Y)$ tel que,
pour tout espace algébrique $Y'$, tout morphisme $Y' \to Y$
et tout faisceau quasi-cohérent $\mathcal{F}'$ sur $X' = Y' \times_Y X$,
les images directes supérieures $R^if'_*\mathcal{F}'$ soient nulles pour tout
$i \geq n(X \to Y)$.
\end{lemma}

\begin{proof}
Soit $V \to Y$ un morphisme \'etale surjectif, où $V$ est un schéma
affine ; voir
Propriétés d'espaces, lemme
\ref{spaces-properties-lemma-quasi-compact-affine-cover}.
Supposons le résultat prouvé pour le changement de base $f_V : V \times_Y X \to V$.
Le résultat vaut alors pour $f$ avec $n(X \to Y) = n(X_V \to V)$.
En effet, si $Y' \to Y$ et $\mathcal{F}'$ sont comme dans le lemme, alors
$R^if'_*\mathcal{F}'|_{V \times_Y Y'}$ est égal à
$R^if'_{V, *}\mathcal{F}'|_{X'_V}$, où
$f'_V : X'_V = V \times_Y Y' \times_Y X \to V \times_Y Y' = Y'_V$ ; voir
Propriétés d'espaces,
lemme \ref{spaces-properties-lemma-pushforward-etale-base-change-modules}.
On peut donc supposer que $Y$ est un schéma affine.

\medskip\noindent
De plus, pour démontrer l'annulation pour tout $Y' \to Y$ et toute
$\mathcal{F}'$, il suffit de le faire lorsque $Y'$ est un schéma affine.
Dans ce cas, $R^if'_*\mathcal{F}'$ est quasi-cohérent d'après
le lemme \ref{lemma-higher-direct-image}.
Il suffit donc de démontrer que $H^i(X', \mathcal{F}') = 0$, car
$H^i(X', \mathcal{F}') = H^0(Y', R^if'_*\mathcal{F}')$ d'après
Cohomologie sur les sites, lemme \ref{sites-cohomology-lemma-apply-Leray}
et le fait que la cohomologie des faisceaux quasi-cohérents
sur les espaces algébriques affines est nulle en degrés supérieurs
(proposition \ref{proposition-vanishing}).

\medskip\noindent
Choisissons $U \to X$, $d$, $V_p \to U_p$ et $d_p$ comme dans
le lemme \ref{lemma-vanishing-quasi-separated}.
Pour tout schéma affine $Y'$ et tout morphisme $Y' \to Y$, notons
$X' = Y' \times_Y X$, $U' = Y' \times_Y U$, $V'_p = Y' \times_Y V_p$.
Alors $U' \to X'$, $d' = d$, $V'_p \to U'_p$ et $d'_p = d$
forment un ensemble de choix comme dans
le lemme \ref{lemma-vanishing-quasi-separated}
pour l'espace algébrique $X'$ (détails omis).
On voit donc que $H^i(X', \mathcal{F}') = 0$ pour $i \geq \max(p + d_p)$,
et l'on gagne.
\end{proof}

\begin{lemma}
\label{lemma-affine-vanishing-higher-direct-images}
Soit $S$ un schéma. Soit $f : X \to Y$ un morphisme affine
d'espaces algébriques sur $S$. Alors
$R^if_*\mathcal{F} = 0$ pour $i > 0$ et tout
$\mathcal{O}_X$-module quasi-cohérent $\mathcal{F}$.
\end{lemma}

\begin{proof}
Rappelons qu'un morphisme affine d'espaces algébriques est représentable.
L'assertion résulte donc de (\ref{equation-representable-higher-direct-image}) et du
lemme \ref{coherent-lemma-relative-affine-vanishing} du chapitre Cohomologie des schémas.
\end{proof}

\begin{lemma}
\label{lemma-relative-affine-cohomology}
Soit $S$ un schéma. Soit $f : X \to Y$ un morphisme affine
d'espaces algébriques sur $S$.
Soit $\mathcal{F}$ un $\mathcal{O}_X$-module quasi-cohérent.
Alors $H^i(X, \mathcal{F}) = H^i(Y, f_*\mathcal{F})$ pour tout $i \geq 0$.
\end{lemma}

\begin{proof}
L'assertion résulte du lemme \ref{lemma-affine-vanishing-higher-direct-images}
et de la suite spectrale de Leray. Voir
Cohomologie sur les sites, lemme \ref{sites-cohomology-lemma-apply-Leray}.
\end{proof}






\section{Cohomologie à support dans un sous-espace fermé}
\label{section-cohomology-support}

\noindent
Cette section est l'analogue des sections du chapitre Cohomologie
\ref{cohomology-section-cohomology-support} et
\ref{cohomology-section-cohomology-support-bis}
et de la section du chapitre Cohomologie \'etale
\ref{etale-cohomology-section-cohomology-support}
pour les faisceaux abéliens sur les espaces algébriques.

\medskip\noindent
Soit $S$ un schéma.
Soit $X$ un espace algébrique sur $S$ et soit $Z \subset X$ un
sous-espace fermé. Soit $\mathcal{F}$ un faisceau abélien sur $X_\etale$. On pose
$$
\Gamma_Z(X, \mathcal{F}) =
\{s \in \mathcal{F}(X) \mid \text{Supp}(s) \subset Z\}
$$
l'ensemble des sections à support dans $Z$
(Propriétés d'espaces, définition \ref{spaces-properties-definition-support}).
C'est un foncteur exact à gauche qui n'est pas exact en général.
On obtient donc un foncteur dérivé
$$
R\Gamma_Z(X, -) : D(X_\etale) \longrightarrow D(\textit{Ab})
$$
et les groupes de cohomologie à support dans $Z$ définis par
$H^q_Z(X, \mathcal{F}) = R^q\Gamma_Z(X, \mathcal{F})$.

\medskip\noindent
Soit $\mathcal{I}$ un faisceau abélien injectif sur $X_\etale$. Soit
$U \subset X$ le sous-espace ouvert complémentaire de $Z$.
Alors l'application de restriction $\mathcal{I}(X) \to \mathcal{I}(U)$ est surjective
(Cohomologie sur les sites, lemme
\ref{sites-cohomology-lemma-restriction-along-monomorphism-surjective})
et a pour noyau $\Gamma_Z(X, \mathcal{I})$. Il s'ensuit immédiatement que
pour $K \in D(X_\etale)$ il existe un triangle distingué
$$
R\Gamma_Z(X, K) \to R\Gamma(X, K) \to R\Gamma(U, K) \to R\Gamma_Z(X, K)[1]
$$
dans $D(\textit{Ab})$. On obtient ainsi une suite exacte longue
de cohomologie
$$
\ldots \to H^i_Z(X, K) \to H^i(X, K) \to H^i(U, K) \to
H^{i + 1}_Z(X, K) \to \ldots
$$
pour tout $K$ dans $D(X_\etale)$.

\medskip\noindent
Pour un faisceau abélien $\mathcal{F}$ sur $X_\etale$, on peut considérer
le {\it sous-faisceau des sections à support dans $Z$}, noté
$\mathcal{H}_Z(\mathcal{F})$, défini par la formule
$$
\mathcal{H}_Z(\mathcal{F})(U) =
\{s \in \mathcal{F}(U) \mid \text{Supp}(s) \subset U \times_X Z\}
$$
Nous utilisons ici le support d'une section défini dans
Propriétés d'espaces, définition \ref{spaces-properties-definition-support}.
D'après l'équivalence de Morphismes d'espaces, lemme
\ref{spaces-morphisms-lemma-closed-immersion-push-pull}
on peut considérer $\mathcal{H}_Z(\mathcal{F})$ comme un faisceau abélien sur
$Z_\etale$. On obtient ainsi un foncteur
$$
\textit{Ab}(X_\etale) \longrightarrow \textit{Ab}(Z_\etale),\quad
\mathcal{F} \longmapsto \mathcal{H}_Z(\mathcal{F})
$$
qui est exact à gauche, mais n'est pas exact en général.

\begin{lemma}
\label{lemma-sections-with-support-acyclic}
Soit $S$ un schéma.
Soit $i : Z \to X$ une immersion fermée d'espaces algébriques sur $S$.
Soit $\mathcal{I}$ un faisceau abélien injectif sur $X_\etale$.
Alors $\mathcal{H}_Z(\mathcal{I})$ est un faisceau abélien injectif
sur $Z_\etale$.
\end{lemma}

\begin{proof}
Remarquons que, pour tout faisceau abélien $\mathcal{G}$ sur $Z_\etale$,
on a
$$
\Hom_Z(\mathcal{G}, \mathcal{H}_Z(\mathcal{F})) =
\Hom_X(i_*\mathcal{G}, \mathcal{F})
$$
car toute section de $i_*\mathcal{G}$ est à support dans $Z$.
Comme $i_*$ est exact (lemme \ref{lemma-finite-higher-direct-image-zero})
et que $\mathcal{I}$ est injectif sur $X_\etale$, on conclut que
$\mathcal{H}_Z(\mathcal{I})$ est injectif sur $Z_\etale$.
\end{proof}

\noindent
Notons
$$
R\mathcal{H}_Z : D(X_\etale) \longrightarrow D(Z_\etale)
$$
le foncteur dérivé. On pose
$\mathcal{H}^q_Z(\mathcal{F}) = R^q\mathcal{H}_Z(\mathcal{F})$, de sorte que
$\mathcal{H}^0_Z(\mathcal{F}) = \mathcal{H}_Z(\mathcal{F})$.
D'après le lemme précédent, on a une suite spectrale de Grothendieck
$$
E_2^{p, q} = H^p(Z, \mathcal{H}^q_Z(\mathcal{F}))
\Rightarrow H^{p + q}_Z(X, \mathcal{F})
$$

\begin{lemma}
\label{lemma-cohomology-with-support-sheaf-on-support}
Soit $S$ un schéma. Soit $i : Z \to X$ une immersion fermée
d'espaces algébriques sur $S$. Soit $\mathcal{G}$ un faisceau abélien injectif
sur $Z_\etale$. Alors $\mathcal{H}^p_Z(i_*\mathcal{G}) = 0$ pour $p > 0$.
\end{lemma}

\begin{proof}
Cela résulte du fait que le foncteur $i_*$ est exact
(lemme \ref{lemma-finite-higher-direct-image-zero}) et transforme
les faisceaux abéliens injectifs en faisceaux abéliens injectifs
(Cohomologie sur les sites, lemme
\ref{sites-cohomology-lemma-pushforward-injective-flat}).
\end{proof}

\begin{lemma}
\label{lemma-etale-localization-sheaf-with-support}
Soit $S$ un schéma. Soit $f : X \to Y$ un morphisme \'etale
d'espaces algébriques sur $S$. Soit $Z \subset Y$ un sous-espace fermé
tel que $f^{-1}(Z) \to Z$ soit un isomorphisme d'espaces algébriques.
Soit $\mathcal{F}$ un faisceau abélien sur $X$. Alors
$$
\mathcal{H}^q_Z(\mathcal{F}) = \mathcal{H}^q_{f^{-1}(Z)}(f^{-1}\mathcal{F})
$$
en tant que faisceaux abéliens sur $Z = f^{-1}(Z)$ et on
a $H^q_Z(Y, \mathcal{F}) = H^q_{f^{-1}(Z)}(X, f^{-1}\mathcal{F})$.
\end{lemma}

\begin{proof}
Comme $f$ est \'etale, l'image inverse d'une résolution injective de $\mathcal{F}$
est une résolution injective de $f^{-1}\mathcal{F}$.
Il suffit donc de vérifier l'égalité pour $\mathcal{H}_Z(-)$,
ce qui résulte des définitions. La démonstration pour la cohomologie à
support est la même. Certains détails sont omis.
\end{proof}

\noindent
Soit $S$ un schéma et soit $X$ un espace algébrique sur $S$.
Soit $T \subset |X|$ un sous-ensemble fermé.
On note $D_T(X_\etale)$ la
sous-catégorie triangulée strictement pleine et saturée de $D(X_\etale)$
formée des objets dont les faisceaux de cohomologie sont à support dans $T$.

\begin{lemma}
\label{lemma-complexes-with-support-on-closed}
Soit $S$ un schéma.
Soit $i : Z \to X$ une immersion fermée d'espaces algébriques sur $S$.
L'application $Ri_* = i_* : D(Z_\etale) \to D(X_\etale)$
induit une équivalence $D(Z_\etale) \to D_{|Z|}(X_\etale)$ de quasi-inverse
$$
i^{-1}|_{D_Z(X_\etale)} = R\mathcal{H}_Z|_{D_{|Z|}(X_\etale)}
$$
\end{lemma}

\begin{proof}
Rappelons que $i^{-1}$ et $i_*$ forment une paire de foncteurs adjoints
exacts telle que $i^{-1}i_*$ est isomorphe au foncteur identité
sur les faisceaux abéliens. Voir
Propriétés d'espaces, lemme
\ref{spaces-properties-lemma-stalk-pullback} et
Morphismes d'espaces, lemme
\ref{spaces-morphisms-lemma-closed-immersion-push-pull}.
Ainsi $i_* : D(Z_\etale) \to D_Z(X_\etale)$ est pleinement fidèle et
$i^{-1}$ fournit
un inverse à gauche. D'autre part, supposons que $K$ soit un objet de
$D_Z(X_\etale)$ et considérons le morphisme d'adjonction
$K \to i_*i^{-1}K$.
Par exactitude de $i_*$ et de $i^{-1}$,
celui-ci induit les morphismes d'adjonction
$H^n(K) \to i_*i^{-1}H^n(K)$ sur les faisceaux de cohomologie.
Comme ces faisceaux de cohomologie
sont à support dans $Z$, ces morphismes d'adjonction sont des isomorphismes
et on conclut que $D(Z_\etale) \to D_Z(X_\etale)$ est une équivalence.

\medskip\noindent
Pour terminer la démonstration, il faut montrer que $R\mathcal{H}_Z(K) = i^{-1}K$
si $K$ est un objet de $D_Z(X_\etale)$. Pour cela, on peut utiliser l'égalité
$K = i_*i^{-1}K$
que nous venons de démontrer. On peut alors
choisir un représentant K-injectif $\mathcal{I}^\bullet$ de
$i^{-1}K$.
Puisque $i_*$ est l'adjoint à droite du foncteur exact
$i^{-1}$, le
complexe $i_*\mathcal{I}^\bullet$ est K-injectif
(Catégories dérivées, lemme \ref{derived-lemma-adjoint-preserve-K-injectives}).
On voit que $R\mathcal{H}_Z(K)$ est calculé par
$\mathcal{H}_Z(i_*\mathcal{I}^\bullet) = \mathcal{I}^\bullet$,
comme voulu.
\end{proof}






\section{Annulation au-delà de la dimension}
\label{section-vanishing-above-dimension}

\noindent
Soit $S$ un schéma. Soit $X$ un espace algébrique quasi-compact et quasi-séparé
sur $S$. Dans ce cas, $|X|$ est un espace spectral, voir
Propriétés d'espaces, lemme
\ref{spaces-properties-lemma-quasi-compact-quasi-separated-spectral}.
De plus, la dimension de $X$ (telle qu'elle est définie dans
Propriétés d'espaces, définition \ref{spaces-properties-definition-dimension})
est égale à la dimension de Krull de $|X|$, voir
Espaces décents, lemme \ref{decent-spaces-lemma-dimension-decent-space}.
Nous allons montrer que, pour les faisceaux quasi-cohérents sur $X$, la
cohomologie s'annule au-delà de la dimension. Ce résultat est déjà intéressant pour les
espaces algébriques quasi-séparés de type fini sur un corps.

\begin{lemma}
\label{lemma-vanishing-above-dimension}
Soit $S$ un schéma. Soit $X$ un espace algébrique quasi-compact et quasi-séparé
sur $S$. Supposons que $\dim(X) \leq d$ pour un entier $d$.
Soit $\mathcal{F}$ un faisceau quasi-cohérent $\mathcal{F}$ sur $X$.
\begin{enumerate}
\item $H^q(X, \mathcal{F}) = 0$ pour $q > d$,
\item $H^d(X, \mathcal{F}) \to H^d(U, \mathcal{F})$ est surjective
pour tout ouvert quasi-compact $U \subset X$,
\item $H^q_Z(X, \mathcal{F}) = 0$ pour $q > d$ pour tout sous-espace fermé
$Z \subset X$ dont le complémentaire est quasi-compact.
\end{enumerate}
\end{lemma}

\begin{proof}
D'après Propriétés d'espaces, lemme
\ref{spaces-properties-lemma-dimension-decent-invariant-under-etale}
tout espace algébrique $Y$ \'etale sur $X$ est de dimension $\leq d$.
Si $Y$ est quasi-séparé, la dimension de $Y$ est égale à la
dimension de Krull de $|Y|$ d'après 
Espaces décents, lemme \ref{decent-spaces-lemma-dimension-decent-space}.
De plus, si $Y$ est un schéma, alors la cohomologie \'etale de $\mathcal{F}$
sur $Y$, resp.\ la cohomologie \'etale de $\mathcal{F}$ à support dans un
sous-schéma fermé coïncide avec la cohomologie usuelle de $\mathcal{F}$,
resp.\ avec la cohomologie usuelle à support dans le sous-schéma fermé.
Voir
Descente, proposition \ref{descent-proposition-same-cohomology-quasi-coherent}
et
Cohomologie \'etale, lemme
\ref{etale-cohomology-lemma-cohomology-with-support-quasi-coherent}.
Nous utiliserons ces faits sans autre mention.

\medskip\noindent
D'après Espaces décents, lemme
\ref{decent-spaces-lemma-filter-quasi-compact-quasi-separated}
il existe un entier $n$ et des sous-espaces ouverts
$$
\emptyset = U_{n + 1} \subset
U_n \subset U_{n - 1} \subset \ldots \subset U_1 = X
$$
ayant la propriété suivante : si l'on pose $T_p = U_p \setminus U_{p + 1}$
(muni de la structure réduite induite), il existe un schéma séparé
quasi-compact $V_p$ et un morphisme \'etale surjectif $f_p : V_p \to U_p$
tel que $f_p^{-1}(T_p) \to T_p$ soit un isomorphisme.

\medskip\noindent
Comme $U_n = V_n$ est un schéma, nos remarques initiales montrent que la cohomologie de
$\mathcal{F}$ sur $U_n$ s'annule en degrés $> d$ d'après
Cohomologie, proposition
\ref{cohomology-proposition-cohomological-dimension-spectral}.
Supposons avoir montré, par récurrence, que
$H^q(U_{p + 1}, \mathcal{F}|_{U_{p + 1}}) = 0$ pour $q > d$.
Il suffit de montrer que $H_{T_p}^q(U_p, \mathcal{F})$ est nul pour
$q > d$ afin d'en déduire l'annulation de la cohomologie
de $\mathcal{F}$ sur $U_p$ en degrés $> d$.
Cependant, on a
$$
H^q_{T_p}(U_p, \mathcal{F}) = H^q_{f_p^{-1}(T_p)}(V_p, \mathcal{F})
$$
d'après le lemme \ref{lemma-etale-localization-sheaf-with-support}
et, puisque $V_p$ est un schéma, l'annulation voulue résulte de
Cohomologie, proposition
\ref{cohomology-proposition-cohomological-dimension-spectral}.
On conclut ainsi que (1) est vraie.

\medskip\noindent
Pour démontrer (2), soit $U \subset X$ un sous-espace ouvert quasi-compact.
Considérons le sous-espace ouvert $U' = U \cup U_n$. Posons $Z = U' \setminus U$.
Alors $g : U_n \to U'$ est un morphisme \'etale tel que
$g^{-1}(Z) \to Z$ soit un isomorphisme. Par conséquent, d'après le
lemme \ref{lemma-etale-localization-sheaf-with-support},
on a $H^q_Z(U', \mathcal{F}) = H^q_Z(U_n, \mathcal{F})$,
qui s'annule en degré $> d$ puisque $U_n$ est un schéma
et que l'on peut appliquer
Cohomologie, proposition
\ref{cohomology-proposition-cohomological-dimension-spectral}.
On conclut que $H^d(U', \mathcal{F}) \to H^d(U, \mathcal{F})$
est surjective. Supposons, par récurrence, avoir ramené
notre problème au cas où $U$ contient $U_{p + 1}$.
On pose alors $U' = U \cup U_p$, on pose $Z = U' \setminus U$, et
on raisonne au moyen du morphisme $f_p : V_p \to U'$ qui est \'etale
et tel que $f_p^{-1}(Z) \to Z$ soit un isomorphisme.
Autrement dit, on voit de nouveau que
$$
H^q_Z(U', \mathcal{F}) = H^q_{f_p^{-1}(Z)}(V_p, \mathcal{F})
$$
et l'on voit de nouveau que ce groupe s'annule en degrés $> d$.
On conclut que $H^d(U', \mathcal{F}) \to H^d(U, \mathcal{F})$
est surjective. On atteint finalement le stade où $U_1 = X \subset U$,
ce qui achève la démonstration.

\medskip\noindent
Un argument formel montre que (2) implique (3).
\end{proof}








\section{Cohomologie et changement de base, I}
\label{section-cohomology-and-base-change}

\noindent
Soit $S$ un schéma.
Soit $f : X \to Y$ un morphisme d'espaces algébriques au-dessus de $S$.
Soit $\mathcal{F}$ un faisceau quasi-cohérent sur $X$.
Supposons en outre que $g : Y' \to Y$ soit un morphisme d'espaces algébriques au-dessus de
$S$. Notons $X' = X_{Y'} = Y' \times_Y X$ le changement de base de $X$ et notons
$f' : X' \to Y'$ le changement de base de $f$.
Notons aussi $g' : X' \to X$ la projection,
et posons $\mathcal{F}' = (g')^*\mathcal{F}$.
La situation est représentée par le diagramme suivant :
\begin{equation}
\label{equation-base-change-diagram}
\vcenter{
\xymatrix{
\mathcal{F}' = (g')^*\mathcal{F} &
X' \ar[r]_{g'} \ar[d]_{f'} &
X \ar[d]^f &
\mathcal{F} \\
Rf'_*\mathcal{F}' &
Y' \ar[r]^g &
Y &
Rf_*\mathcal{F}
}
}
\end{equation}
Voici le cas le plus simple de la propriété de changement de base que nous avons en vue.

\begin{lemma}
\label{lemma-affine-base-change}
Soit $S$ un schéma. Soit $f : X \to Y$ un morphisme affine d'espaces
algébriques au-dessus de $S$. Soit $\mathcal{F}$ un $\mathcal{O}_X$-module quasi-cohérent.
Dans ce cas, $f_*\mathcal{F} \cong Rf_*\mathcal{F}$ est un faisceau quasi-cohérent,
et, pour tout diagramme (\ref{equation-base-change-diagram}),
on a
$$
g^*f_*\mathcal{F} = f'_*(g')^*\mathcal{F}.
$$
\end{lemma}

\begin{proof}
D'après la discussion autour de
(\ref{equation-representable-higher-direct-image}),
on se ramène au cas d'un morphisme affine de schémas, qui
est traité dans Cohomologie des schémas, lemme
\ref{coherent-lemma-affine-base-change}.
\end{proof}

\begin{lemma}[Changement de base plat]
\label{lemma-flat-base-change-cohomology}
Soit $S$ un schéma. Considérons un diagramme cartésien d'espaces algébriques
$$
\xymatrix{
X' \ar[d]_{f'} \ar[r]_{g'} & X \ar[d]^f \\
Y' \ar[r]^g & Y
}
$$
au-dessus de $S$.
Soit $\mathcal{F}$ un $\mathcal{O}_X$-module quasi-cohérent
dont l'image inverse est $\mathcal{F}' = (g')^*\mathcal{F}$.
Supposons que $g$ soit plat et que $f$ soit quasi-compact et quasi-séparé.
Pour tout $i \geq 0$,
\begin{enumerate}
\item la flèche de changement de base de
Cohomologie sur les sites, lemme
\ref{sites-cohomology-lemma-base-change-map-flat-case}
est un isomorphisme
$$
g^*R^if_*\mathcal{F} \longrightarrow R^if'_*\mathcal{F}',
$$
\item si $Y = \Spec(A)$ et $Y' = \Spec(B)$, alors
$H^i(X, \mathcal{F}) \otimes_A B = H^i(X', \mathcal{F}')$.
\end{enumerate}
\end{lemma}

\begin{proof}
Le morphisme $g'$ est plat d'après
Morphismes d'espaces, lemme \ref{spaces-morphisms-lemma-base-change-flat}.
Notons que la platitude de $g$ et de $g'$ équivaut à la platitude
des morphismes de petits sites étales annelés, voir
Morphismes d'espaces, lemme \ref{spaces-morphisms-lemma-flat-morphism-sites}.
On peut donc appliquer
Cohomologie sur les sites, lemme
\ref{sites-cohomology-lemma-base-change-map-flat-case}
pour obtenir une flèche de changement de base
$$
g^*R^pf_*\mathcal{F} \longrightarrow R^pf'_*\mathcal{F}'
$$
Pour montrer que cette flèche est un isomorphisme, on peut travailler localement pour la
topologie étale sur $Y'$. On peut donc supposer que $Y$ et $Y'$ sont des schémas
affines. Posons $Y = \Spec(A)$ et $Y' = \Spec(B)$.
Dans ce cas, il s'agit en fait de montrer que la flèche
$$
H^p(X, \mathcal{F}) \otimes_A B \longrightarrow H^p(X_B, \mathcal{F}_B)
$$
est un isomorphisme, où $X_B = \Spec(B) \times_{\Spec(A)} X$ et
$\mathcal{F}_B$ est l'image inverse de $\mathcal{F}$ sur $X_B$.
Autrement dit, il suffit de démontrer (2).

\medskip\noindent
Fixons un homomorphisme plat d'anneaux $A \to B$ et soit $X$ un espace algébrique
quasi-compact et quasi-séparé sur $A$. Notons que $g' : X_B \to X$
est affine en tant que changement de base de $\Spec(B) \to \Spec(A)$. Par conséquent,
les images directes supérieures $R^i(g')_*\mathcal{F}_B$ sont nulles d'après le
lemme \ref{lemma-affine-vanishing-higher-direct-images}.
Ainsi, $H^p(X_B, \mathcal{F}_B) = H^p(X, g'_*\mathcal{F}_B)$, voir
Cohomologie sur les sites, lemme \ref{sites-cohomology-lemma-apply-Leray}.
De plus, on a
$$
g'_*\mathcal{F}_B = \mathcal{F} \otimes_{\underline{A}} \underline{B}
$$
où $\underline{A}$ et $\underline{B}$ désignent les faisceaux constants d'anneaux
de valeurs respectives $A$ et $B$. En effet, il est clair qu'il existe une flèche
du membre de droite vers celui de gauche. Pour tout schéma affine $U$ étale sur $X$, on a
\begin{align*}
g'_*\mathcal{F}_B(U) & = \mathcal{F}_B(\Spec(B) \times_{\Spec(A)} U) \\
& =
\Gamma(\Spec(B) \times_{\Spec(A)} U,
(\Spec(B) \times_{\Spec(A)} U \to U)^*\mathcal{F}|_U) \\
& =
B \otimes_A \mathcal{F}(U)
\end{align*}
donc cette flèche est un isomorphisme. Écrivons $B = \colim M_i$ comme limite
inductive filtrante de $A$-modules libres de type fini $M_i$ en utilisant le théorème de Lazard, voir
Algèbre, théorème \ref{algebra-theorem-lazard}.
On en déduit que
\begin{align*}
H^p(X, g'_*\mathcal{F}_B) &
= H^p(X, \mathcal{F} \otimes_{\underline{A}} \underline{B}) \\
& = H^p(X, \colim_i \mathcal{F} \otimes_{\underline{A}} \underline{M_i}) \\
& = \colim_i H^p(X, \mathcal{F} \otimes_{\underline{A}} \underline{M_i}) \\
& = \colim_i H^p(X, \mathcal{F}) \otimes_A M_i \\
& = H^p(X, \mathcal{F}) \otimes_A \colim_i M_i \\
& = H^p(X, \mathcal{F}) \otimes_A B
\end{align*}
La première égalité vient de ce que
$g'_*\mathcal{F}_B = \mathcal{F} \otimes_{\underline{A}} \underline{B}$,
comme on l'a vu ci-dessus.
La deuxième vient de ce que $\otimes$ commute aux limites inductives.
La troisième égalité vient de ce que la cohomologie sur $X$ commute aux
limites inductives (voir le
lemme \ref{lemma-colimits}).
La quatrième égalité vient de ce que $M_i$ est libre de type fini (c'est-à-dire de ce que la cohomologie
commute aux sommes directes finies).
La cinquième vient de ce que $\otimes$ commute aux limites inductives.
La sixième résulte du choix de notre système.
\end{proof}

\begin{lemma}
\label{lemma-etale-pull-push-split}
Soit $f : X \to Y$ un morphisme quasi-compact, séparé et \'etale
d'espaces algébriques. Alors, pour tout $\mathcal{O}_X$-module quasi-cohérent
$\mathcal{F}$, le morphisme $f^*f_*\mathcal{F} \to \mathcal{F}$ est scindé.
\end{lemma}

\begin{proof}
Considérons le diagramme cartésien
$$
\xymatrix{
X \times_Y X \ar[r]_-p \ar[d]_q & X \ar[d]^f \\
X \ar[r]^f & Y
}
$$
D'après le Lemme \ref{lemma-affine-base-change}, on a
$f^*f_*\mathcal{F} = q_*p^*\mathcal{F}$. Le morphisme
$\Delta : X \to X \times_Y X$ est ouvert (car c'est un morphisme entre
espaces algébriques \'etales sur $Y$) et fermé puisque $f$ est séparé.
On voit donc que $p^*\mathcal{F}$ est somme directe de
$\Delta_*\mathcal{F}$ et d'un module quasi-cohérent à support
dans le complémentaire (ouvert et fermé) de $\Delta(X)$.
En suivant les morphismes, le lecteur vérifie que
$$
\mathcal{F} = q_*\Delta_*\mathcal{F} \to
q_*p^*\mathcal{F} = f^*f_*\mathcal{F} \to
\mathcal{F}
$$
est le morphisme identité. Nous omettons les détails.
\end{proof}



\section{Modules cohérents sur les espaces algébriques localement noethériens}
\label{section-coherent}

\noindent
Cette section est l'analogue de
Cohomologie des schémas, section \ref{coherent-section-coherent-sheaves}.
Dans Modules sur les sites, Définition \ref{sites-modules-definition-site-local},
nous avons défini les modules cohérents sur tout topos annelé. Nous employons cette notion
pour définir les modules cohérents sur les espaces algébriques localement noethériens.
Bien qu'il soit possible de travailler plus généralement avec des modules cohérents,
nous nous abstiendrons de le faire.

\begin{definition}
\label{definition-coherent}
Soit $S$ un schéma. Soit $X$ un espace algébrique localement noethérien sur $S$.
Un module quasi-cohérent $\mathcal{F}$ sur $X$ est dit {\it cohérent}
si $\mathcal{F}$ est un $\mathcal{O}_X$-module cohérent sur le site
$X_\etale$ au sens de
Modules sur les sites, Définition \ref{sites-modules-definition-site-local}.
\end{definition}

\noindent
Cette définition est compatible avec la notion déjà existante de
module cohérent sur un schéma localement noethérien; voir l'assertion (5) de
Propriétés des espaces, section
\ref{spaces-properties-section-properties-modules} (ou, plus directement,
Descente, Lemme \ref{descent-lemma-equivalence-quasi-coherent-properties}).
Désormais, si $X$ est un schéma localement noethérien sur $S$,
nous ne distinguerons donc pas un module cohérent sur $X$ considéré
comme schéma d'un module cohérent sur $X$ considéré comme espace algébrique;
ceci est compatible avec les identifications correspondantes des catégories
de modules quasi-cohérents décrites dans
Propriétés des espaces, section \ref{spaces-properties-section-quasi-coherent}.

\medskip\noindent
Cela étant, le lemme suivant donne une caractérisation concrète
des modules cohérents sur les espaces algébriques localement
noethériens.

\begin{lemma}
\label{lemma-coherent-Noetherian}
Soit $S$ un schéma.
Soit $X$ un espace algébrique localement noethérien sur $S$.
Soit $\mathcal{F}$ un $\mathcal{O}_X$-module.
Les conditions suivantes sont équivalentes :
\begin{enumerate}
\item $\mathcal{F}$ est cohérent,
\item $\mathcal{F}$ est un $\mathcal{O}_X$-module quasi-cohérent de type fini,
\item $\mathcal{F}$ est un $\mathcal{O}_X$-module de présentation finie,
\item pour tout morphisme \'etale $\varphi : U \to X$, où $U$ est un schéma,
l'image inverse $\varphi^*\mathcal{F}$ est un module cohérent sur $U$, et
\item il existe un morphisme \'etale surjectif $\varphi : U \to X$,
où $U$ est un schéma, tel que l'image inverse $\varphi^*\mathcal{F}$ soit
un module cohérent sur $U$.
\end{enumerate}
En particulier, $\mathcal{O}_X$ est cohérent, tout
$\mathcal{O}_X$-module inversible est cohérent et, plus généralement, tout
$\mathcal{O}_X$-module localement libre de type fini est cohérent.
\end{lemma}

\begin{proof}
Notons que, si $X$ est un espace algébrique localement noethérien et si
$U \to X$ est un morphisme \'etale, alors $U$ est localement noethérien; voir
Propriétés des espaces, section \ref{spaces-properties-section-types-properties}.
Le lemme résulte alors des points (1) -- (5) établis dans
Propriétés des espaces, section \ref{spaces-properties-section-properties-modules},
et du résultat correspondant pour les modules cohérents sur les schémas
localement noethériens; voir
Cohomologie des schémas, Lemme \ref{coherent-lemma-coherent-Noetherian}.
\end{proof}

\begin{lemma}
\label{lemma-coherent-abelian-Noetherian}
Soit $S$ un schéma. Soit $X$ un espace algébrique localement noethérien sur $S$.
La catégorie des $\mathcal{O}_X$-modules cohérents est abélienne. Plus précisément,
le noyau et le conoyau d'un morphisme de $\mathcal{O}_X$-modules cohérents sont
cohérents. Toute extension de faisceaux cohérents est cohérente.
\end{lemma}

\begin{proof}
Choisissons un schéma $U$ et un morphisme \'etale surjectif $f : U \to X$.
L'image inverse $f^*$ est un foncteur exact, car elle est un foncteur de restriction; voir
Propriétés des espaces, Équation
(\ref{spaces-properties-equation-restrict-modules}).
D'après le
Lemme \ref{lemma-coherent-Noetherian}, pour vérifier qu'un
$\mathcal{O}_X$-module $\mathcal{F}$ est
cohérent, il suffit de vérifier que $f^*\mathcal{F}$ est cohérent. Le
lemme résulte donc du cas des schémas, qui est
Cohomologie des schémas, Lemme \ref{coherent-lemma-coherent-abelian-Noetherian}.
\end{proof}

\noindent
Les modules cohérents forment une sous-catégorie de Serre de la
catégorie des $\mathcal{O}_X$-modules quasi-cohérents. Cela n'est pas vrai
pour les modules sur un topos annelé quelconque.

\begin{lemma}
\label{lemma-coherent-Noetherian-quasi-coherent-sub-quotient}
Soit $S$ un schéma.
Soit $X$ un espace algébrique localement noethérien sur $S$.
Soit $\mathcal{F}$ un $\mathcal{O}_X$-module cohérent.
Tout sous-module quasi-cohérent de $\mathcal{F}$ est cohérent.
Tout module quotient quasi-cohérent de $\mathcal{F}$ est cohérent.
\end{lemma}

\begin{proof}
Choisissons un schéma $U$ et un morphisme \'etale surjectif $f : U \to X$.
L'image inverse $f^*$ est un foncteur exact, car elle est un foncteur de restriction; voir
Propriétés des espaces, Équation
(\ref{spaces-properties-equation-restrict-modules}).
D'après le
Lemme \ref{lemma-coherent-Noetherian}, pour vérifier qu'un
$\mathcal{O}_X$-module $\mathcal{G}$ est
cohérent, il suffit de vérifier que $f^*\mathcal{H}$ est cohérent. Le
lemme résulte donc du cas des schémas, qui est
Cohomologie des schémas, Lemme
\ref{coherent-lemma-coherent-Noetherian-quasi-coherent-sub-quotient}.
\end{proof}

\begin{lemma}
\label{lemma-tensor-hom-coherent}
Soit $S$ un schéma.
Soit $X$ un espace algébrique localement noethérien sur $S$.
Soient $\mathcal{F}$, $\mathcal{G}$ des $\mathcal{O}_X$-modules cohérents.
Les $\mathcal{O}_X$-modules $\mathcal{F} \otimes_{\mathcal{O}_X} \mathcal{G}$
et $\SheafHom_{\mathcal{O}_X}(\mathcal{F}, \mathcal{G})$ sont
cohérents.
\end{lemma}

\begin{proof}
Par le Lemme \ref{lemma-coherent-Noetherian}, cela résulte du résultat
pour les schémas; voir
Cohomologie des schémas, Lemme \ref{coherent-lemma-tensor-hom-coherent}.
\end{proof}

\begin{lemma}
\label{lemma-local-isomorphism}
Soit $S$ un schéma. Soit $X$ un espace algébrique localement noethérien sur $S$.
Soient $\mathcal{F}$, $\mathcal{G}$ des $\mathcal{O}_X$-modules cohérents.
Soit $\varphi : \mathcal{G} \to \mathcal{F}$ un homomorphisme
de $\mathcal{O}_X$-modules. Soit $\overline{x}$ un point géométrique de $X$
au-dessus de $x \in |X|$.
\begin{enumerate}
\item Si $\mathcal{F}_{\overline{x}} = 0$, il existe un voisinage
ouvert $X' \subset X$ de $x$ tel que $\mathcal{F}|_{X'} = 0$.
\item Si $\varphi_{\overline{x}} : \mathcal{G}_{\overline{x}} \to
\mathcal{F}_{\overline{x}}$ est injectif, il existe un voisinage
ouvert $X' \subset X$ de $x$ tel que $\varphi|_{X'}$ soit injectif.
\item Si $\varphi_{\overline{x}} : \mathcal{G}_{\overline{x}} \to
\mathcal{F}_{\overline{x}}$ est surjectif, il existe un voisinage
ouvert $X' \subset X$ de $x$ tel que $\varphi|_{X'}$ soit surjectif.
\item Si $\varphi_{\overline{x}} : \mathcal{G}_{\overline{x}} \to
\mathcal{F}_{\overline{x}}$ est bijectif, il existe un voisinage
ouvert $X' \subset X$ de $x$ tel que $\varphi|_{X'}$ soit un isomorphisme.
\end{enumerate}
\end{lemma}

\begin{proof}
Soit $\varphi : U \to X$ un morphisme \'etale, où $U$ est un schéma, et
soit $u \in U$ un point d'image $x$. D'après
Propriétés des espaces, Lemmes
\ref{spaces-properties-lemma-stalk-quasi-coherent} et
\ref{spaces-properties-lemma-describe-etale-local-ring}
ainsi que
Compléments d'algèbre, Lemme \ref{more-algebra-lemma-dumb-properties-henselization},
on voit que $\varphi_{\overline{x}}$ est injectif, surjectif ou bijectif
si et seulement si $\varphi_u : \varphi^*\mathcal{F}_u \to \varphi^*\mathcal{G}_u$
possède la propriété correspondante. On peut donc appliquer la version de
ce lemme pour les schémas pour voir que (quitte à restreindre $U$) le morphisme
$\varphi^*\mathcal{F} \to \varphi^*\mathcal{G}$ est injectif, surjectif
ou un isomorphisme. Soit $X' \subset X$ le sous-espace ouvert correspondant
à $|\varphi|(|U|) \subset |X|$; voir
Propriétés des espaces, Lemme \ref{spaces-properties-lemma-open-subspaces}.
Puisque $\{U \to X'\}$ est un recouvrement pour la topologie \'etale, on conclut
que $\varphi|_{X'}$ est injectif, surjectif ou un isomorphisme, ce qui donne le résultat voulu.
Enfin, remarquons que (1) résulte de (2) en considérant le morphisme
$\mathcal{F} \to 0$.
\end{proof}

\begin{lemma}
\label{lemma-coherent-support-closed}
Soit $S$ un schéma. Soit $X$ un espace algébrique localement noethérien sur $S$.
Soit $\mathcal{F}$ un $\mathcal{O}_X$-module cohérent. Soit $i : Z \to X$
le support schématique de $\mathcal{F}$, et soit $\mathcal{G}$
le $\mathcal{O}_Z$-module quasi-cohérent tel que
$i_*\mathcal{G} = \mathcal{F}$; voir
Morphismes d'espaces, Définition
\ref{spaces-morphisms-definition-scheme-theoretic-support}.
Alors $\mathcal{G}$ est un $\mathcal{O}_Z$-module cohérent.
\end{lemma}

\begin{proof}
L'énoncé du lemme a un sens, car un module cohérent est en
particulier de type fini. De plus, comme $Z \to X$ est une immersion fermée,
ce morphisme est localement de type fini, et $Z$ est donc localement noethérien; voir
Morphismes d'espaces, Lemmes
\ref{spaces-morphisms-lemma-immersion-locally-finite-type} et
\ref{spaces-morphisms-lemma-locally-finite-type-locally-noetherian}.
Enfin, puisque $\mathcal{G}$ est de type fini, c'est un
$\mathcal{O}_Z$-module cohérent d'après le
Lemme \ref{lemma-coherent-Noetherian}.
\end{proof}

\begin{lemma}
\label{lemma-i-star-equivalence}
Soit $S$ un schéma. Soit $i : Z \to X$ une immersion fermée d'espaces
algébriques localement noethériens sur $S$.
Soit $\mathcal{I} \subset \mathcal{O}_X$ le faisceau quasi-cohérent d'idéaux
définissant $Z$. Le foncteur $i_*$ induit une équivalence entre la
catégorie des $\mathcal{O}_X$-modules cohérents annulés par $\mathcal{I}$
et la catégorie des $\mathcal{O}_Z$-modules cohérents.
\end{lemma}

\begin{proof}
Le foncteur est pleinement fidèle d'après
Morphismes d'espaces, Lemme \ref{spaces-morphisms-lemma-i-star-equivalence}.
Soit $\mathcal{F}$ un $\mathcal{O}_X$-module cohérent
annulé par $\mathcal{I}$. D'après
Morphismes d'espaces, Lemme \ref{spaces-morphisms-lemma-i-star-equivalence},
on peut écrire $\mathcal{F} = i_*\mathcal{G}$ pour un faisceau quasi-cohérent
$\mathcal{G}$ sur $Z$. Pour vérifier que $\mathcal{G}$ est cohérent,
on peut travailler localement pour la topologie \'etale (Lemme \ref{lemma-coherent-Noetherian}).
En choisissant un recouvrement \'etale par un schéma, on conclut que
$\mathcal{G}$ est cohérent par le cas des schémas
(Cohomologie des schémas, Lemme \ref{coherent-lemma-i-star-equivalence}).
Ainsi le foncteur est pleinement fidèle, ce qui termine la démonstration.
\end{proof}

\begin{lemma}
\label{lemma-finite-pushforward-coherent}
Soit $S$ un schéma. Soit $f : X \to Y$ un morphisme fini d'espaces
algébriques sur $S$, avec $Y$ localement noethérien. Soit $\mathcal{F}$ un
$\mathcal{O}_X$-module cohérent. Supposons $f$ fini et $Y$ localement
noethérien. Alors $R^pf_*\mathcal{F} = 0$ pour $p > 0$, et
$f_*\mathcal{F}$ est cohérent.
\end{lemma}

\begin{proof}
Choisissons un schéma $V$ et un morphisme \'etale surjectif $V \to Y$.
Alors $V \times_Y X \to V$ est un morphisme fini de schémas localement
noethériens. D'après (\ref{equation-representable-higher-direct-image}), on se ramène
au cas des schémas, qui est
Cohomologie des schémas, Lemme \ref{coherent-lemma-finite-pushforward-coherent}.
\end{proof}






\section{Faisceaux cohérents sur les espaces noethériens}
\label{section-coherent-quasi-compact}

\noindent
Dans cette section, nous donnons quelques propriétés des faisceaux cohérents sur
les espaces algébriques noethériens.

\begin{lemma}
\label{lemma-acc-coherent}
Soit $S$ un schéma. Soit $X$ un espace algébrique noethérien sur $S$.
Soit $\mathcal{F}$ un $\mathcal{O}_X$-module cohérent.
Les sous-modules quasi-cohérents de $\mathcal{F}$ vérifient la condition
de chaîne ascendante. Autrement dit, pour toute suite
$$
\mathcal{F}_1 \subset \mathcal{F}_2 \subset \ldots \subset \mathcal{F}
$$
de sous-modules quasi-cohérents, on a
$\mathcal{F}_n = \mathcal{F}_{n + 1} = \ldots $ pour un certain $n \geq 0$.
\end{lemma}

\begin{proof}
Choisissons un schéma affine $U$ et un morphisme \'etale surjectif $U \to X$
(voir Propriétés des espaces, Lemme
\ref{spaces-properties-lemma-quasi-compact-affine-cover}).
Alors $U$ est un schéma noethérien (d'après
Morphismes d'espaces, Lemme
\ref{spaces-morphisms-lemma-locally-finite-type-locally-noetherian}).
Si $\mathcal{F}_n|_U = \mathcal{F}_{n + 1}|_U = \ldots$,
alors $\mathcal{F}_n = \mathcal{F}_{n + 1} = \ldots$.
Le résultat découle donc du cas des schémas; voir
Cohomologie des schémas, Lemme \ref{coherent-lemma-acc-coherent}.
\end{proof}

\begin{lemma}
\label{lemma-power-ideal-kills-sheaf}
Soit $S$ un schéma. Soit $X$ un espace algébrique noethérien sur $S$.
Soit $\mathcal{F}$ un faisceau cohérent sur $X$. Soit
$\mathcal{I} \subset \mathcal{O}_X$ un faisceau quasi-cohérent d'idéaux
correspondant à un sous-espace fermé $Z \subset X$. Il existe alors un
$n \geq 0$ tel que $\mathcal{I}^n\mathcal{F} = 0$ si et seulement si
$\text{Supp}(\mathcal{F}) \subset Z$ ensemblistement.
\end{lemma}

\begin{proof}
Choisissons un schéma affine $U$ et un morphisme \'etale surjectif $U \to X$
(voir Propriétés des espaces, Lemme
\ref{spaces-properties-lemma-quasi-compact-affine-cover}).
Alors $U$ est un schéma noethérien (d'après
Morphismes d'espaces, Lemme
\ref{spaces-morphisms-lemma-locally-finite-type-locally-noetherian}).
Notons que $\mathcal{I}^n\mathcal{F}|_U = 0$ si et seulement si
$\mathcal{I}^n\mathcal{F} = 0$, et il en va de même pour la condition sur
le support. Le résultat découle donc du cas des schémas; voir
Cohomologie des schémas, Lemme \ref{coherent-lemma-power-ideal-kills-sheaf}.
\end{proof}

\begin{lemma}[Artin-Rees]
\label{lemma-Artin-Rees}
Soit $S$ un schéma. Soit $X$ un espace algébrique noethérien sur $S$.
Soit $\mathcal{F}$ un faisceau cohérent sur $X$. Soit
$\mathcal{G} \subset \mathcal{F}$ un sous-faisceau quasi-cohérent.
Soit $\mathcal{I} \subset \mathcal{O}_X$ un faisceau quasi-cohérent
d'idéaux. Il existe alors un $c \geq 0$ tel que, pour tout $n \geq c$, on
ait
$$
\mathcal{I}^{n - c}(\mathcal{I}^c\mathcal{F} \cap \mathcal{G})
=
\mathcal{I}^n\mathcal{F} \cap \mathcal{G}
$$
\end{lemma}

\begin{proof}
Choisissons un schéma affine $U$ et un morphisme \'etale surjectif $U \to X$
(voir Propriétés des espaces, Lemme
\ref{spaces-properties-lemma-quasi-compact-affine-cover}).
Alors $U$ est un schéma noethérien (d'après
Morphismes d'espaces, Lemme
\ref{spaces-morphisms-lemma-locally-finite-type-locally-noetherian}).
L'égalité du lemme est vérifiée si et seulement si elle l'est après
restriction à $U$. Le résultat découle donc du cas des schémas; voir
Cohomologie des schémas, Lemme \ref{coherent-lemma-Artin-Rees}.
\end{proof}

\begin{lemma}
\label{lemma-homs-over-open}
Soit $S$ un schéma. Soit $X$ un espace algébrique noethérien sur $S$.
Soit $\mathcal{F}$ un $\mathcal{O}_X$-module quasi-cohérent.
Soit $\mathcal{G}$ un $\mathcal{O}_X$-module cohérent.
Soit $\mathcal{I} \subset \mathcal{O}_X$ un faisceau quasi-cohérent
d'idéaux. Notons $Z \subset X$ le sous-espace fermé correspondant et
posons $U = X \setminus Z$. Il existe un isomorphisme canonique
$$
\colim_n \Hom_{\mathcal{O}_X}(\mathcal{I}^n\mathcal{G}, \mathcal{F})
\longrightarrow
\Hom_{\mathcal{O}_U}(\mathcal{G}|_U, \mathcal{F}|_U).
$$
En particulier, on a un isomorphisme
$$
\colim_n \Hom_{\mathcal{O}_X}(\mathcal{I}^n, \mathcal{F})
\longrightarrow
\Gamma(U, \mathcal{F}).
$$
\end{lemma}

\begin{proof}
Soit $W$ un schéma affine et soit $W \to X$ un morphisme \'etale
surjectif (voir Propriétés des espaces, Lemme
\ref{spaces-properties-lemma-quasi-compact-affine-cover}).
Posons $R = W \times_X W$. Alors $W$ et $R$ sont des schémas noethériens; voir
Morphismes d'espaces, Lemme
\ref{spaces-morphisms-lemma-locally-finite-type-locally-noetherian}.
Le résultat vaut donc pour les restrictions de $\mathcal{F}$, $\mathcal{G}$,
ainsi que de $\mathcal{I}$, $U$, $Z$, à $W$ et à $R$, d'après
Cohomologie des schémas, Lemme \ref{coherent-lemma-homs-over-open}.
Il en résulte formellement que le résultat vaut sur $X$.
\end{proof}






\section{Dévissage des faisceaux cohérents}
\label{section-devissage}

\noindent
Cette section est l'analogue de
Cohomologie des schémas, section \ref{coherent-section-devissage}.

\begin{lemma}
\label{lemma-prepare-filter-support}
Soit $S$ un schéma. Soit $X$ un espace algébrique noethérien sur $S$.
Soit $\mathcal{F}$ un faisceau cohérent sur $X$. Supposons que
$\text{Supp}(\mathcal{F}) = Z \cup Z'$, où $Z$, $Z'$ sont fermés.
Il existe alors une suite exacte courte de faisceaux cohérents
$$
0 \to \mathcal{G}' \to \mathcal{F} \to \mathcal{G} \to 0
$$
avec $\text{Supp}(\mathcal{G}') \subset Z'$ et
$\text{Supp}(\mathcal{G}) \subset Z$.
\end{lemma}

\begin{proof}
Soit $\mathcal{I} \subset \mathcal{O}_X$ le faisceau d'idéaux
définissant sur $Z$ la structure réduite induite de sous-espace fermé; voir
Propriétés des espaces, Lemme
\ref{spaces-properties-lemma-reduced-closed-subspace}.
Considérons les sous-faisceaux
$\mathcal{G}'_n = \mathcal{I}^n\mathcal{F}$ et les
quotients $\mathcal{G}_n = \mathcal{F}/\mathcal{I}^n\mathcal{F}$.
Pour tout $n$, on a une suite exacte courte
$$
0 \to \mathcal{G}'_n \to \mathcal{F} \to \mathcal{G}_n \to 0
$$
Pour tout point géométrique $\overline{x}$ de $Z' \setminus Z$, on a
$\mathcal{I}_{\overline{x}} = \mathcal{O}_{X, \overline{x}}$
et donc $\mathcal{G}_{n, \overline{x}} = 0$. Ainsi, on voit que
$\text{Supp}(\mathcal{G}_n) \subset Z$. Notons que $X \setminus Z'$
est un espace algébrique noethérien. D'après le
Lemme \ref{lemma-power-ideal-kills-sheaf},
il existe donc un $n$ tel que $\mathcal{G}'_n|_{X \setminus Z'} =
\mathcal{I}^n\mathcal{F}|_{X \setminus Z'} = 0$.
Pour un tel $n$, on voit que $\text{Supp}(\mathcal{G}'_n) \subset Z'$.
Il suffit donc de poser $\mathcal{G}' = \mathcal{G}'_n$ et $\mathcal{G} = \mathcal{G}_n$
pour conclure.
\end{proof}

\noindent
Dans la suite, nous emploierons librement le support schématique des
modules de type fini défini dans Morphismes d'espaces, Définition
\ref{spaces-morphisms-definition-scheme-theoretic-support}.

\begin{lemma}
\label{lemma-prepare-filter-irreducible}
Soit $S$ un schéma. Soit $X$ un espace algébrique noethérien sur $S$.
Soit $\mathcal{F}$ un faisceau cohérent sur $X$. Supposons que le support
schématique de $\mathcal{F}$ soit un sous-espace réduit $Z \subset X$ tel que
$|Z|$ soit irréductible. Il existe alors un entier $r > 0$, un
faisceau non nul d'idéaux $\mathcal{I} \subset \mathcal{O}_Z$ et un
morphisme injectif de faisceaux cohérents
$$
i_*\left(\mathcal{I}^{\oplus r}\right) \to \mathcal{F}
$$
dont le conoyau est à support dans un sous-espace fermé strict de $Z$.
\end{lemma}

\begin{proof}
Par hypothèse, il existe un $\mathcal{O}_Z$-module cohérent
$\mathcal{G}$ de support $Z$ tel que $\mathcal{F} \cong i_*\mathcal{G}$; voir
Lemme \ref{lemma-coherent-support-closed}. Il suffit donc de démontrer le
lemme dans le cas $Z = X$ et $i = \text{id}$.

\medskip\noindent
Par Propriétés des espaces, Proposition
\ref{spaces-properties-proposition-locally-quasi-separated-open-dense-scheme}
il existe un sous-espace ouvert dense $U \subset X$ qui est un schéma. Notons que
$U$ est un schéma intègre noethérien. Quitte à restreindre $U$, on peut supposer
que $\mathcal{F}|_U \cong \mathcal{O}_U^{\oplus r}$ (par exemple d'après
Cohomologie des schémas, Lemme \ref{coherent-lemma-prepare-filter-irreducible},
ou par un argument direct d'algèbre). Soit $\mathcal{I} \subset \mathcal{O}_X$
un faisceau quasi-cohérent d'idéaux dont le sous-espace fermé associé
est le complémentaire de $U$ dans $X$ (voir, par exemple,
Propriétés des espaces, section \ref{spaces-properties-section-reduced}). 
D'après le Lemme \ref{lemma-homs-over-open}, il existe un $n \geq 0$ et un
morphisme $\mathcal{I}^n(\mathcal{O}_X^{\oplus r}) \to \mathcal{F}$
qui redonne notre isomorphisme sur $U$. Puisque
$\mathcal{I}^n(\mathcal{O}_X^{\oplus r}) = (\mathcal{I}^n)^{\oplus r}$,
on obtient un morphisme comme dans le lemme. Il est injectif : en effet, si $\sigma$ est
une section non nulle de $\mathcal{I}^{\oplus r}$ sur un schéma $W$ \'etale
sur $X$, alors, puisque $X$, et donc $W$, est réduit, le support de $\sigma$
contient un ouvert non vide de $W$. Or le noyau de
$(\mathcal{I}^n)^{\oplus r} \to \mathcal{F}$ est nul
sur un ouvert dense; $\sigma$ ne peut donc pas être une section du noyau.
\end{proof}

\begin{lemma}
\label{lemma-coherent-filter}
Soit $S$ un schéma. Soit $X$ un espace algébrique noethérien sur $S$.
Soit $\mathcal{F}$ un faisceau cohérent sur $X$. Il existe une filtration
$$
0 = \mathcal{F}_0 \subset \mathcal{F}_1 \subset
\ldots \subset \mathcal{F}_m = \mathcal{F}
$$
par des sous-faisceaux cohérents telle que, pour chaque $j = 1, \ldots, m$,
il existe un sous-espace fermé réduit $Z_j \subset X$ tel que $|Z_j|$
soit irréductible, et un faisceau d'idéaux $\mathcal{I}_j \subset \mathcal{O}_{Z_j}$
tels que
$$
\mathcal{F}_j/\mathcal{F}_{j - 1}
\cong (Z_j \to X)_* \mathcal{I}_j
$$
\end{lemma}

\begin{proof}
Considérons l'ensemble
$$
\mathcal{T} =
\left\{
\begin{matrix}
T \subset |X|
\text{ fermé tel qu'il existe un faisceau cohérent }
\mathcal{F} \\
\text{ avec }
\text{Supp}(\mathcal{F}) = T
\text{ pour lequel le lemme est faux}
\end{matrix}
\right\}
$$
Nous voulons montrer que $\mathcal{T}$ est vide. Sinon,
puisque $|X|$ est noethérien (Propriétés des espaces, Lemme
\ref{spaces-properties-lemma-Noetherian-topology})
on peut choisir un élément minimal $T \in \mathcal{T}$. Cela signifie qu'il
existe un faisceau cohérent $\mathcal{F}$ sur $X$, de support $T$,
pour lequel le lemme n'est pas vrai. Il est clair que $T \not = \emptyset$, car
le seul faisceau de support vide est le faisceau nul, pour lequel le
lemme est vrai (avec $m = 0$).

\medskip\noindent
Si $T$ n'est pas irréductible, on peut écrire $T = Z_1 \cup Z_2$,
où $Z_1, Z_2$ sont fermés et strictement contenus dans $T$.
On peut alors appliquer le Lemme \ref{lemma-prepare-filter-support}
pour obtenir une suite exacte courte de faisceaux cohérents
$$
0 \to
\mathcal{G}_1 \to
\mathcal{F} \to
\mathcal{G}_2 \to 0
$$
avec $\text{Supp}(\mathcal{G}_i) \subset Z_i$. Par minimalité de
$T$, chacun des $\mathcal{G}_i$ admet une filtration comme dans l'énoncé
du lemme. En considérant la filtration induite sur $\mathcal{F}$,
on aboutit à une contradiction. On en conclut
que $T$ est irréductible.

\medskip\noindent
Supposons $T$ irréductible. Soit $\mathcal{J}$ le faisceau d'idéaux
définissant sur $T$ la structure réduite induite de sous-espace fermé;
voir Propriétés des espaces, Lemme
\ref{spaces-properties-lemma-reduced-closed-subspace}.
D'après le Lemme \ref{lemma-power-ideal-kills-sheaf}, il existe
un $n \geq 0$ tel que $\mathcal{J}^n\mathcal{F} = 0$. On obtient donc
une filtration
$$
0 = \mathcal{I}^n\mathcal{F} \subset \mathcal{I}^{n - 1}\mathcal{F}
\subset \ldots \subset \mathcal{I}\mathcal{F} \subset \mathcal{F}
$$
dont chacun des sous-quotients successifs est annulé par $\mathcal{J}$.
Si chacun de ces sous-quotients admet une filtration comme dans l'énoncé
du lemme, il en va de même pour $\mathcal{F}$. Autrement dit, on peut
supposer que $\mathcal{J}$ annule $\mathcal{F}$.

\medskip\noindent
Supposons $T$ irréductible et $\mathcal{J}\mathcal{F} = 0$, où
$\mathcal{J}$ est comme ci-dessus. Alors le support schématique de
$\mathcal{F}$ est $T$; voir
Morphismes d'espaces, Lemme \ref{spaces-morphisms-lemma-i-star-equivalence}.
On peut donc appliquer le Lemme \ref{lemma-prepare-filter-irreducible}.
On obtient une suite exacte courte
$$
0 \to
i_*(\mathcal{I}^{\oplus r}) \to
\mathcal{F} \to
\mathcal{Q} \to 0
$$
où le support de $\mathcal{Q}$ est une partie fermée stricte de $T$.
On voit donc que $\mathcal{Q}$ admet une filtration du type voulu
par minimalité de $T$. Mais il en va alors manifestement de même pour $\mathcal{F}$,
d'où la contradiction finale.
\end{proof}

\begin{lemma}
\label{lemma-property-initial}
Soit $S$ un schéma. Soit $X$ un espace algébrique noethérien sur $S$.
Soit $\mathcal{P}$ une propriété des faisceaux cohérents sur $X$. Supposons que
\begin{enumerate}
\item pour toute suite exacte courte de faisceaux cohérents
$$
0 \to \mathcal{F}_1 \to \mathcal{F} \to \mathcal{F}_2 \to 0
$$
si les $\mathcal{F}_i$, $i = 1, 2$, vérifient la propriété $\mathcal{P}$,
il en va de même pour $\mathcal{F}$;
\item pour tout sous-espace fermé réduit $Z \subset X$ tel que $|Z|$ soit irréductible
et tout faisceau quasi-cohérent d'idéaux $\mathcal{I} \subset \mathcal{O}_Z$,
la propriété $\mathcal{P}$ est vérifiée par $i_*\mathcal{I}$.
\end{enumerate}
Alors la propriété $\mathcal{P}$ est vérifiée par tout faisceau cohérent sur $X$.
\end{lemma}

\begin{proof}
Notons d'abord que, si $\mathcal{F}$ est un faisceau cohérent muni d'une filtration
$$
0 = \mathcal{F}_0 \subset \mathcal{F}_1 \subset
\ldots \subset \mathcal{F}_m = \mathcal{F}
$$
par des sous-faisceaux cohérents telle que chaque $\mathcal{F}_i/\mathcal{F}_{i - 1}$
vérifie la propriété $\mathcal{P}$, il en va de même pour $\mathcal{F}$.
Cela résulte de la propriété (1) de $\mathcal{P}$.
D'autre part, d'après le Lemme \ref{lemma-coherent-filter},
on peut munir tout $\mathcal{F}$ d'une filtration
dont les sous-quotients successifs sont comme dans (2).
Le lemme en résulte.
\end{proof}

\noindent
Voici une variante plus utile du lemme précédent.

\begin{lemma}
\label{lemma-property-higher-rank-cohomological}
Soit $S$ un schéma. Soit $X$ un espace algébrique noethérien sur $S$.
Soit $\mathcal{P}$ une propriété des faisceaux cohérents sur $X$. Supposons que
\begin{enumerate}
\item pour toute suite exacte courte de faisceaux cohérents
$$
0 \to \mathcal{F}_1 \to \mathcal{F} \to \mathcal{F}_2 \to 0
$$
si les $\mathcal{F}_i$, $i = 1, 2$, vérifient la propriété $\mathcal{P}$,
il en va de même pour $\mathcal{F}$;
\item si $\mathcal{P}$ est vérifiée par $\mathcal{F}^{\oplus r}$ pour
un certain $r \geq 1$, elle est vérifiée par $\mathcal{F}$;
\item pour tout sous-espace fermé réduit $i : Z \to X$ tel que
$|Z|$ soit irréductible, il existe un faisceau cohérent $\mathcal{G}$ sur $Z$
tel que
\begin{enumerate}
\item $\text{Supp}(\mathcal{G}) = Z$,
\item pour tout faisceau quasi-cohérent non nul d'idéaux
$\mathcal{I} \subset \mathcal{O}_Z$, il existe un sous-faisceau quasi-cohérent
$\mathcal{G}' \subset \mathcal{I}\mathcal{G}$ tel que
$\text{Supp}(\mathcal{G}/\mathcal{G}')$ soit une partie fermée stricte de $|Z|$
et que la propriété $\mathcal{P}$ soit vérifiée par $i_*\mathcal{G}'$.
\end{enumerate}
\end{enumerate}
Alors la propriété $\mathcal{P}$ est vérifiée par tout faisceau cohérent sur $X$.
\end{lemma}

\begin{proof}
Considérons l'ensemble
$$
\mathcal{T} =
\left\{
\begin{matrix}
T \subset |X|
\text{ fermé non vide tel qu'il existe un faisceau cohérent } \\
\mathcal{F}
\text{ avec }
\text{Supp}(\mathcal{F}) = T
\text{ pour lequel le lemme est faux}
\end{matrix}
\right\}
$$
Nous voulons montrer que $\mathcal{T}$ est vide. Sinon,
puisque $|X|$ est noethérien (Propriétés des espaces, Lemme
\ref{spaces-properties-lemma-Noetherian-topology})
on peut choisir un élément minimal $T \in \mathcal{T}$. Cela signifie qu'il
existe un faisceau cohérent $\mathcal{F}$ sur $X$, de support $T$,
pour lequel le lemme n'est pas vrai.

\medskip\noindent
Si $T$ n'est pas irréductible, on peut écrire $T = Z_1 \cup Z_2$,
où $Z_1, Z_2$ sont fermés et strictement contenus dans $T$.
On peut alors appliquer le Lemme \ref{lemma-prepare-filter-support}
pour obtenir une suite exacte courte de faisceaux cohérents
$$
0 \to
\mathcal{G}_1 \to
\mathcal{F} \to
\mathcal{G}_2 \to 0
$$
avec $\text{Supp}(\mathcal{G}_i) \subset Z_i$. Par minimalité de
$T$, chacun des $\mathcal{G}_i$ vérifie $\mathcal{P}$. Ainsi $\mathcal{F}$
vérifie la propriété $\mathcal{P}$ d'après (1), ce qui est une contradiction.

\medskip\noindent
Supposons $T$ irréductible. Soit $\mathcal{J}$ le faisceau d'idéaux
définissant sur $T$ la structure réduite induite de sous-espace fermé;
voir Propriétés des espaces, Lemme
\ref{spaces-properties-lemma-reduced-closed-subspace}.
D'après le Lemme \ref{lemma-power-ideal-kills-sheaf}, il existe
un $n \geq 0$ tel que $\mathcal{J}^n\mathcal{F} = 0$. On obtient donc
une filtration
$$
0 = \mathcal{J}^n\mathcal{F} \subset \mathcal{J}^{n - 1}\mathcal{F}
\subset \ldots \subset \mathcal{J}\mathcal{F} \subset \mathcal{F}
$$
dont chacun des sous-quotients successifs est annulé par $\mathcal{J}$.
Si chacun de ces sous-quotients admet une filtration comme dans l'énoncé
du lemme, il en va de même pour $\mathcal{F}$ d'après (1). Autrement dit, on peut
supposer que $\mathcal{J}$ annule $\mathcal{F}$.

\medskip\noindent
Supposons $T$ irréductible et $\mathcal{J}\mathcal{F} = 0$, où
$\mathcal{J}$ est comme ci-dessus. Notons $i : Z \to X$ le sous-espace fermé
correspondant à $\mathcal{J}$. Alors $\mathcal{F} = i_*\mathcal{H}$
pour un certain $\mathcal{O}_Z$-module cohérent $\mathcal{H}$; voir
Morphismes d'espaces, Lemme \ref{spaces-morphisms-lemma-i-star-equivalence},
et Lemme \ref{lemma-coherent-support-closed}.
Soit $\mathcal{G}$ le faisceau cohérent sur $Z$ vérifiant
(3)(a) et (3)(b). Appliquons le Lemme \ref{lemma-prepare-filter-irreducible}
pour obtenir des morphismes injectifs
$$
\mathcal{I}_1^{\oplus r_1} \to \mathcal{H}
\quad\text{et}\quad
\mathcal{I}_2^{\oplus r_2} \to \mathcal{G}
$$
dont les conoyaux sont à support dans des parties fermées strictes de $Z$. On trouve donc
un ouvert non vide $V \subset Z$ tel que
$$
\mathcal{H}^{\oplus r_2}_V \cong \mathcal{G}^{\oplus r_1}_V
$$
Soit $\mathcal{I} \subset \mathcal{O}_Z$ un faisceau quasi-cohérent d'idéaux
définissant $Z \setminus V$; d'après
le Lemme \ref{lemma-homs-over-open}, on obtient
un morphisme
$$
\mathcal{I}^n\mathcal{G}^{\oplus r_1} \longrightarrow \mathcal{H}^{\oplus r_2}
$$
qui est un isomorphisme sur $V$. Le noyau est à support dans $Z \setminus V$,
donc est annulé par une puissance de $\mathcal{I}$; voir
Lemme \ref{lemma-power-ideal-kills-sheaf}. Ainsi, quitte à augmenter
$n$, on peut supposer que le morphisme affiché est injectif; voir
Lemme \ref{lemma-Artin-Rees}. En appliquant (3)(b), on trouve
$\mathcal{G}' \subset \mathcal{I}^n\mathcal{G}$ tel que
$$
(i_*\mathcal{G}')^{\oplus r_1} \longrightarrow
i_*\mathcal{H}^{\oplus r_2} = \mathcal{F}^{\oplus r_2}
$$
soit injectif, avec un conoyau à support dans une partie fermée stricte de $Z$,
et que la propriété $\mathcal{P}$ soit vérifiée par $i_*\mathcal{G}'$.
D'après (1), la propriété $\mathcal{P}$ est vérifiée par $(i_*\mathcal{G}')^{\oplus r_1}$.
D'après (1) et la minimalité de $T = |Z|$, la propriété $\mathcal{P}$ est vérifiée par
$\mathcal{F}^{\oplus r_2}$. Enfin, d'après (2), la propriété $\mathcal{P}$
est vérifiée par $\mathcal{F}$, ce qui est la contradiction cherchée.
\end{proof}

\begin{lemma}
\label{lemma-property-higher-rank-cohomological-variant}
Soit $S$ un schéma. Soit $X$ un espace algébrique noethérien sur $S$.
Soit $\mathcal{P}$ une propriété des faisceaux cohérents sur $X$. Supposons que
\begin{enumerate}
\item pour toute suite exacte courte de faisceaux cohérents sur $X$,
si deux des trois termes vérifient la propriété $\mathcal{P}$, le troisième la vérifie aussi;
\item si $\mathcal{P}$ est vérifiée par $\mathcal{F}^{\oplus r}$ pour
un certain $r \geq 1$, elle est vérifiée par $\mathcal{F}$;
\item pour tout sous-espace fermé réduit $i : Z \to X$ tel que
$|Z|$ soit irréductible, il existe un faisceau cohérent $\mathcal{G}$ sur $X$
dont le support schématique est $Z$ et tel que $\mathcal{P}$ soit vérifiée par
$\mathcal{G}$.
\end{enumerate}
Alors la propriété $\mathcal{P}$ est vérifiée par tout faisceau cohérent sur $X$.
\end{lemma}

\begin{proof}
Nous allons montrer que les conditions (1) et (2) du
Lemme \ref{lemma-property-initial} sont vérifiées. C'est clair pour la condition (1).
Pour établir (2), posons
$$
\mathcal{T} =
\left\{
\begin{matrix}
i : Z \to X \text{ sous-espace fermé réduit tel que }|Z|\text{ soit irréductible et}\\
\text{ que }i_*\mathcal{I}\text{ ne vérifie pas }\mathcal{P}
\text{ pour un certain faisceau quasi-cohérent }\mathcal{I} \subset \mathcal{O}_Z
\end{matrix}
\right\}
$$
Si $\mathcal{T}$ n'est pas vide, comme $X$ est noethérien, on peut
trouver un $i : Z \to X$ minimal dans $\mathcal{T}$. Nous allons montrer
que cela conduit à une contradiction.

\medskip\noindent
Soit $\mathcal{G}$ le faisceau de support schématique $Z$ dont
l'existence est postulée dans l'hypothèse (3). Soit
$\varphi : i_*\mathcal{I}^{\oplus r} \to \mathcal{G}$ comme dans le
Lemme \ref{lemma-prepare-filter-irreducible}. Soit
$$
0 = \mathcal{F}_0 \subset \mathcal{F}_1 \subset
\ldots \subset \mathcal{F}_m = \Coker(\varphi)
$$
une filtration comme dans le Lemme \ref{lemma-coherent-filter}. Par minimalité
de $Z$ et d'après l'hypothèse (1), on voit que $\Coker(\varphi)$ vérifie
la propriété $\mathcal{P}$. Puisque $\varphi$ est injectif, l'hypothèse (1)
montre encore que $i_*\mathcal{I}^{\oplus r}$ vérifie la propriété
$\mathcal{P}$. L'hypothèse (2) montre alors que $i_*\mathcal{I}$
vérifie la propriété $\mathcal{P}$.

\medskip\noindent
Enfin, si $\mathcal{J} \subset \mathcal{O}_Z$ est un second faisceau quasi-cohérent
d'idéaux, posons $\mathcal{K} = \mathcal{I} \cap \mathcal{J}$
et considérons les suites exactes courtes
$$
0 \to \mathcal{K} \to \mathcal{I} \to \mathcal{I}/\mathcal{K} \to 0
\quad
\text{et}
\quad
0 \to \mathcal{K} \to \mathcal{J} \to \mathcal{J}/\mathcal{K} \to 0
$$
En raisonnant comme ci-dessus et en utilisant la minimalité de $Z$, on voit que
$i_*\mathcal{I}/\mathcal{K}$ et $i_*\mathcal{J}/\mathcal{K}$
vérifient $\mathcal{P}$. L'hypothèse (1) montre donc que
$i_*\mathcal{K}$, puis $i_*\mathcal{J}$, vérifient $\mathcal{P}$.
Autrement dit, $Z$ n'est pas un élément de $\mathcal{T}$, ce qui est
la contradiction cherchée.
\end{proof}






\section{Limites de modules cohérents}
\label{section-limits}

\noindent
Une limite inductive de modules cohérents (sur un espace algébrique
localement noethérien) n'est en général pas cohérente. Elle est toutefois quasi-cohérente,
car toute limite inductive de modules quasi-cohérents sur un espace algébrique est
quasi-cohérente, voir Propriétés des espaces, Lemme
\ref{spaces-properties-lemma-properties-quasi-coherent}.
Réciproquement, si l'espace algébrique est noethérien, tout module quasi-cohérent
est une limite inductive filtrante de modules cohérents.

\begin{lemma}
\label{lemma-directed-colimit-coherent}
Soit $S$ un schéma. Soit $X$ un espace algébrique noethérien sur $S$.
Tout $\mathcal{O}_X$-module quasi-cohérent est la limite inductive filtrante
de ses sous-modules cohérents.
\end{lemma}

\begin{proof}
Soit $\mathcal{F}$ un $\mathcal{O}_X$-module quasi-cohérent.
Si $\mathcal{G}, \mathcal{H} \subset \mathcal{F}$ sont des
sous-$\mathcal{O}_X$-modules cohérents, alors l'image de
$\mathcal{G} \oplus \mathcal{H} \to \mathcal{F}$ est encore un
sous-$\mathcal{O}_X$-module cohérent qui les contient tous deux
(voir les Lemmes \ref{lemma-coherent-abelian-Noetherian} et
\ref{lemma-coherent-Noetherian-quasi-coherent-sub-quotient}).
On voit ainsi que le système est filtrant.
Il suffit donc de montrer que $\mathcal{F}$ peut s'écrire comme
une limite inductive filtrante de modules cohérents : en prenant alors les
images de ces modules dans $\mathcal{F}$, on conclut qu'il y en a
suffisamment.

\medskip\noindent
Soit $U$ un schéma affine et $U \to X$ un morphisme \'etale surjectif.
Posons $R = U \times_X U$, de sorte que $X = U/R$ comme d'habitude. D'après
Propriétés des espaces, Proposition
\ref{spaces-properties-proposition-quasi-coherent},
on a $\QCoh(\mathcal{O}_X) = \QCoh(U, R, s, t, c)$.
On se ramène donc à établir l'assertion correspondante pour
$\QCoh(U, R, s, t, c)$. Le résultat découle alors du
Lemme plus général \ref{groupoids-lemma-colimit-coherent} du chapitre Groupoïdes.
\end{proof}

\begin{lemma}
\label{lemma-direct-colimit-finite-presentation}
Soit $S$ un schéma. Soit $f : X \to Y$ un morphisme affine d'espaces
algébriques sur $S$, avec $Y$ noethérien. Alors tout
$\mathcal{O}_X$-module quasi-cohérent est une limite inductive filtrante de
$\mathcal{O}_X$-modules de présentation finie.
\end{lemma}

\begin{proof}
Soit $\mathcal{F}$ un $\mathcal{O}_X$-module quasi-cohérent.
Écrivons $f_*\mathcal{F} = \colim \mathcal{H}_i$, où $\mathcal{H}_i$ est
un $\mathcal{O}_Y$-module cohérent, voir le
Lemme \ref{lemma-directed-colimit-coherent}.
D'après le Lemme \ref{lemma-coherent-Noetherian}, les modules $\mathcal{H}_i$
sont des $\mathcal{O}_Y$-modules de présentation finie. Par conséquent,
$f^*\mathcal{H}_i$ est un $\mathcal{O}_X$-module de présentation finie, voir
Propriétés des espaces, section
\ref{spaces-properties-section-properties-modules}.
Nous affirmons que le morphisme
$$
\colim f^*\mathcal{H}_i = f^*f_*\mathcal{F} \to \mathcal{F}
$$
est surjectif puisque $f$ est supposé affine. En effet, choisissons
un schéma $V$ et un morphisme \'etale surjectif $V \to Y$. Posons
$U = X \times_Y V$. Alors $U$ est un schéma, $f' : U \to V$ est affine, et
$U \to X$ est \'etale surjectif. D'après
Propriétés des espaces, Lemme
\ref{spaces-properties-lemma-pushforward-etale-base-change-modules},
on a $f'_*(\mathcal{F}|_U) = f_*\mathcal{F}|_V$, et de même
pour les images réciproques. Ainsi, la restriction de $f^*f_*\mathcal{F} \to \mathcal{F}$
à $U$ est le morphisme
$$
f^*f_*\mathcal{F}|_U = (f')^*(f_*\mathcal{F})|_V) =
(f')^*f'_*(\mathcal{F}|_U) \to \mathcal{F}|_U
$$
qui est surjectif puisque $f'$ est un morphisme affine de schémas.
L'affirmation est donc démontrée.

\medskip\noindent
On conclut que tout module quasi-cohérent sur $X$ est quotient d'une
limite inductive filtrante de modules de présentation finie. En particulier,
$\mathcal{F}$ est le conoyau d'un morphisme
$$
\colim_{j \in J} \mathcal{G}_j \longrightarrow \colim_{i \in I} \mathcal{H}_i
$$
où les $\mathcal{G}_j$ et les $\mathcal{H}_i$ sont de présentation finie. Notons
que, pour tout $j \in I$, il existe $i \in I$ et un morphisme
$\alpha : \mathcal{G}_j \to \mathcal{H}_i$ tels que le diagramme
$$
\xymatrix{
\mathcal{G}_j \ar[r]_\alpha \ar[d] & \mathcal{H}_i \ar[d] \\
\colim_{j \in J} \mathcal{G}_j \ar[r] &
\colim_{i \in I} \mathcal{H}_i
}
$$
soit commutatif, voir le
Lemme \ref{lemma-finite-presentation-quasi-compact-colimit}.
Dans cette situation, $\Coker(\alpha)$ est un
$\mathcal{O}_X$-module de présentation finie muni d'un morphisme
$\Coker(\alpha) \to \mathcal{F}$. Considérons l'ensemble $K$ des
triplets $(i, j, \alpha)$ comme ci-dessus. On pose
$(i, j, \alpha) \leq (i', j', \alpha')$ si et seulement si
$i \leq i'$, $j \leq j'$, et si le diagramme
$$
\xymatrix{
\mathcal{G}_j \ar[r]_\alpha \ar[d] & \mathcal{H}_i \ar[d] \\
\mathcal{G}_{j'} \ar[r]^{\alpha'} &
\mathcal{H}_{i'}
}
$$
est commutatif. Il résulte de ce qui précède que $K$ est
un ensemble ordonné filtrant,
$$
\mathcal{F} = \colim_{(i, j, \alpha) \in K} \Coker(\alpha),
$$
ce qui conclut la démonstration.
\end{proof}






\section{Annulation de la cohomologie}
\label{section-vanishing}

\noindent
Dans cette section, nous montrons qu'un espace algébrique quasi-compact et quasi-séparé
est affine si sa cohomologie supérieure s'annule
pour tout faisceau quasi-cohérent. Nous procédons par une suite de lemmes
qui deviendront tous superflus une fois démontrée la
Proposition \ref{proposition-vanishing-affine}.

\begin{situation}
\label{situation-vanishing}
Ici, $S$ est un schéma et $X$ est un espace algébrique quasi-compact et quasi-séparé
sur $S$ ayant la propriété suivante : pour tout
$\mathcal{O}_X$-module quasi-cohérent $\mathcal{F}$, on a
$H^1(X, \mathcal{F}) = 0$. Posons $A = \Gamma(X, \mathcal{O}_X)$.
\end{situation}

\noindent
Nous voudrions montrer que le morphisme canonique
$$
p : X \longrightarrow \Spec(A)
$$
(voir Propriétés des espaces, Lemme
\ref{spaces-properties-lemma-morphism-to-affine-scheme}) est un isomorphisme.
Si $M$ est un $A$-module, nous notons $M \otimes_A \mathcal{O}_X$
le module quasi-cohérent $p^*\tilde M$.

\begin{lemma}
\label{lemma-vanishing-compute}
Dans la Situation \ref{situation-vanishing}, pour tout $A$-module $M$, on a
$p_*(M \otimes_A \mathcal{O}_X) = \widetilde{M}$ et
$\Gamma(X, M \otimes_A \mathcal{O}_X) = M$.
\end{lemma}

\begin{proof}
L'égalité $p_*(M \otimes_A \mathcal{O}_X) = \widetilde{M}$ résulte
de l'égalité $\Gamma(X, M \otimes_A \mathcal{O}_X) = M$, car
$p_*(M \otimes_A \mathcal{O}_X)$ est un module quasi-cohérent sur
$\Spec(A)$ d'après Morphismes d'espaces, Lemme
\ref{spaces-morphisms-lemma-pushforward}.
Observons que $\Gamma(X, \bigoplus_{i \in I} \mathcal{O}_X) =
\bigoplus_{i \in I} A$ d'après le Lemme \ref{lemma-colimits}. Le
lemme est donc vrai pour les modules libres. Choisissons une suite exacte courte
$F_1 \to F_0 \to M$, où $F_0, F_1$ sont des $A$-modules libres. Puisque
$H^1(X, -)$ est nul, le foncteur des sections globales est exact à droite.
De plus, l'image réciproque $p^*$ est elle aussi exacte à droite. On voit donc
que
$$
\Gamma(X, F_1 \otimes_A \mathcal{O}_X) \to
\Gamma(X, F_0 \otimes_A \mathcal{O}_X) \to
\Gamma(X, M \otimes_A \mathcal{O}_X) \to 0
$$
est exacte. Le résultat en découle.
\end{proof}

\noindent
Le lemme suivant montre que la Situation \ref{situation-vanishing}
est préservée par le changement de base de $X \to \Spec(A)$ par $\Spec(A') \to \Spec(A)$.

\begin{lemma}
\label{lemma-vanishing-base-change}
Dans la Situation \ref{situation-vanishing}.
\begin{enumerate}
\item Pour tout morphisme affine $X' \to X$ d'espaces algébriques, on a
$H^1(X', \mathcal{F}') = 0$ pour tout
$\mathcal{O}_{X'}$-module quasi-cohérent $\mathcal{F}'$.
\item Pour toute $A$-algèbre $A'$, en posant $X' = X \times_{\Spec(A)} \Spec(A')$,
le morphisme $X' \to X$ est affine et $\Gamma(X', \mathcal{O}_{X'}) = A'$.
\end{enumerate}
\end{lemma}

\begin{proof}
L'assertion (1) résulte du Lemme \ref{lemma-affine-vanishing-higher-direct-images}
et de la suite spectrale de Leray (Cohomologie sur les sites, Lemme
\ref{sites-cohomology-lemma-Leray}). Soit $A \to A'$ comme en (2).
Alors $X' \to X$ est affine, car les morphismes affines sont stables par
changement de base (Morphismes d'espaces, Lemme
\ref{spaces-morphisms-lemma-base-change-affine}) et parce que
tout morphisme de schémas affines est affine. L'égalité
$\Gamma(X', \mathcal{O}_{X'}) = A'$ résulte de
$(X' \to X)_*\mathcal{O}_{X'} = A' \otimes_A \mathcal{O}_X$
d'après le Lemme \ref{lemma-affine-base-change}, puis de
$$
\Gamma(X', \mathcal{O}_{X'}) =
\Gamma(X, (X' \to X)_*\mathcal{O}_{X'}) =
\Gamma(X, A' \otimes_A \mathcal{O}_X) = A'
$$
d'après le Lemme \ref{lemma-vanishing-compute}.
\end{proof}

\begin{lemma}
\label{lemma-vanishing-separate-closed}
Dans la Situation \ref{situation-vanishing}, soient $Z_0, Z_1 \subset |X|$
deux parties fermées disjointes. Il existe alors $a \in A$ tel que
$Z_0 \subset V(a)$ et $Z_1 \subset V(a - 1)$.
\end{lemma}

\begin{proof}
On peut, et l'on va, munir $Z_0$, $Z_1$ de la structure d'espace réduit induite
(Propriétés des espaces, Définition
\ref{spaces-properties-definition-reduced-induced-space}); notons
$i_0 : Z_0 \to X$ et $i_1 : Z_1 \to X$ les immersions fermées correspondantes.
Comme $Z_0 \cap Z_1 = \emptyset$, le morphisme canonique de
$\mathcal{O}_X$-modules quasi-cohérents
$$
\mathcal{O}_X
\longrightarrow
i_{0, *}\mathcal{O}_{Z_0} \oplus i_{1, *}\mathcal{O}_{Z_1}
$$
est surjectif (il suffit de regarder les germes aux points géométriques). Puisque $H^1(X, -)$
est nul sur le noyau de ce morphisme, le morphisme induit sur les sections globales est
surjectif. On peut donc trouver $a \in A$ dont l'image est la section globale
$(0, 1)$ du membre de droite.
\end{proof}

\begin{lemma}
\label{lemma-vanishing-injective}
Dans la Situation \ref{situation-vanishing}, le morphisme $p : X \to \Spec(A)$ est
universellement injectif.
\end{lemma}

\begin{proof}
Soit $A \to k$ un homomorphisme d'anneaux, où $k$ est un corps. Il suffit de
montrer que $\Spec(k) \times_{\Spec(A)} X$ a au plus un point (voir
Morphismes d'espaces, Lemme
\ref{spaces-morphisms-lemma-universally-injective-local}).
En utilisant le Lemme \ref{lemma-vanishing-base-change}, on peut supposer que $A$
est un corps, et il faut montrer que $|X|$ a au plus un point.

\medskip\noindent
Considérons $X$ comme un espace algébrique sur $\Spec(k)$ et utilisons
la notation $X(K)$ pour désigner les $K$-points de $X$
pour toute extension $K/k$, voir
Morphismes d'espaces, section \ref{spaces-morphisms-section-points-fields}.
Si $K/k$ est une extension algébriquement close
de degré de transcendance assez grand, alors $X(K) \to |X|$
est surjective, voir Morphismes d'espaces, Lemme
\ref{spaces-morphisms-lemma-large-enough}. Après avoir remplacé $k$
par $K$, il suffit donc de prouver que $X(k)$ est réduit à un seul élément
(dans le cas $A = k)$.

\medskip\noindent
Soient $x, x' \in X(k)$. D'après Espaces décents, Lemme
\ref{decent-spaces-lemma-algebraic-residue-field-extension-closed-point},
$x$ et $x'$ sont des points fermés de $|X|$. Ainsi $x$ et $x'$
auraient des images distinctes dans $\Spec(k)$ si $x \not = x'$, d'après le
Lemme \ref{lemma-vanishing-separate-closed}. On conclut que
$x = x'$, comme voulu.
\end{proof}

\begin{lemma}
\label{lemma-vanishing-separated}
Dans la Situation \ref{situation-vanishing}, le morphisme $p : X \to \Spec(A)$ est
séparé.
\end{lemma}

\begin{proof}
D'après Espaces décents, Lemme
\ref{decent-spaces-lemma-there-is-a-scheme-integral-over},
il existe un schéma $Y$ et un morphisme entier surjectif
$Y \to X$. Comme un morphisme entier est affine, on peut appliquer le
Lemme \ref{lemma-vanishing-base-change}
pour voir que $H^1(Y, \mathcal{G}) = 0$ pour tout
$\mathcal{O}_Y$-module quasi-cohérent $\mathcal{G}$.
Puisque $Y \to X$ est quasi-compact et que $X$ est quasi-compact,
$Y$ est quasi-compact.
Comme $Y$ est un schéma, on peut appliquer
Cohomologie des schémas, Lemme
\ref{coherent-lemma-quasi-compact-h1-zero-covering}
pour voir que $Y$ est affine. Ainsi $Y$ est séparé.
Notons qu'un morphisme entier est affine et universellement fermé, voir
Morphismes d'espaces, Lemme
\ref{spaces-morphisms-lemma-integral-universally-closed}.
D'après Morphismes d'espaces, Lemme
\ref{spaces-morphisms-lemma-image-universally-closed-separated},
$X$ est un espace algébrique séparé.
\end{proof}

\begin{proposition}
\label{proposition-vanishing-affine}
\begin{slogan}
Le critère d'affinité de Serre pour les espaces algébriques.
\end{slogan}
Un espace algébrique quasi-compact et quasi-séparé est affine
si et seulement si tous les groupes de cohomologie supérieure des faisceaux quasi-cohérents
s'annulent. Plus précisément, tout espace algébrique comme dans la
Situation \ref{situation-vanishing} est un schéma affine.
\end{proposition}

\begin{proof}
Choisissons un schéma affine $U = \Spec(B)$ et un morphisme \'etale surjectif
$\varphi : U \to X$. Posons $R = U \times_X U$. Comme $p$ est séparé
(Lemme \ref{lemma-vanishing-separated}), $R$ est un
sous-schéma fermé de $U \times_{\Spec(A)} U = \Spec(B \otimes_A B)$.
Ainsi $R = \Spec(C)$ est lui aussi affine, et l'homomorphisme d'anneaux
$$
B \otimes_A B \longrightarrow C
$$
est surjectif. Notons comme d'habitude $s, t : B \to C$ les deux homomorphismes. Choisissons
$g_1, \ldots, g_m \in B$ tels que $s(g_1), \ldots, s(g_m)$ engendrent $C$
sur $t : B \to C$ (ce qui est possible puisque $t : B \to C$ est de présentation
finie et que l'homomorphisme affiché est surjectif). Alors $g_1, \ldots, g_m$
définissent des sections globales de $\varphi_*\mathcal{O}_U$, et le morphisme
$$
\mathcal{O}_X[z_1, \ldots, z_n] \longrightarrow \varphi_*\mathcal{O}_U,
\quad
z_j \longmapsto g_j
$$
est surjectif : on peut le vérifier après restriction à $U$.
En effet, $\varphi^*\varphi_*\mathcal{O}_U = t_*\mathcal{O}_R$
(d'après le Lemme \ref{lemma-flat-base-change-cohomology});
on obtient donc exactement la condition que les $s(g_i)$ engendrent $C$
sur $t : B \to C$. L'annulation de $H^1$ du noyau montre alors que
$$
\Gamma(X, \mathcal{O}_X[x_1, \ldots, x_n]) =
A[x_1, \ldots, x_n] \longrightarrow
\Gamma(X, \varphi_*\mathcal{O}_U) = \Gamma(U, \mathcal{O}_U) = B
$$
est surjectif. On conclut que $B$ est une $A$-algèbre de type fini.
Ainsi $X \to \Spec(A)$ est de type fini et séparé.
D'après le Lemme \ref{lemma-vanishing-injective}
et le
Lemme \ref{spaces-morphisms-lemma-locally-quasi-finite} de Morphismes d'espaces,
il est aussi localement quasi-fini. Par conséquent, $X \to \Spec(A)$ est représentable d'après
Morphismes d'espaces, Lemme
\ref{spaces-morphisms-lemma-locally-quasi-finite-separated-representable},
et $X$ est un schéma. Enfin, $X$ est affine, donc égal à $\Spec(A)$,
par application de Cohomologie des schémas, Lemme
\ref{coherent-lemma-quasi-compact-h1-zero-covering}.
\end{proof}

\begin{lemma}
\label{lemma-Noetherian-h1-zero}
Soit $S$ un schéma. Soit $X$ un espace algébrique noethérien sur $S$.
Supposons que, pour tout $\mathcal{O}_X$-module cohérent
$\mathcal{F}$, on ait $H^1(X, \mathcal{F}) = 0$.
Alors $X$ est un schéma affine.
\end{lemma}

\begin{proof}
L'hypothèse entraîne que $H^1(X, \mathcal{F}) = 0$ pour tout
$\mathcal{O}_X$-module quasi-cohérent $\mathcal{F}$, d'après les
Lemmes \ref{lemma-directed-colimit-coherent} et \ref{lemma-colimits}.
Alors $X$ est affine d'après la
Proposition \ref{proposition-vanishing-affine}.
\end{proof}

\begin{lemma}
\label{lemma-Noetherian-h1-zero-invertible}
Soit $S$ un schéma. Soit $X$ un espace algébrique noethérien sur $S$.
Soit $\mathcal{L}$ un $\mathcal{O}_X$-module inversible.
Supposons que, pour tout $\mathcal{O}_X$-module cohérent
$\mathcal{F}$, il existe $n \geq 1$ tel que
$H^1(X, \mathcal{F} \otimes_{\mathcal{O}_X} \mathcal{L}^{\otimes n}) = 0$.
Alors $X$ est un schéma et $\mathcal{L}$ est ample sur $X$.
\end{lemma}

\begin{proof}
Soit $s \in H^0(X, \mathcal{L}^{\otimes d})$ une section globale.
Soit $U \subset X$ le sous-espace ouvert sur lequel $s$ engendre
$\mathcal{L}^{\otimes d}$. En particulier, on a
$\mathcal{L}^{\otimes d}|_U \cong \mathcal{O}_U$.
Nous affirmons que $U$ est affine.

\medskip\noindent
Démonstration de l'affirmation. Montrons que $H^1(U, \mathcal{F}) = 0$
pour tout $\mathcal{O}_U$-module quasi-cohérent $\mathcal{F}$.
Cela démontrera l'affirmation d'après la Proposition \ref{proposition-vanishing-affine}.
Notons $j : U \to X$ le morphisme d'inclusion.
Comme le morphisme $j$ est \'etale-localement affine
(d'après Morphismes, Lemme \ref{morphisms-lemma-affine-s-open}),
on voit que $j$ est affine (Morphismes d'espaces, Lemme
\ref{spaces-morphisms-lemma-affine-local}).
On a donc
$$
H^1(U, \mathcal{F}) = H^1(X, j_*\mathcal{F})
$$
d'après le Lemme \ref{lemma-affine-vanishing-higher-direct-images}
(et Cohomologie sur les sites, Lemme \ref{sites-cohomology-lemma-apply-Leray}).
Écrivons $j_*\mathcal{F} = \colim \mathcal{F}_i$ comme limite inductive filtrante
de $\mathcal{O}_X$-modules cohérents, voir le
Lemme \ref{lemma-directed-colimit-coherent}. Alors
$$
H^1(X, j_*\mathcal{F}) = \colim H^1(X, \mathcal{F}_i)
$$
d'après le Lemme \ref{lemma-colimits}.
Il suffit donc de montrer que $H^1(X, \mathcal{F}_i)$ s'envoie
sur zéro dans $H^1(U, j^*\mathcal{F}_i)$. Par hypothèse, il existe
$n \geq 1$ tel que
$$
H^1(X,
\mathcal{F}_i \otimes_{\mathcal{O}_X}
(\mathcal{O}_X \oplus \mathcal{L} \oplus \ldots
\oplus \mathcal{L}^{\otimes d - 1})
\otimes_{\mathcal{O}_X} \mathcal{L}^{\otimes n}) = 0
$$
Il existe donc $a \geq 0$ tel que
$H^1(X, \mathcal{F}_i \otimes_{\mathcal{O}_X} \mathcal{L}^{\otimes ad}) = 0$.
D'autre part, le morphisme
$$
s^a : \mathcal{F}_i \longrightarrow
\mathcal{F}_i \otimes_{\mathcal{O}_X} \mathcal{L}^{\otimes ad}
$$
est un isomorphisme après restriction à $U$. En considérant le
diagramme commutatif
$$
\xymatrix{
H^1(X, \mathcal{F}_i) \ar[r] \ar[d]_{s^a} & H^1(U, j^*\mathcal{F}_i)
\ar[d]^{\cong} \\
H^1(X, \mathcal{F}_i \otimes_{\mathcal{O}_X} \mathcal{L}^{\otimes ad}) \ar[r] &
H^1(U,
j^*(\mathcal{F}_i \otimes_{\mathcal{O}_X} \mathcal{L}^{\otimes ad}))
}
$$
on conclut que le morphisme $H^1(X, \mathcal{F}_i) \to H^1(U, j^*\mathcal{F}_i)$
est nul, et l'affirmation est démontrée.

\medskip\noindent
Soit $x \in |X|$ un point fermé. D'après Espaces décents, Lemme
\ref{decent-spaces-lemma-decent-space-closed-point},
on peut représenter $x$ par une immersion fermée $i : \Spec(k) \to X$
(cela utilise aussi le fait qu'un espace algébrique quasi-séparé est
décent, voir Espaces décents, section
\ref{decent-spaces-section-reasonable-decent}).
Ainsi $\mathcal{O}_X \to i_*\mathcal{O}_{\Spec(k)}$ est surjectif.
Soit $\mathcal{I} \subset \mathcal{O}_X$ son noyau, et choisissons
$d \geq 1$ tel que
$H^1(X, \mathcal{I} \otimes_{\mathcal{O}_X} \mathcal{L}^{\otimes d}) = 0$.
Alors
$$
H^0(X, \mathcal{L}^{\otimes d}) \to
H^0(X,
i_*\mathcal{O}_{\Spec(k)} \otimes_{\mathcal{O}_X} \mathcal{L}^{\otimes d}) =
H^0(\Spec(k), i^*\mathcal{L}^{\otimes d}) \cong k
$$
est surjectif d'après la suite exacte longue de cohomologie. Il existe donc
$s \in H^0(X, \mathcal{L}^{\otimes d})$
tel que $x \in U$, où $U$ est le sous-espace ouvert associé à $s$
comme ci-dessus. D'après notre affirmation, $x$ appartient donc au lieu schématique
(voir Propriétés des espaces, Lemme \ref{spaces-properties-lemma-subscheme})
de $X$.

\medskip\noindent
Pour conclure que $X$ est un schéma, il suffit de montrer que
toute partie ouverte de $|X|$ qui contient tous les points fermés
est égale à $|X|$. Cela résulte du fait que $|X|$
est un espace topologique noethérien, voir
Propriétés des espaces, Lemme \ref{spaces-properties-lemma-Noetherian-sober}.
Enfin, lorsque $X$ est un schéma, on peut appliquer
Cohomologie des schémas, Lemme
\ref{coherent-lemma-quasi-compact-h1-zero-invertible}
pour conclure que $\mathcal{L}$ est ample.
\end{proof}









\section{Morphismes finis et espaces affines}
\label{section-finite-affine}

\noindent
Cette section est l'analogue de la
section \ref{coherent-section-finite-affine} du chapitre Cohomologie des schémas.

\begin{lemma}
\label{lemma-finite-morphism-Noetherian}
Soit $S$ un schéma. Soit $f : Y \to X$ un morphisme d'espaces algébriques
sur $S$. Supposons $f$ fini et surjectif, et $X$ localement noethérien.
Soit $i : Z \to X$ une immersion fermée. Notons $i' : Z' \to Y$
l'image réciproque de $Z$
(Morphismes d'espaces, section \ref{spaces-morphisms-section-closed-immersions})
et $f' : Z' \to Z$
le morphisme induit. Alors $\mathcal{G} = f'_*\mathcal{O}_{Z'}$
est un $\mathcal{O}_Z$-module cohérent dont le support est $Z$.
\end{lemma}

\begin{proof}
Observons que $f'$ est le changement de base de $f$; il est donc fini
et surjectif d'après Morphismes d'espaces, Lemmes
\ref{spaces-morphisms-lemma-base-change-surjective} et
\ref{spaces-morphisms-lemma-base-change-integral}.
Notons que $Y$, $Z$ et $Z'$ sont localement noethériens d'après
Morphismes d'espaces, Lemme
\ref{spaces-morphisms-lemma-locally-finite-type-locally-noetherian}
(et parce que les immersions fermées et les morphismes finis
sont de type fini). D'après le Lemme \ref{lemma-finite-pushforward-coherent},
$\mathcal{G}$ est un $\mathcal{O}_Z$-module cohérent.
Le support de $\mathcal{G}$ est fermé dans $|Z|$, voir
Morphismes d'espaces, Lemme
\ref{spaces-morphisms-lemma-support-finite-type}.
Si le support de $\mathcal{G}$ n'était pas égal à $|Z|$,
on pourrait donc, après avoir remplacé $X$ par un sous-espace ouvert, supposer
$\mathcal{G} = 0$ mais $Z \not = \emptyset$.
On aurait alors $f'_*\mathcal{O}_{Z'} = 0$. En particulier, la section
$1 \in \Gamma(Z', \mathcal{O}_{Z'}) = \Gamma(Z, f'_*\mathcal{O}_{Z'})$
serait nulle, ce qui entraînerait que $Z' = \emptyset$ est l'espace
algébrique vide. C'est impossible puisque $Z' \to Z$ est surjectif.
\end{proof}

\begin{lemma}
\label{lemma-affine-morphism-projection-ideal}
Soit $S$ un schéma. Soit $f : Y \to X$ un morphisme d'espaces algébriques
sur $S$. Soit $\mathcal{F}$ un faisceau quasi-cohérent sur $Y$.
Soit $\mathcal{I}$ un faisceau quasi-cohérent d'idéaux sur $X$.
Si $f$ est affine, alors
$\mathcal{I}f_*\mathcal{F} = f_*(f^{-1}\mathcal{I}\mathcal{F})$
(avec les notations expliquées dans la démonstration).
\end{lemma}

\begin{proof}
Les notations signifient ce qui suit. Puisque $f^{-1}$ est un foncteur exact,
$f^{-1}\mathcal{I}$ est un faisceau
d'idéaux de $f^{-1}\mathcal{O}_X$. Par le morphisme
$f^\sharp : f^{-1}\mathcal{O}_X \to \mathcal{O}_Y$
sur $Y_\etale$, il opère sur $\mathcal{F}$.
Alors $f^{-1}\mathcal{I}\mathcal{F}$ est le sous-faisceau
engendré par les sommes de sections locales de la forme $as$, où $a$
est une section locale de $f^{-1}\mathcal{I}$ et $s$ une section locale
de $\mathcal{F}$. C'est un sous-$\mathcal{O}_Y$-module quasi-cohérent
de $\mathcal{F}$, car c'est aussi l'image d'un morphisme naturel
$f^*\mathcal{I} \otimes_{\mathcal{O}_Y} \mathcal{F} \to \mathcal{F}$.

\medskip\noindent
Cela étant, la démonstration est immédiate. En effet, la question est
\'etale-locale sur $X$, et l'on peut donc supposer que $X$ est un schéma affine.
Dans ce cas, le résultat découle du résultat correspondant
pour les schémas, voir Cohomologie des schémas, Lemme
\ref{coherent-lemma-affine-morphism-projection-ideal}.
\end{proof}

\begin{lemma}
\label{lemma-image-affine-finite-morphism-affine-Noetherian}
Soit $S$ un schéma. Soit $f : Y \to X$ un morphisme d'espaces
algébriques sur $S$. Supposons que
\begin{enumerate}
\item $f$ soit fini,
\item $f$ soit surjectif,
\item $Y$ soit affine, et
\item $X$ soit noethérien.
\end{enumerate}
Alors $X$ est affine.
\end{lemma}

\begin{proof}
Nous allons montrer que, sous les hypothèses du lemme, pour tout
$\mathcal{O}_X$-module cohérent $\mathcal{F}$, on a $H^1(X, \mathcal{F}) = 0$.
Cela entraîne que $H^1(X, \mathcal{F}) = 0$ pour tout
$\mathcal{O}_X$-module quasi-cohérent $\mathcal{F}$ d'après les
Lemmes \ref{lemma-directed-colimit-coherent} et \ref{lemma-colimits}.
Il s'ensuit alors que $X$ est affine d'après la
Proposition \ref{proposition-vanishing-affine}.

\medskip\noindent
Soit $\mathcal{P}$ la propriété des faisceaux cohérents
$\mathcal{F}$ sur $X$ définie par
$$
\mathcal{P}(\mathcal{F}) \Leftrightarrow H^1(X, \mathcal{F}) = 0.
$$
Nous allons appliquer le Lemme \ref{lemma-property-higher-rank-cohomological}.
Il faut donc vérifier les propriétés (1), (2) et (3) de ce lemme pour $\mathcal{P}$.
La propriété (1) résulte de la suite exacte longue de cohomologie associée
à une suite exacte courte de faisceaux. La propriété (2) résulte du fait que
$H^1(X, -)$ est un foncteur additif. Pour vérifier (3), soit $i : Z \to X$
un sous-espace fermé réduit tel que $|Z|$ soit irréductible. Soient $i' : Z' \to Y$
et $f' : Z' \to Z$ comme dans le Lemme \ref{lemma-finite-morphism-Noetherian},
et posons $\mathcal{G} = f'_*\mathcal{O}_{Z'}$. Nous affirmons que
$\mathcal{G}$ vérifie les propriétés (3)(a) et (3)(b) du
Lemme \ref{lemma-property-higher-rank-cohomological},
ce qui achèvera la démonstration. La propriété (3)(a) a été établie dans le
Lemme \ref{lemma-finite-morphism-Noetherian}. Pour vérifier (3)(b), soit
$\mathcal{I}$ un faisceau quasi-cohérent non nul d'idéaux sur $Z$.
Notons $\mathcal{I}' \subset \mathcal{O}_{Z'}$ l'idéal quasi-cohérent
$(f')^{-1}\mathcal{I} \mathcal{O}_{Z'}$, c'est-à-dire l'
image de $(f')^*\mathcal{I} \to \mathcal{O}_{Z'}$.
D'après le Lemme \ref{lemma-affine-morphism-projection-ideal}, on a
$f_*\mathcal{I}' = \mathcal{I} \mathcal{G}$.
Nous affirmons que la valeur commune
$\mathcal{G}' = \mathcal{I} \mathcal{G} = f'_*\mathcal{I}'$
satisfait la condition énoncée en (3)(b).
Tout d'abord, il est clair que le support de $\mathcal{G}/\mathcal{G}'$
est contenu dans le support de $\mathcal{O}_Z/\mathcal{I}$,
qui est un sous-espace propre de $|Z|$, puisque $\mathcal{I}$ est un
faisceau d'idéaux non nul sur l'espace algébrique réduit et irréductible $Z$.
Le morphisme $f'$ est affine, donc $R^1f'_*\mathcal{I}' = 0$ d'après le
Lemme \ref{lemma-affine-vanishing-higher-direct-images}.
Comme $Z'$ est affine (en tant que sous-schéma fermé d'un schéma affine),
on a $H^1(Z', \mathcal{I}') = 0$. La suite spectrale de Leray
(sous la forme du
Lemme \ref{sites-cohomology-lemma-apply-Leray} de Cohomologie sur les sites)
entraîne donc que $H^1(Z, f'_*\mathcal{I}') = 0$.
Comme $i : Z \to X$ est affine, on conclut que
$R^1i_*f'_*\mathcal{I}' = 0$, puis que $H^1(X, i_*f'_*\mathcal{I}') = 0$
par une nouvelle application de Leray. Autrement dit, $H^1(X, i_*\mathcal{G}') = 0$,
comme voulu.
\end{proof}








\section{Une version faible du lemme de Chow}
\label{section-weak-chow}

\noindent
Dans cette section, nous démontrons rapidement le lemme suivant afin
de faciliter la démonstration des résultats fondamentaux sur la cohomologie des
modules cohérents sur les espaces algébriques propres.

\begin{lemma}
\label{lemma-weak-chow}
Soit $A$ un anneau. Soit $X$ un espace algébrique sur $\Spec(A)$
dont le morphisme structural $X \to \Spec(A)$ est séparé et de type fini.
Il existe alors un morphisme propre surjectif $X' \to X$,
où $X'$ est un schéma H-quasi-projectif sur $\Spec(A)$.
\end{lemma}

\begin{proof}
Soit $W$ un schéma affine et soit $f : W \to X$ un morphisme
\'etale surjectif. Il existe un entier $d$ tel que toutes les fibres géométriques
de f aient $\leq d$ points (car $X$ est un espace algébrique séparé,
donc raisonnable, voir
Espaces décents, Lemme \ref{decent-spaces-lemma-bounded-fibres}).
En choisissant $d$ minimal, on obtient un ouvert non vide $U \subset X$ tel que
$f^{-1}(U) \to U$ soit fini \'etale de degré $d$, voir
Espaces décents, Lemme
\ref{decent-spaces-lemma-quasi-compact-reasonable-stratification}.
Soit
$$
V \subset W \times_X W \times_X \ldots \times_X W
$$
(avec $d$ facteurs dans le produit fibré) le complémentaire de toutes les diagonales.
Comme $W \to X$ est séparé, la diagonale $W \to W \times_X W$ est une
immersion fermée. Comme $W \to X$ est \'etale, la diagonale
$W \to W \times_X W$ est une immersion ouverte, voir
Morphismes d'espaces, Lemmes
\ref{spaces-morphisms-lemma-etale-unramified} et
\ref{spaces-morphisms-lemma-diagonal-unramified-morphism}.
Les diagonales sont donc des sous-schémas ouverts et fermés
du schéma quasi-compact $W \times_X \ldots \times_X W$.
On conclut en particulier que $V$ est un schéma quasi-compact.
Choisissons une immersion ouverte $W \subset Y$, où $Y$ est H-projectif sur
$A$ (c'est possible puisque $W$ est affine et de type fini sur $A$;
on peut par exemple utiliser
Morphismes, Lemmes
\ref{morphisms-lemma-quasi-affine-finite-type-over-S} et
\ref{morphisms-lemma-H-quasi-projective-open-H-projective}).
Soit
$$
Z \subset Y \times_A Y \times_A \ldots \times_A Y
$$
l'image schématique du composé
$V \to W \times_X \ldots \times_X W \to Y \times_A \ldots \times_A Y$.
Observons que ce morphisme est quasi-compact, puisque $V$ est quasi-compact
et que $Y \times_A \ldots \times_A Y$ est séparé.
Notons que $V \to Z$ est une immersion ouverte, car
$V \to Y \times_A \ldots \times_A Y$ est une immersion, voir
Morphismes, Lemme \ref{morphisms-lemma-quasi-compact-immersion}.
Les morphismes de projection donnent $d$ morphismes $g_i : Z \to Y$.
Ces morphismes $g_i$ sont projectifs puisque $Y$ est projectif sur $A$, voir
les résultats de Morphismes, section \ref{morphisms-section-projective}.
Posons
$$
X' = \bigcup g_i^{-1}(W) \subset Z
$$
Il existe un morphisme $X' \to X$ dont la restriction à $g_i^{-1}(W)$ est
le composé $g_i^{-1}(W) \to W \to X$.
En effet, ces morphismes coïncident sur $V$, donc sur
$g_i^{-1}(W) \cap g_j^{-1}(W)$ d'après
Morphismes d'espaces, Lemme \ref{spaces-morphisms-lemma-equality-of-morphisms}.
Affirmation : le morphisme $X' \to X$ est propre.

\medskip\noindent
Si l'affirmation est vraie, le lemme résulte d'une récurrence sur $d$.
En effet, par construction, $X'$ est H-quasi-projectif sur $\Spec(A)$.
L'image de $X' \to X$ contient l'ouvert $U$, puisque $V$ se projette surjectivement sur $U$.
Notons $T$ l'espace algébrique muni de la structure réduite induite sur $X \setminus U$.
Alors $T \times_X W$ est un sous-schéma fermé de $W$, donc est affine.
De plus, le morphisme $T \times_X W \to T$ est \'etale et toute fibre géométrique
a $< d$ points. Par hypothèse de récurrence, il existe un morphisme propre
surjectif $T' \to T$, où $T'$ est un schéma H-quasi-projectif
sur $\Spec(A)$. Comme $T$ est un sous-espace fermé de $X$,
$T' \to X$ est un morphisme propre. Le lemme résulte donc du
morphisme propre surjectif $X' \amalg T' \to X$.

\medskip\noindent
Démonstration de l'affirmation. Par construction, le morphisme $X' \to X$ est séparé
et de type fini. Nous allons vérifier les conditions (1) -- (4) de
Morphismes d'espaces, Lemme
\ref{spaces-morphisms-lemma-refined-valuative-criterion-universally-closed},
pour les morphismes $V \to X'$ et $X' \to X$.
Les conditions (1) et (2) ont été établies ci-dessus.
La condition (3) est satisfaite parce que $X' \to X$ est séparé (son
espace algébrique source étant séparé). Il suffit donc de vérifier
l'existence d'un relèvement dans $X'$ pour les diagrammes
$$
\xymatrix{
\Spec(K) \ar[r] \ar[d] & V \ar[d] \\
\Spec(R) \ar[r] & X
}
$$
où $R$ est un anneau de valuation de corps des fractions $K$.
Notons que le morphisme horizontal supérieur est donné par $d$ points à valeurs dans
$K$, deux à deux distincts, $w_1, \ldots, w_d$, de $W$. En fait, ils
sont exactement les antécédents du point $x \in X(K)$
provenant du diagramme. Comme $W \to X$ est surjectif,
on peut, après avoir éventuellement remplacé $R$ par une extension d'anneaux de valuation,
relever le morphisme $\Spec(R) \to X$ en un morphisme $w : \Spec(R) \to W$, voir
Morphismes d'espaces, Lemme
\ref{spaces-morphisms-lemma-lift-valuation-ring-through-flat-morphism}.
Puisque $w_1, \ldots, w_d$ forment l'ensemble de tous les antécédents de
$x$, la restriction $w|_{\Spec(K)}$ est égale à l'un d'eux, disons $w_i$.
On obtient ainsi un diagramme commutatif
$$
\xymatrix{
\Spec(K) \ar[r] \ar[d] &  Z \ar[d]_{g_i}\\
\Spec(R) \ar[r]^w & Y
}
$$
D'après le critère valuatif de propreté appliqué au morphisme projectif
$g_i$, on peut relever $w$ en $z : \Spec(R) \to Z$, voir
Morphismes, Lemme \ref{morphisms-lemma-locally-projective-proper}
et
Schémas, Proposition \ref{schemes-proposition-characterize-universally-closed}.
L'image de $z$ est contenue dans $g_i^{-1}(W) \subset X'$, ce qui achève la démonstration.
\end{proof}







\section{Critère valuatif noethérien}
\label{section-Noetherian-valuative-criterion}

\noindent
Nous démontrons une version du critère valuatif de propreté
faisant intervenir des anneaux de valuation discrète. On trouvera des versions plus précises
(et donc plus techniques) dans
Limites d'espaces, section
\ref{spaces-limits-section-Noetherian-valuative-criterion}.

\begin{lemma}
\label{lemma-check-separated-dvr}
Soit $S$ un schéma. Soit $f : X \to Y$ un morphisme d'espaces algébriques
sur $S$. Supposons que
\begin{enumerate}
\item $Y$ soit localement noethérien,
\item $f$ soit localement de type fini et quasi-séparé,
\item pour tout diagramme commutatif
$$
\xymatrix{
\Spec(K) \ar[r] \ar[d] & X \ar[d] \\
\Spec(A) \ar[r] \ar@{-->}[ru] & Y
}
$$
où $A$ est un anneau de valuation discrète et $K$ son corps des fractions,
il existe au plus une flèche pointillée rendant le diagramme commutatif.
\end{enumerate}
Alors $f$ est séparé.
\end{lemma}

\begin{proof}
Il faut montrer que la diagonale $\Delta : X \to X \times_Y X$ est une immersion
fermée. On sait déjà que $\Delta$ est représentable, séparée,
un monomorphisme et localement de type fini, voir Morphismes d'espaces, Lemme
\ref{spaces-morphisms-lemma-properties-diagonal}.
Choisissons un schéma affine $U$ et un morphisme \'etale
$U \to X \times_Y X$. Posons $V = X \times_{\Delta, X \times_Y X} U$.
Il suffit de montrer que $V \to U$ est une immersion fermée
(Morphismes d'espaces, Lemme
\ref{spaces-morphisms-lemma-closed-immersion-local}).
Comme $X \times_Y X$ est localement de type fini sur $Y$,
$U$ est noethérien (utiliser Morphismes d'espaces, Lemmes
\ref{spaces-morphisms-lemma-composition-finite-type},
\ref{spaces-morphisms-lemma-base-change-finite-type} et
\ref{spaces-morphisms-lemma-locally-finite-type-locally-noetherian}).
Notons que $V$ est un schéma, puisque $\Delta$ est représentable.
En outre, $V$ est quasi-compact puisque $f$ est quasi-séparé.
Ainsi $V \to U$ est de type fini. Considérons un diagramme commutatif
$$
\xymatrix{
\Spec(K) \ar[r] \ar[d] & V \ar[d] \\
\Spec(A) \ar[r] \ar@{-->}[ru] & U
}
$$
de morphismes de schémas,
où $A$ est un anneau de valuation discrète de corps des fractions $K$.
On peut interpréter le composé $\Spec(A) \to U \to X \times_Y X$
comme une paire de morphismes $a, b : \Spec(A) \to X$ qui coïncident comme morphismes
dans $Y$ et dont les restrictions à $\Spec(K)$ sont égales. L'hypothèse
(3) garantit donc que $a = b$, et l'on obtient la flèche pointillée du diagramme.
D'après Limites, Lemme \ref{limits-lemma-Noetherian-dvr-valuative-proper},
on conclut que $V \to U$ est propre. Autrement dit, $\Delta$ est propre.
Comme $\Delta$ est un monomorphisme, on en déduit que $\Delta$ est une
immersion fermée (Morphismes \'etales, Lemme
\ref{etale-lemma-characterize-closed-immersion}), comme voulu.
\end{proof}

\begin{lemma}
\label{lemma-check-proper-dvr}
Soit $S$ un schéma. Soit $f : X \to Y$ un morphisme d'espaces algébriques
sur $S$. Supposons que
\begin{enumerate}
\item $Y$ soit localement noethérien,
\item $f$ soit de type fini et quasi-séparé,
\item pour tout diagramme commutatif
$$
\xymatrix{
\Spec(K) \ar[r] \ar[d] & X \ar[d] \\
\Spec(A) \ar[r] \ar@{-->}[ru] & Y
}
$$
où $A$ est un anneau de valuation discrète et $K$ son corps des fractions,
il existe une unique flèche pointillée rendant le diagramme commutatif.
\end{enumerate}
Alors $f$ est propre.
\end{lemma}

\begin{proof}
Il suffit de prouver que $f$ est universellement fermé, car $f$ est séparé d'après le
Lemme \ref{lemma-check-separated-dvr}.
Pour cela, on peut travailler \'etale-localement sur $Y$
(Morphismes d'espaces, Lemme
\ref{spaces-morphisms-lemma-universally-closed-local}).
On peut donc supposer que $Y = \Spec(A)$ est un schéma affine noethérien.
Choisissons $X' \to X$ comme dans la version faible du lemme de Chow
(Lemme \ref{lemma-weak-chow}). Nous affirmons que $X' \to \Spec(A)$
est universellement fermé. L'affirmation entraîne le lemme d'après
Morphismes d'espaces, Lemme \ref{spaces-morphisms-lemma-image-proper-is-proper}.
Pour la démontrer, d'après
Limites, Lemme \ref{limits-lemma-check-universally-closed-Noetherian},
il suffit de prouver que, dans tout diagramme commutatif à flèches pleines
$$
\xymatrix{
\Spec(K) \ar[r] \ar[d] & X' \ar[r] & X \ar[d] \\
\Spec(A) \ar[rr] \ar@{-->}[ru]^a \ar@{-->}[rru]_b & & Y
}
$$
où $A$ est un anneau de valuation discrète de corps des fractions $K$, on peut trouver la
flèche pointillée $a$. Par hypothèse, on peut trouver la flèche pointillée $b$.
Alors le morphisme $X' \times_{X, b} \Spec(A) \to \Spec(A)$
est un morphisme propre de schémas et, d'après le critère valuatif
pour les morphismes de schémas, le morphisme $b$ admet le relèvement voulu $a$.
\end{proof}

\begin{remark}[Variante pour les anneaux de valuation discrète complets]
\label{remark-variant}
Dans les Lemmes \ref{lemma-check-separated-dvr} et \ref{lemma-check-proper-dvr},
il suffit de considérer les anneaux de valuation discrète complets.
Plus précisément, dans le Lemme \ref{lemma-check-separated-dvr}, on peut remplacer
la condition (3) par la condition suivante : pour tout diagramme commutatif
$$
\xymatrix{
\Spec(K) \ar[r] \ar[d] & X \ar[d] \\
\Spec(A) \ar[r] \ar@{-->}[ru] & Y
}
$$
où $A$ est un anneau de valuation discrète complet de corps des fractions $K$,
il existe au plus une flèche pointillée rendant le diagramme commutatif. En effet, pour
tout diagramme comme dans le Lemme \ref{lemma-check-separated-dvr} (3),
le complété $A^\wedge$ est un anneau de valuation discrète
(Compléments d'algèbre, Lemme \ref{more-algebra-lemma-completion-dvr}),
et l'unicité de la flèche $\Spec(A^\wedge) \to X$
entraîne celle de la flèche $\Spec(A) \to X$,
par exemple d'après Propriétés des espaces, Proposition
\ref{spaces-properties-proposition-sheaf-fpqc}.
De même, dans le Lemme \ref{lemma-check-proper-dvr},
on peut remplacer la condition (3) par la condition suivante :
pour tout diagramme commutatif
$$
\xymatrix{
\Spec(K) \ar[r] \ar[d] & X \ar[d] \\
\Spec(A) \ar[r] & Y
}
$$
où $A$ est un anneau de valuation discrète complet de corps des fractions $K$,
il existe une extension $A \subset A'$ d'anneaux de valuation discrète complets
induisant une extension des corps des fractions $K \subset K'$ telle qu'il existe une
unique flèche $\Spec(A') \to X$ rendant le diagramme
$$
\xymatrix{
\Spec(K') \ar[r] \ar[d] & \Spec(K) \ar[r] & X \ar[d] \\
\Spec(A') \ar[r] \ar[rru] & \Spec(A) \ar[r] & Y
}
$$
commutatif. En effet, pour tout diagramme comme dans le Lemme \ref{lemma-check-proper-dvr}
partie (3), l'existence d'un diagramme commutatif
$$
\xymatrix{
\Spec(L) \ar[r] \ar[d] & \Spec(K) \ar[r] & X \ar[d] \\
\Spec(B) \ar[r] \ar[rru] & \Spec(A) \ar[r] & Y
}
$$
pour {\it toute} extension $A \subset B$ d'anneaux de valuation discrète
entraîne l'existence d'une flèche $\Spec(A) \to X$ s'insérant dans
le diagramme. Cela a été démontré dans
Morphismes d'espaces, Lemme \ref{spaces-morphisms-lemma-push-down-solution}.
En fait, ces considérations montrent qu'il suffit de chercher
les flèches pointillées dans les diagrammes pour toute classe d'anneaux de valuation discrète
telle que, pour tout anneau de valuation discrète, il en existe une extension
qui appartienne à la classe. On pourrait, par exemple, prendre les anneaux de valuation
discrète complets à corps résiduel algébriquement clos.
\end{remark}







\section{Images directes supérieures des faisceaux cohérents}
\label{section-proper-pushforward}

\noindent
Dans cette section, nous démontrons le fait fondamental que les images
directes supérieures d'un faisceau cohérent par un morphisme propre
sont cohérentes. Commençons par un lemme auxiliaire.

\begin{lemma}
\label{lemma-kill-by-twisting}
Soit $S$ un schéma. Considérons un diagramme commutatif
$$
\xymatrix{
X \ar[r]_i \ar[rd]_f & \mathbf{P}^n_Y \ar[d] \\
& Y
}
$$
d'espaces algébriques au-dessus de $S$. Supposons que $i$ soit une immersion fermée
et que $Y$ soit noethérien. Posons $\mathcal{L} = i^*\mathcal{O}_{\mathbf{P}^n_Y}(1)$.
Soit $\mathcal{F}$ un module cohérent sur $X$.
Il existe alors un entier $d_0$ tel que, pour tout $d \geq d_0$, on ait
$R^pf_*(\mathcal{F} \otimes_{\mathcal{O}_X} \mathcal{L}^{\otimes d}) = 0$
pour tout $p > 0$.
\end{lemma}

\begin{proof}
La nullité de $R^pf_*(\mathcal{F} \otimes \mathcal{L}^{\otimes d})$
se vérifie localement pour la topologie \'etale sur $Y$, voir
l'équation (\ref{equation-representable-higher-direct-image}).
On peut donc supposer que $Y$ est le spectre d'un anneau noethérien. Dans ce cas,
$X$ est un schéma et le résultat découle de
Cohomologie des schémas, lemme \ref{coherent-lemma-kill-by-twisting}.
\end{proof}

\begin{lemma}
\label{lemma-proper-pushforward-coherent}
Soit $S$ un schéma. Soit $f : X \to Y$ un morphisme propre
d'espaces algébriques au-dessus de $S$, avec $Y$ localement noethérien.
Soit $\mathcal{F}$ un $\mathcal{O}_X$-module cohérent.
Alors $R^if_*\mathcal{F}$ est un $\mathcal{O}_Y$-module cohérent
pour tout $i \geq 0$.
\end{lemma}

\begin{proof}
Remarquons d'abord que $X$ est un espace algébrique localement noethérien
d'après Morphismes d'espaces, lemme
\ref{spaces-morphisms-lemma-locally-finite-type-locally-noetherian}.
L'énoncé du lemme a donc un sens. De plus, la formation de
$R^if_*\mathcal{F}$ commute à la localisation \'etale sur $Y$
(Propriétés des espaces, lemme
\ref{spaces-properties-lemma-pushforward-etale-base-change-modules})
et la cohérence de $R^if_*\mathcal{F}$ se vérifie
localement pour la topologie \'etale sur $Y$ (lemme \ref{lemma-coherent-Noetherian}).
On peut donc supposer que $Y = \Spec(A)$ est un schéma affine noethérien.

\medskip\noindent
Supposons que $Y = \Spec(A)$ soit un schéma affine. Notons que $f$ est localement
de présentation finie
(Morphismes d'espaces, lemme
\ref{spaces-morphisms-lemma-noetherian-finite-type-finite-presentation}).
Il est donc de présentation finie, et par suite $X$ est noethérien
(Morphismes d'espaces, lemme
\ref{spaces-morphisms-lemma-finite-presentation-noetherian}).
Le lemme \ref{lemma-property-higher-rank-cohomological-variant}
s'applique donc à la catégorie des modules cohérents sur $X$.
Pour un faisceau cohérent $\mathcal{F}$ sur $X$, disons que $\mathcal{P}$ est vérifiée
si et seulement si $R^if_*\mathcal{F}$ est un module cohérent sur $\Spec(A)$.
Nous allons montrer que les conditions (1), (2) et (3) du
lemme \ref{lemma-property-higher-rank-cohomological-variant} sont vérifiées
pour cette propriété, ce qui achèvera la démonstration du lemme.

\medskip\noindent
Vérification de la condition (1). Soit
$$
0 \to \mathcal{F}_1 \to \mathcal{F}_2 \to \mathcal{F}_3 \to 0
$$
une suite exacte courte de faisceaux cohérents sur $X$.
Considérons la suite exacte longue des images directes supérieures
$$
R^{p - 1}f_*\mathcal{F}_3 \to
R^pf_*\mathcal{F}_1 \to
R^pf_*\mathcal{F}_2 \to
R^pf_*\mathcal{F}_3 \to
R^{p + 1}f_*\mathcal{F}_1
$$
Il est alors clair que, si deux des trois faisceaux $\mathcal{F}_i$
vérifient la propriété $\mathcal{P}$, les images directes supérieures du
troisième sont encadrées, dans ce complexe exact, par deux faisceaux
cohérents. Ces images directes supérieures sont donc elles aussi cohérentes d'après les
lemmes \ref{lemma-coherent-abelian-Noetherian} et
\ref{lemma-coherent-Noetherian-quasi-coherent-sub-quotient}.
Ainsi, la propriété $\mathcal{P}$ est également vérifiée pour le troisième.

\medskip\noindent
Vérification de la condition (2). Elle résulte immédiatement de ce que
$R^if_*(\mathcal{F}_1 \oplus \mathcal{F}_2) =
R^if_*\mathcal{F}_1 \oplus R^if_*\mathcal{F}_2$ et de ce qu'un facteur direct
d'un module cohérent est cohérent (voir les lemmes cités ci-dessus).

\medskip\noindent
Vérification de la condition (3). Soit $i : Z \to X$ une immersion fermée,
où $Z$ est réduit et $|Z|$ irréductible. Posons $g = f \circ i : Z \to \Spec(A)$.
Soit $\mathcal{G}$ un module cohérent sur $Z$ dont le support schématique
est égal à $Z$, tel que $R^pg_*\mathcal{G}$ soit cohérent pour tout $p$.
Alors $\mathcal{F} = i_*\mathcal{G}$ est un module cohérent sur
$X$, de support schématique $Z$, tel que
$R^pf_*\mathcal{F} = R^pg_*\mathcal{G}$. Pour le voir, utilisons
la suite spectrale de Leray
(Cohomologie sur les sites, lemme \ref{sites-cohomology-lemma-relative-Leray})
et le fait que $R^qi_*\mathcal{G} = 0$ pour $q > 0$ d'après le
lemme \ref{lemma-affine-vanishing-higher-direct-images},
ainsi que le fait qu'une immersion fermée est affine
(Morphismes d'espaces, lemme
\ref{spaces-morphisms-lemma-closed-immersion-affine}).
On se ramène donc à trouver un faisceau cohérent $\mathcal{G}$ sur $Z$
de support égal à $Z$, tel que $R^pg_*\mathcal{G}$ soit cohérent
pour tout $p$.

\medskip\noindent
Appliquons le lemme \ref{lemma-weak-chow} au morphisme $Z \to \Spec(A)$.
On obtient ainsi un diagramme
$$
\xymatrix{
Z \ar[rd]_g & Z' \ar[d]^-{g'} \ar[l]^\pi \ar[r]_i & \mathbf{P}^n_A \ar[dl] \\
& \Spec(A) &
}
$$
où $\pi : Z' \to Z$ est propre et surjectif, et $i$ une immersion.
Puisque $Z \to \Spec(A)$ est propre, on conclut que $g'$ est propre
(Morphismes d'espaces, lemme \ref{spaces-morphisms-lemma-composition-proper}).
Par suite, $i$ est une immersion fermée
(Morphismes d'espaces, lemmes
\ref{spaces-morphisms-lemma-universally-closed-permanence} et
\ref{spaces-morphisms-lemma-immersion-when-closed}).
Il s'ensuit que le morphisme
$i' = (i, \pi) : \mathbf{P}^n_A \times_{\Spec(A)} Z' = \mathbf{P}^n_Z$ est
une immersion fermée
(Morphismes d'espaces, lemme \ref{spaces-morphisms-lemma-semi-diagonal}).
Posons
$$
\mathcal{L} =
i^*\mathcal{O}_{\mathbf{P}^n_A}(1) =
(i')^*\mathcal{O}_{\mathbf{P}^n_Z}(1)
$$
On peut appliquer le lemme \ref{lemma-kill-by-twisting}
à $\mathcal{L}$ et $\pi$, ainsi qu'à $\mathcal{L}$ et $g'$.
Ainsi, pour tout $d \gg 0$, on a
$R^p\pi_*\mathcal{L}^{\otimes d} = 0$ pour tout $p > 0$ et
$R^p(g')_*\mathcal{L}^{\otimes d} = 0$ pour tout $p > 0$.
Posons $\mathcal{G} = \pi_*\mathcal{L}^{\otimes d}$.
D'après la suite spectrale de Leray
(Cohomologie sur les sites, lemme \ref{sites-cohomology-lemma-relative-Leray}),
on a
$$
E_2^{p, q} = R^pg_* R^q\pi_*\mathcal{L}^{\otimes d}
\Rightarrow
R^{p + q}(g')_*\mathcal{L}^{\otimes d}
$$
et, par le choix de $d$, les seuls termes non nuls de $E_2^{p, q}$ sont
ceux pour lesquels $q = 0$, tandis que les seuls termes non nuls de
$R^{p + q}(g')_*\mathcal{L}^{\otimes d}$ sont ceux pour lesquels $p = q = 0$.
Il en résulte que $R^pg_*\mathcal{G} = 0$ pour $p > 0$ et
que $g_*\mathcal{G} = (g')_*\mathcal{L}^{\otimes d}$.
En appliquant
Cohomologie des schémas, lemme
\ref{coherent-lemma-locally-projective-pushforward}
on voit que $g_*\mathcal{G} = (g')_*\mathcal{L}^{\otimes d}$ est
cohérent.

\medskip\noindent
Il reste à vérifier que le support de $\mathcal{G}$ est $Z$.
Cela résulte du fait que $\mathcal{L}^{\otimes d}$ possède
beaucoup de sections globales. Donnons les détails.
Notons que $\mathcal{L}^{\otimes d}$ est engendré par ses sections globales pour tout $d \geq 0$,
car il en est de même de $\mathcal{O}_{\mathbf{P}^n}(d)$.
Choisissons un point $z \in Z'$ s'envoyant sur le point générique $\xi$ de $Z$ ;
c'est possible puisque $\pi$ est surjectif.
(Observons que $Z$ possède bien un point générique, car $|Z|$ est irréductible
et $Z$ est noethérien, donc quasi-séparé ; ainsi $|Z|$ est un espace topologique
sobre d'après
Propriétés des espaces, lemme
\ref{spaces-properties-lemma-quasi-separated-sober}.)
Choisissons $s \in \Gamma(Z', \mathcal{L}^{\otimes d})$ qui ne s'annule pas en $z$.
Puisque $\Gamma(Z, \mathcal{G}) = \Gamma(Z', \mathcal{L}^{\otimes d})$,
on peut considérer $s$ comme une section globale de $\mathcal{G}$.
Choisissons un point géométrique $\overline{z}$ de $Z'$ au-dessus de $z$
et notons $\overline{\xi} = g' \circ \overline{z}$
le point géométrique correspondant de $Z$. L'application d'adjonction
$$
(g')^*\mathcal{G} =
(g')^*g'_*\mathcal{L}^{\otimes d} \longrightarrow \mathcal{L}^{\otimes d}
$$
induit une application entre les fibres
$\mathcal{G}_{\overline{\xi}} \to \mathcal{L}_{\overline{z}}$, voir
Propriétés des espaces, lemme
\ref{spaces-properties-lemma-stalk-pullback-quasi-coherent}.
De plus, l'application d'adjonction envoie l'image réciproque de $s$ (vu comme section
de $\mathcal{G}$) sur $s$ (vu comme section de $\mathcal{L}^{\otimes d}$).
Ainsi, l'image de $s$ dans l'espace vectoriel source de la flèche
$$
\mathcal{G}_{\overline{\xi}} \otimes \kappa(\overline{\xi})
\longrightarrow
\mathcal{L}^{\otimes d}_{\overline{z}} \otimes \kappa(\overline{z})
$$
n'est pas nulle, puisque, par le choix de $s$, son image dans la cible de la flèche
est non nulle. Ainsi, $\xi$ appartient au support de $\mathcal{G}$
(Morphismes d'espaces, lemme \ref{spaces-morphisms-lemma-support-finite-type}).
Puisque $|Z|$ est irréductible et $Z$ réduit, on conclut que
le support schématique de $\mathcal{G}$ est $Z$ tout entier, comme voulu.
\end{proof}

\begin{lemma}
\label{lemma-proper-over-affine-cohomology-finite}
Soit $A$ un anneau noethérien.
Soit $f : X \to \Spec(A)$ un morphisme propre d'espaces algébriques.
Soit $\mathcal{F}$ un $\mathcal{O}_X$-module cohérent.
Alors $H^i(X, \mathcal{F})$ est un $A$-module de type fini pour tout $i \geq 0$.
\end{lemma}

\begin{proof}
Il s'agit simplement du cas affine du lemme \ref{lemma-proper-pushforward-coherent}.
En effet, le lemme \ref{lemma-higher-direct-image} montre que
$R^if_*\mathcal{F}$ est un faisceau quasi-cohérent. C'est donc le faisceau
quasi-cohérent associé au $A$-module
$\Gamma(\Spec(A), R^if_*\mathcal{F}) = H^i(X, \mathcal{F})$.
L'égalité résulte de
Cohomologie sur les sites, lemme \ref{sites-cohomology-lemma-apply-Leray},
et de l'annulation des groupes de cohomologie supérieurs des modules quasi-cohérents
sur les schémas affines (Cohomologie des schémas, lemme
\ref{coherent-lemma-quasi-coherent-affine-cohomology-zero}).
D'après le lemme \ref{lemma-coherent-Noetherian}, le faisceau $R^if_*\mathcal{F}$ est
cohérent si et seulement si $H^i(X, \mathcal{F})$
est un $A$-module de type fini. Le
lemme \ref{lemma-proper-pushforward-coherent} donne donc la conclusion.
\end{proof}

\begin{lemma}
\label{lemma-graded-finiteness}
Soit $A$ un anneau noethérien.
Soit $B$ une $A$-algèbre graduée de type fini.
Soit $f : X \to \Spec(A)$ un morphisme propre d'espaces algébriques.
Posons $\mathcal{B} = f^*\widetilde B$.
Soit $\mathcal{F}$ un $\mathcal{B}$-module gradué
quasi-cohérent de type fini.
Pour tout $p \geq 0$, le $B$-module gradué $H^p(X, \mathcal{F})$
est un $B$-module de type fini.
\end{lemma}

\begin{proof}
Pour le démontrer, considérons le diagramme de produit fibré
$$
\xymatrix{
X' = \Spec(B) \times_{\Spec(A)} X
\ar[r]_-\pi \ar[d]_{f'} &
X \ar[d]^f \\
\Spec(B) \ar[r] &
\Spec(A)
}
$$
Notons que $f'$ est un morphisme propre, voir
Morphismes d'espaces, lemme \ref{spaces-morphisms-lemma-base-change-proper}.
De plus, $B$ est une $A$-algèbre de type fini, et est donc
noethérienne (Algèbre, lemme \ref{algebra-lemma-Noetherian-permanence}).
Il en résulte que $X'$ est un espace algébrique noethérien
(Morphismes d'espaces, lemme
\ref{spaces-morphisms-lemma-finite-presentation-noetherian}).
Notons que $X'$ est le spectre relatif de la
$\mathcal{O}_X$-algèbre quasi-cohérente $\mathcal{B}$ d'après
Morphismes d'espaces, lemme
\ref{spaces-morphisms-lemma-affine-equivalence-algebras}.
Puisque $\mathcal{F}$ est un $\mathcal{B}$-module quasi-cohérent,
il existe un unique
$\mathcal{O}_{X'}$-module quasi-cohérent $\mathcal{F}'$ tel que
$\pi_*\mathcal{F}' = \mathcal{F}$, voir
Morphismes d'espaces, lemme
\ref{spaces-morphisms-lemma-affine-equivalence-modules}.
Puisque $\mathcal{F}$ est de type fini comme $\mathcal{B}$-module, on
conclut que $\mathcal{F}'$ est un
$\mathcal{O}_{X'}$-module de type fini (détails omis). Autrement dit,
$\mathcal{F}'$ est un $\mathcal{O}_{X'}$-module cohérent
(lemme \ref{lemma-coherent-Noetherian}).
Puisque le morphisme $\pi : X' \to X$ est affine, on a
$$
H^p(X, \mathcal{F}) = H^p(X', \mathcal{F}')
$$
d'après le
lemme \ref{lemma-affine-vanishing-higher-direct-images}
et
Cohomologie sur les sites, lemme \ref{sites-cohomology-lemma-apply-Leray}.
Le lemme résulte donc du
lemme \ref{lemma-proper-over-affine-cohomology-finite}.
\end{proof}









\section{Faisceaux inversibles amples et cohomologie}
\label{section-ample-cohomology}

\noindent
Voici un critère d'amplitude pour les espaces algébriques propres sur une base affine,
exprimé par l'annulation de la cohomologie après tensorisation.

\begin{lemma}
\label{lemma-vanshing-gives-ample}
Soit $R$ un anneau noethérien. Soit $X$ un espace algébrique propre
sur $R$. Soit $\mathcal{L}$ un $\mathcal{O}_X$-module inversible.
Les conditions suivantes sont équivalentes :
\begin{enumerate}
\item $X$ est un schéma et $\mathcal{L}$ est ample sur $X$ ;
\item pour tout $\mathcal{O}_X$-module cohérent $\mathcal{F}$, il existe
un $n_0 \geq 0$ tel que
$H^p(X, \mathcal{F} \otimes \mathcal{L}^{\otimes n}) = 0$ pour tout $n \geq n_0$
et tout $p > 0$ ; et
\item pour tout $\mathcal{O}_X$-module cohérent $\mathcal{F}$, il existe
un $n \geq 1$ tel que
$H^1(X, \mathcal{F} \otimes \mathcal{L}^{\otimes n}) = 0$.
\end{enumerate}
\end{lemma}

\begin{proof}
L'implication (1) $\Rightarrow$ (2) résulte de
Cohomologie des schémas, lemme \ref{coherent-lemma-vanshing-gives-ample}.
L'implication (2) $\Rightarrow$ (3) est triviale.
L'implication (3) $\Rightarrow$ (1) est le
lemme \ref{lemma-Noetherian-h1-zero-invertible}.
\end{proof}

\begin{lemma}
\label{lemma-surjective-finite-morphism-ample}
Soit $R$ un anneau noethérien. Soit $f : Y \to X$ un morphisme entre espaces
algébriques propres sur $R$. Soit $\mathcal{L}$ un
$\mathcal{O}_X$-module inversible. Supposons $f$ fini et surjectif.
Les conditions suivantes sont équivalentes :
\begin{enumerate}
\item $X$ est un schéma et $\mathcal{L}$ est ample ; et
\item $Y$ est un schéma et $f^*\mathcal{L}$ est ample.
\end{enumerate}
\end{lemma}

\begin{proof}
Supposons (1). Alors $Y$ est un schéma, car un morphisme fini est représentable
(par des schémas), voir Morphismes d'espaces, lemme
\ref{spaces-morphisms-lemma-integral-local}. Par suite, (2)
résulte de Cohomologie des schémas, lemme
\ref{coherent-lemma-surjective-finite-morphism-ample}.

\medskip\noindent
Supposons (2). Soit $P$ la propriété suivante des
$\mathcal{O}_X$-modules cohérents $\mathcal{F}$ : il existe $n_0$
tel que $H^p(X, \mathcal{F} \otimes \mathcal{L}^{\otimes n}) = 0$
pour tout $n \geq n_0$ et tout $p > 0$. Nous allons montrer que $P$ est vérifiée
par tout $\mathcal{O}_X$-module cohérent $\mathcal{F}$, ce qui entraîne que
$\mathcal{L}$ est ample d'après le lemme \ref{lemma-vanshing-gives-ample}.
Nous allons appliquer le lemme \ref{lemma-property-higher-rank-cohomological}.
Il faut donc vérifier les conditions (1), (2) et (3) de ce lemme pour $P$.
La propriété (1) résulte de la suite exacte longue de cohomologie associée
à une suite exacte courte de faisceaux et du fait que la tensorisation par
un faisceau inversible est un foncteur exact. La propriété (2) résulte de ce que
$H^p(X, -)$ est un foncteur additif.

\medskip\noindent
Pour vérifier (3), soit $i : Z \to X$ un sous-espace fermé réduit tel que $|Z|$
soit irréductible. Soient $i' : Z' \to Y$ et $f' : Z' \to Z$ comme dans le
lemme \ref{lemma-finite-morphism-Noetherian}, et posons
$\mathcal{G} = f'_*\mathcal{O}_{Z'}$. Nous affirmons que
$\mathcal{G}$ vérifie les propriétés (3)(a) et (3)(b) du
lemme \ref{lemma-property-higher-rank-cohomological},
ce qui achèvera la démonstration. La propriété (3)(a) a été établie au
lemme \ref{lemma-finite-morphism-Noetherian}. Pour voir (3)(b), soit
$\mathcal{I}$ un faisceau quasi-cohérent d'idéaux non nul sur $Z$.
Notons $\mathcal{I}' \subset \mathcal{O}_{Z'}$ l'idéal quasi-cohérent
$(f')^{-1}\mathcal{I} \mathcal{O}_{Z'}$, c'est-à-dire
l'image de $(f')^*\mathcal{I} \to \mathcal{O}_{Z'}$.
D'après le lemme \ref{lemma-affine-morphism-projection-ideal}, on a
$f_*\mathcal{I}' = \mathcal{I} \mathcal{G}$. Nous affirmons que la valeur commune
$\mathcal{G}' = \mathcal{I} \mathcal{G} = f'_*\mathcal{I}'$
vérifie la condition énoncée en (3)(b).
D'abord, il est clair que le support de $\mathcal{G}/\mathcal{G}'$
est contenu dans celui de $\mathcal{O}_Z/\mathcal{I}$,
qui est un sous-espace propre de $|Z|$, puisque $\mathcal{I}$ est un
faisceau d'idéaux non nul sur l'espace algébrique réduit et irréductible $Z$.
Rappelons que $f'_*$, $i_*$ et $i'_*$ transforment les modules cohérents en
modules cohérents, voir les lemmes \ref{lemma-finite-pushforward-coherent} et
\ref{lemma-i-star-equivalence}.
Puisque $Y$ est un schéma et que $\mathcal{L}$ est ample,
le lemme \ref{lemma-vanshing-gives-ample} montre
qu'il existe $n_0$ tel que
$$
H^p(Y, i'_*\mathcal{I}' \otimes_{\mathcal{O}_Y} f^*\mathcal{L}^{\otimes n}) = 0
$$
pour $n \geq n_0$ et $p > 0$. On obtient alors
\begin{align*}
H^p(X, i_*\mathcal{G}' \otimes_{\mathcal{O}_X} \mathcal{L}^{\otimes n})
& =
H^p(Z, \mathcal{G'} \otimes_{\mathcal{O}_Z} i^*\mathcal{L}^{\otimes n}) \\
& =
H^p(Z,
f'_*\mathcal{I}' \otimes_{\mathcal{O}_Z} i^*\mathcal{L}^{\otimes n})) \\
& =
H^p(Z, f'_*(\mathcal{I}' \otimes_{\mathcal{O}_{Z'}}
(f')^*i^*\mathcal{L}^{\otimes n})) \\
& =
H^p(Z, f'_*(\mathcal{I}' \otimes_{\mathcal{O}_{Z'}}
(i')^*f^*\mathcal{L}^{\otimes n})) \\
& =
H^p(Z',
\mathcal{I}' \otimes_{\mathcal{O}_{Z'}} (i')^*f^*\mathcal{L}^{\otimes n})) \\
& =
H^p(Y, i'_*\mathcal{I}' \otimes_{\mathcal{O}_Y} f^*\mathcal{L}^{\otimes n}) = 0
\end{align*}
Nous avons utilisé ici la formule de projection et la suite spectrale de Leray
(voir Cohomologie sur les sites, sections
\ref{sites-cohomology-section-projection-formula} et
\ref{sites-cohomology-section-leray})
ainsi que le lemme \ref{lemma-finite-higher-direct-image-zero}.
Cela vérifie la propriété (3)(b) du
lemme \ref{lemma-property-higher-rank-cohomological}, comme voulu.
\end{proof}










\section{Le théorème des fonctions formelles}
\label{section-theorem-formal-functions}

\noindent
Cette section est l'analogue de
Cohomologie des schémas, section \ref{coherent-section-theorem-formal-functions}.
Nous invitons le lecteur à lire d'abord cette section.

\begin{situation}
\label{situation-formal-functions}
Ici, $A$ est un anneau noethérien et $I \subset A$ un idéal.
De plus, $f : X \to \Spec(A)$ est un morphisme propre d'espaces algébriques
et $\mathcal{F}$ un faisceau cohérent sur $X$.
\end{situation}

\noindent
Dans cette situation, on note $I^n\mathcal{F}$ le sous-module quasi-cohérent
de $\mathcal{F}$ engendré, comme $\mathcal{O}_X$-module,
par les produits de sections locales de $\mathcal{F}$ et d'éléments de $I^n$.
Autrement dit, c'est l'image de l'application
$f^*\widetilde{I} \otimes_{\mathcal{O}_X} \mathcal{F} \to \mathcal{F}$.

\begin{lemma}
\label{lemma-cohomology-powers-ideal-times-F}
Dans la situation \ref{situation-formal-functions}.
Posons $B = \bigoplus_{n \geq 0} I^n$.
Alors, pour tout $p \geq 0$, le $B$-module gradué
$\bigoplus_{n \geq 0} H^p(X, I^n\mathcal{F})$ est
un $B$-module de type fini.
\end{lemma}

\begin{proof}
Posons $\mathcal{B} = \bigoplus I^n\mathcal{O}_X = f^*\widetilde{B}$.
Alors $\bigoplus I^n\mathcal{F}$ est un $\mathcal{B}$-module
gradué de type fini. Le résultat découle donc
du lemme \ref{lemma-graded-finiteness}.
\end{proof}

\begin{lemma}
\label{lemma-cohomology-powers-ideal-application}
Dans la situation \ref{situation-formal-functions}.
Pour tout $p \geq 0$, il existe un entier $c \geq 0$ tel que
\begin{enumerate}
\item l'application de multiplication
$I^{n - c} \otimes H^p(X, I^c\mathcal{F}) \to H^p(X, I^n\mathcal{F})$
soit surjective pour tout $n \geq c$ ; et
\item l'image de $H^p(X, I^{n + m}\mathcal{F}) \to H^p(X, I^n\mathcal{F})$
soit contenue dans le sous-module $I^{m - c} H^p(X, I^n\mathcal{F})$
pour tout $n \geq 0$ et tout $m \geq c$.
\end{enumerate}
\end{lemma}

\begin{proof}
D'après le lemme \ref{lemma-cohomology-powers-ideal-times-F},
on peut trouver $d_1, \ldots, d_t \geq 0$ et
$x_i \in H^p(X, I^{d_i}\mathcal{F})$ tels que
$\bigoplus_{n \geq 0} H^p(X, I^n\mathcal{F})$ soit engendré
par $x_1, \ldots, x_t$ sur $B = \bigoplus_{n \geq 0} I^n$.
Posons $c = \max\{d_i\}$. Il est clair que (1) est vérifiée.
Pour (2), posons $b = \max(0, n - c)$.
Considérons le diagramme commutatif de $A$-modules
$$
\xymatrix{
I^{n + m - c - b} \otimes I^b \otimes
H^p(X, I^c\mathcal{F}) \ar[r] \ar[d] &
I^{n + m - c} \otimes H^p(X, I^c\mathcal{F}) \ar[r] &
H^p(X, I^{n + m}\mathcal{F}) \ar[d] \\
I^{n + m - c - b} \otimes H^p(X, I^n\mathcal{F}) \ar[rr] & &
H^p(X, I^n\mathcal{F})
}
$$
D'après la partie (1) du lemme, le composé des flèches horizontales
est surjectif si $n + m \geq c$. D'autre part, il est clair
que $n + m - c - b \geq m - c$. D'où la partie (2).
\end{proof}

\begin{lemma}
\label{lemma-ML-cohomology-powers-ideal}
Dans la situation \ref{situation-formal-functions}.
Fixons $p \geq 0$.
\begin{enumerate}
\item Il existe $c_1 \geq 0$ tel que, pour tout $n \geq c_1$,
on ait
$$
\Ker(
H^p(X, \mathcal{F}) \to H^p(X, \mathcal{F}/I^n\mathcal{F})
)
\subset
I^{n - c_1}H^p(X, \mathcal{F}).
$$
\item Le système projectif
$$
\left(H^p(X, \mathcal{F}/I^n\mathcal{F})\right)_{n \in \mathbf{N}}
$$
vérifie la condition de Mittag-Leffler (voir
Homologie, définition \ref{homology-definition-Mittag-Leffler}).
\item En fait, quels que soient $p$ et $n$, il existe $c_2(n) \geq n$
tel que
$$
\Im(H^p(X, \mathcal{F}/I^k\mathcal{F})
\to H^p(X, \mathcal{F}/I^n\mathcal{F}))
=
\Im(H^p(X, \mathcal{F})
\to H^p(X, \mathcal{F}/I^n\mathcal{F}))
$$
pour tout $k \geq c_2(n)$.
\end{enumerate}
\end{lemma}

\begin{proof}
Posons $c_1 = \max\{c_p, c_{p + 1}\}$, où $c_p, c_{p +1}$ sont les entiers
obtenus dans le lemme \ref{lemma-cohomology-powers-ideal-application} pour
$H^p$ et $H^{p + 1}$. Nous utiliserons cette constante dans les démonstrations de
(1), (2) et (3).

\medskip\noindent
Démontrons la partie (1). Considérons la suite exacte courte
$$
0 \to I^n\mathcal{F} \to \mathcal{F} \to \mathcal{F}/I^n\mathcal{F} \to 0
$$
La suite exacte longue de cohomologie donne
$$
\Ker(
H^p(X, \mathcal{F}) \to H^p(X, \mathcal{F}/I^n\mathcal{F})
)
=
\Im(
H^p(X, I^n\mathcal{F}) \to H^p(X, \mathcal{F})
)
$$
D'après le choix de $c_1$, ce sous-module est donc contenu dans
$I^{n - c_1}H^p(X, \mathcal{F})$ pour $n \geq c_1$.

\medskip\noindent
Notons que la partie (3) entraîne la partie (2) par définition de la condition
de Mittag-Leffler.

\medskip\noindent
Démontrons la partie (3).
Fixons $n$ pour le reste de la démonstration.
Considérons le diagramme commutatif
$$
\xymatrix{
0 \ar[r] &
I^n\mathcal{F} \ar[r] &
\mathcal{F} \ar[r] &
\mathcal{F}/I^n\mathcal{F} \ar[r] &
0 \\
0 \ar[r] &
I^{n + m}\mathcal{F} \ar[r] \ar[u] &
\mathcal{F} \ar[r] \ar[u] &
\mathcal{F}/I^{n + m}\mathcal{F} \ar[r] \ar[u] &
0
}
$$
Il induit le diagramme commutatif suivant
$$
\xymatrix{
H^p(X, I^n\mathcal{F}) \ar[r] &
H^p(X, \mathcal{F}) \ar[r] &
H^p(X, \mathcal{F}/I^n\mathcal{F}) \ar[r]_\delta &
H^{p + 1}(X, I^n\mathcal{F}) \\
H^p(X, I^{n + m}\mathcal{F}) \ar[r] \ar[u] &
H^p(X, \mathcal{F}) \ar[r] \ar[u]^1 &
H^p(X, \mathcal{F}/I^{n + m}\mathcal{F}) \ar[r] \ar[u] &
H^{p + 1}(X, I^{n + m}\mathcal{F}) \ar[u]^a
}
$$
Si $m \geq c_1$, on voit que l'image de $a$ est
contenue dans $I^{m - c_1} H^{p + 1}(X, I^n\mathcal{F})$.
D'après le lemme d'Artin--Rees (voir Algèbre, lemme \ref{algebra-lemma-map-AR}),
il existe un entier $c_3(n)$ tel que
$$
I^N H^{p + 1}(X, I^n\mathcal{F}) \cap \Im(\delta)
\subset
\delta\left(I^{N - c_3(n)}H^p(X, \mathcal{F}/I^n\mathcal{F})\right)
$$
pour tout $N \geq c_3(n)$. Comme $H^p(X, \mathcal{F}/I^n\mathcal{F})$
est annulé par $I^n$, on voit que, si $m \geq c_3(n) + c_1 + n$,
alors
$$
\Im(H^p(X, \mathcal{F}/I^{n + m}\mathcal{F})
\to H^p(X, \mathcal{F}/I^n\mathcal{F}))
=
\Im(H^p(X, \mathcal{F})
\to H^p(X, \mathcal{F}/I^n\mathcal{F}))
$$
Autrement dit, la partie (3) est vérifiée avec $c_2(n) = c_3(n) + c_1 + n$.
\end{proof}

\begin{theorem}[Théorème des fonctions formelles]
\label{theorem-formal-functions}
Dans la situation \ref{situation-formal-functions}. Fixons $p \geq 0$.
Le système d'applications
$$
H^p(X, \mathcal{F})/I^nH^p(X, \mathcal{F})
\longrightarrow
H^p(X, \mathcal{F}/I^n\mathcal{F})
$$
définit un isomorphisme de limites
$$
H^p(X, \mathcal{F})^\wedge
\longrightarrow
\lim_n H^p(X, \mathcal{F}/I^n\mathcal{F})
$$
où le membre de gauche est le complété du $A$-module
$H^p(X, \mathcal{F})$ relativement à l'idéal $I$, voir
Algèbre, section \ref{algebra-section-completion}.
De plus, c'est un homéomorphisme pour les topologies de limite projective.
\end{theorem}

\begin{proof}
Cela résulte en fait immédiatement du
lemme \ref{lemma-ML-cohomology-powers-ideal}. Donnons les détails.
Posons $M = H^p(X, \mathcal{F})$ et $M_n = H^p(X, \mathcal{F}/I^n\mathcal{F})$.
Notons $N_n = \Im(M \to M_n)$.
D'après la description de la limite dans Homologie, section
\ref{homology-section-inverse-systems}, on a
$$
\lim_n M_n
=
\{(x_n) \in \prod M_n \mid \varphi_i(x_n) = x_{n - 1}, \ n = 2, 3, \ldots\}
$$
Choisissons un élément $x = (x_n) \in \lim_n M_n$.
D'après la partie (3) du lemme \ref{lemma-ML-cohomology-powers-ideal},
on a $x_n \in N_n$ pour tout $n$, puisque, par
définition, $x_n$ est l'image d'un certain $x_{n + m} \in M_{n + m}$ pour
tout $m$. D'après la partie (1) du lemme \ref{lemma-ML-cohomology-powers-ideal},
il existe une factorisation
$$
M \to N_n \to M/I^{n - c_1}M
$$
de l'application de réduction. Notons $y_n \in M/I^{n - c_1}M$ l'image de $x_n$
pour $n \geq c_1$. Puisque, pour $n' \geq n$, le composé
$M \to M_{n'} \to M_n$ est l'application donnée $M \to M_n$, on voit que
$y_{n'}$ s'envoie sur $y_n$ par l'application canonique
$M/I^{n' - c_1}M \to M/I^{n - c_1}M$. Ainsi, $y = (y_{n + c_1})$
définit un élément de $\lim_n M/I^nM$.
Nous omettons de vérifier que $y$ s'envoie sur $x$ par
l'application
$$
M^\wedge = \lim_n M/I^nM \longrightarrow \lim_n M_n
$$
du lemme. Nous omettons également la vérification concernant les topologies.
\end{proof}

\begin{lemma}
\label{lemma-spell-out-theorem-formal-functions}
Soit $A$ un anneau. Soit $I \subset A$ un idéal. Supposons $A$
noethérien et complet pour la topologie $I$-adique.
Soit $f : X \to \Spec(A)$ un morphisme propre d'espaces algébriques.
Soit $\mathcal{F}$ un faisceau cohérent sur $X$. Alors
$$
H^p(X, \mathcal{F}) = \lim_n H^p(X, \mathcal{F}/I^n\mathcal{F})
$$
pour tout $p \geq 0$.
\end{lemma}

\begin{proof}
C'est une reformulation du théorème des fonctions formelles
(théorème \ref{theorem-formal-functions}) dans le
cas d'un anneau de base noethérien complet. En effet, dans ce cas, le
$A$-module $H^p(X, \mathcal{F})$ est de type fini
(lemme \ref{lemma-proper-over-affine-cohomology-finite}), donc
complet pour la topologie $I$-adique (Algèbre, lemme \ref{algebra-lemma-completion-tensor}),
et l'on voit qu'il est inutile de compléter le membre de gauche.
\end{proof}

\begin{lemma}
\label{lemma-formal-functions-stalk}
Soit $S$ un schéma. Soit $f : X \to Y$ un morphisme d'espaces algébriques
au-dessus de $S$, et soit $\mathcal{F}$ un faisceau quasi-cohérent sur $X$. Supposons que
\begin{enumerate}
\item $Y$ soit localement noethérien,
\item $f$ soit propre, et
\item $\mathcal{F}$ soit cohérent.
\end{enumerate}
Soit $\overline{y}$ un point géométrique de $Y$.
Considérons les « voisinages infinitésimaux »
$$
\xymatrix{
X_n =
\Spec(\mathcal{O}_{Y, \overline{y}}/\mathfrak m_{\overline{y}}^n) \times_Y X
\ar[r]_-{i_n} \ar[d]_{f_n} &
X \ar[d]^f \\
\Spec(\mathcal{O}_{Y, \overline{y}}/\mathfrak m_{\overline{y}}^n)
\ar[r]^-{c_n} & Y
}
$$
de la fibre $X_1 = X_{\overline{y}}$, et posons
$\mathcal{F}_n = i_n^*\mathcal{F}$. On a alors
$$
\left(R^pf_*\mathcal{F}\right)_{\overline{y}}^\wedge
\cong
\lim_n H^p(X_n, \mathcal{F}_n)
$$
comme $\mathcal{O}_{Y, \overline{y}}^\wedge$-modules.
\end{lemma}

\begin{proof}
Il s'agit simplement d'une reformulation d'un cas particulier du théorème
des fonctions formelles, théorème \ref{theorem-formal-functions}.
Donnons les détails. Notons que $\mathcal{O}_{Y, \overline{y}}$
est un anneau local noethérien, voir
Propriétés des espaces, lemme
\ref{spaces-properties-lemma-Noetherian-local-ring-Noetherian}.
Considérons le morphisme canonique
$c : \Spec(\mathcal{O}_{Y, \overline{y}}) \to Y$.
C'est un morphisme plat, car il identifie les anneaux locaux.
Notons $f' : X' \to \Spec(\mathcal{O}_{Y, \overline{y}})$
le changement de base de $f$ à cet anneau local. On a
$c^*R^pf_*\mathcal{F} = R^pf'_*\mathcal{F}'$ d'après le
lemme \ref{lemma-flat-base-change-cohomology}.
De plus, on dispose d'identifications canoniques $X_n = X'_n$
pour tout $n \geq 1$.

\medskip\noindent
On peut donc supposer que $Y = \Spec(A)$ est le spectre d'un
anneau local noethérien strictement hensélien $A$, d'idéal maximal
$\mathfrak m$, et que $\overline{y} \to Y$ est égal à
$\Spec(A/\mathfrak m) \to Y$. Il en résulte que
$$
\left(R^pf_*\mathcal{F}\right)_{\overline{y}} =
\Gamma(Y, R^pf_*\mathcal{F}) =
H^p(X, \mathcal{F})
$$
car $(Y, \overline{y})$ est un objet initial de la catégorie
des voisinages \'etales de $\overline{y}$.
Les morphismes $c_n$ sont tous des immersions fermées.
Leurs changements de base $i_n$ sont donc eux aussi des immersions fermées.
Notons que $i_{n, *}\mathcal{F}_n = i_{n, *}i_n^*\mathcal{F}
= \mathcal{F}/\mathfrak m^n\mathcal{F}$. La suite spectrale de Leray
pour $i_n$ et le lemme \ref{lemma-finite-pushforward-coherent} donnent
$$
H^p(X_n, \mathcal{F}_n) =
H^p(X, i_{n, *}\mathcal{F}) =
H^p(X, \mathcal{F}/\mathfrak m^n\mathcal{F})
$$
On peut donc bien appliquer le théorème des fonctions formelles pour calculer
la limite figurant dans l'énoncé du lemme, ce qui conclut.
\end{proof}

\noindent
Voici un lemme que nous généraliserons ensuite aux fibres de
dimension $ > 0$, dans le lemme suivant.

\begin{lemma}
\label{lemma-higher-direct-images-zero-finite-fibre}
Soit $S$ un schéma. Soit $f : X \to Y$ un morphisme d'espaces algébriques
au-dessus de $S$. Soit $\overline{y}$ un point géométrique de $Y$.
Supposons que
\begin{enumerate}
\item $Y$ soit localement noethérien,
\item $f$ soit propre, et
\item l'espace topologique sous-jacent à $X_{\overline{y}}$ soit discret.
\end{enumerate}
Alors, pour tout faisceau cohérent $\mathcal{F}$ sur $X$, on a
$(R^pf_*\mathcal{F})_{\overline{y}} = 0$ pour tout $p > 0$.
\end{lemma}

\begin{proof}
Soit $\kappa(\overline{y})$ le corps résiduel de l'anneau
local $\mathcal{O}_{Y, \overline{y}}$. Comme dans le
lemme \ref{lemma-formal-functions-stalk},
posons $X_{\overline{y}} = X_1 = \Spec(\kappa(\overline{y})) \times_Y X$.
D'après Morphismes d'espaces, lemme
\ref{spaces-morphisms-lemma-quasi-finite-at-point}
le morphisme $f : X \to Y$ est quasi-fini en chacun
des points de la fibre de $X \to Y$ au-dessus de $\overline{y}$.
Il en résulte que $X_{\overline{y}} \to \overline{y}$ est séparé et
quasi-fini. Par suite, $X_{\overline{y}}$ est un schéma
d'après Morphismes d'espaces, proposition
\ref{spaces-morphisms-proposition-locally-quasi-finite-separated-over-scheme}.
Comme il est quasi-compact, son espace topologique sous-jacent est un espace
discret fini. C'est donc un schéma affine d'après
Schémas, lemme \ref{schemes-lemma-scheme-finite-discrete-affine}.
Le lemme \ref{lemma-image-affine-finite-morphism-affine-Noetherian}
montre alors que les espaces algébriques $X_n$ sont eux aussi des schémas affines.
De plus, l'espace topologique sous-jacent à chaque $X_n$ est le même
que celui de $X_1$. Il s'ensuit que $H^p(X_n, \mathcal{F}_n) = 0$
pour tout $p > 0$. Ainsi,
$(R^pf_*\mathcal{F})_{\overline{y}}^\wedge = 0$
d'après le lemme \ref{lemma-formal-functions-stalk}.
Notons que $R^pf_*\mathcal{F}$ est cohérent d'après le
lemme \ref{lemma-proper-pushforward-coherent} ;
ainsi $R^pf_*\mathcal{F}_{\overline{y}}$ est un
$\mathcal{O}_{Y, \overline{y}}$-module de type fini.
D'après Algèbre, lemme \ref{algebra-lemma-completion-tensor},
cela implique que $(R^pf_*\mathcal{F})_{\overline{y}} = 0$.
\end{proof}

\begin{lemma}
\label{lemma-higher-direct-images-zero-above-dimension-fibre}
\begin{slogan}
Pour les morphismes propres, les fibres des images directes supérieures sont nulles aux degrés
strictement supérieurs à la dimension de la fibre.
\end{slogan}
Soit $S$ un schéma.
Soit $f : X \to Y$ un morphisme d'espaces algébriques au-dessus de $S$.
Soit $\overline{y}$ un point géométrique de $Y$.
Supposons que
\begin{enumerate}
\item $Y$ soit localement noethérien,
\item $f$ soit propre, et
\item $\dim(X_{\overline{y}}) = d$.
\end{enumerate}
Alors, pour tout faisceau cohérent $\mathcal{F}$ sur $X$, on a
$(R^pf_*\mathcal{F})_{\overline{y}} = 0$ pour tout $p > d$.
\end{lemma}

\begin{proof}
Soit $\kappa(\overline{y})$ le corps résiduel de l'anneau
local $\mathcal{O}_{Y, \overline{y}}$. Comme dans le
lemme \ref{lemma-formal-functions-stalk},
posons $X_{\overline{y}} = X_1 = \Spec(\kappa(\overline{y})) \times_Y X$.
De plus, l'espace topologique sous-jacent à chaque voisinage
infinitésimal $X_n$ est le même que celui de $X_{\overline{y}}$.
Ainsi, $H^p(X_n, \mathcal{F}_n) = 0$ pour tout $p > d$ d'après le
lemme \ref{lemma-vanishing-above-dimension}.
Ainsi, on voit que $(R^pf_*\mathcal{F})_{\overline{y}}^\wedge = 0$
d'après le lemme \ref{lemma-formal-functions-stalk} pour $p > d$.
Notons que $R^pf_*\mathcal{F}$ est cohérent d'après le
lemme \ref{lemma-proper-pushforward-coherent} ;
ainsi $R^pf_*\mathcal{F}_{\overline{y}}$ est un
$\mathcal{O}_{Y, \overline{y}}$-module de type fini.
D'après Algèbre, lemme \ref{algebra-lemma-completion-tensor},
cela implique que $(R^pf_*\mathcal{F})_{\overline{y}} = 0$.
\end{proof}






\section{Applications du théorème des fonctions formelles}
\label{section-applications-formal-functions}

\noindent
Nous compléterons cette section selon les besoins.

\begin{lemma}
\label{lemma-characterize-finite}
(Pour une version plus générale, voir
Compléments sur les morphismes d'espaces, lemme
\ref{spaces-more-morphisms-lemma-characterize-finite}).
Soit $S$ un schéma.
Soit $f : X \to Y$ un morphisme d'espaces algébriques au-dessus de $S$.
Supposons $Y$ localement noethérien.
Les conditions suivantes sont équivalentes :
\begin{enumerate}
\item $f$ est fini ; et
\item $f$ est propre et $|X_k|$ est un espace discret
pour tout morphisme $\Spec(k) \to Y$, où $k$ est un corps.
\end{enumerate}
\end{lemma}

\begin{proof}
Un morphisme fini est propre d'après
Morphismes d'espaces, lemme \ref{spaces-morphisms-lemma-finite-proper}.
Un morphisme fini est quasi-fini d'après
Morphismes d'espaces, lemme \ref{spaces-morphisms-lemma-finite-quasi-finite}.
Un morphisme quasi-fini a des fibres discrètes $X_k$, voir
Morphismes d'espaces, lemme \ref{spaces-morphisms-lemma-locally-quasi-finite}.
Ainsi, un morphisme fini est propre et a des fibres discrètes $X_k$.

\medskip\noindent
Supposons $f$ propre, à fibres discrètes $X_k$.
Nous voulons montrer que $f$ est fini.
Il suffit en fait de montrer que $f$ est affine.
En effet, si $f$ est affine, il s'ensuit que
$f$ est entier d'après
Morphismes d'espaces, lemme
\ref{spaces-morphisms-lemma-integral-universally-closed}
puis, d'après
Morphismes d'espaces, lemme \ref{spaces-morphisms-lemma-finite-integral},
que $f$ est fini.

\medskip\noindent
Pour montrer que $f$ est affine, on peut supposer $Y$ affine ; notre
but est de montrer que $X$ est également affine.
Puisque $f$ est propre, $X$ est séparé et quasi-compact.
Nous allons montrer que, pour tout
$\mathcal{O}_X$-module cohérent $\mathcal{F}$, on a $H^1(X, \mathcal{F}) = 0$.
Cela implique que $H^1(X, \mathcal{F}) = 0$ pour tout
$\mathcal{O}_X$-module quasi-cohérent $\mathcal{F}$ d'après les
lemmes \ref{lemma-directed-colimit-coherent} et \ref{lemma-colimits}.
On en déduit que $X$ est affine d'après la
proposition \ref{proposition-vanishing-affine}. D'après le
lemme \ref{lemma-higher-direct-images-zero-finite-fibre},
on conclut que les fibres de $R^1f_*\mathcal{F}$
sont nulles en tout point géométrique de $Y$.
Autrement dit, $R^1f_*\mathcal{F} = 0$. La suite spectrale de Leray
pour $f$ donne donc
$H^1(X , \mathcal{F}) = H^1(Y, f_*\mathcal{F})$.
Puisque $Y$ est affine et $f_*\mathcal{F}$ quasi-cohérent
(Morphismes d'espaces, lemme \ref{spaces-morphisms-lemma-pushforward}),
on conclut que $H^1(Y, f_*\mathcal{F}) = 0$
d'après
Cohomologie des schémas, lemme
\ref{coherent-lemma-quasi-coherent-affine-cohomology-zero}.
Ainsi, $H^1(X, \mathcal{F}) = 0$, comme voulu.
\end{proof}

\noindent
On en déduit le résultat utile suivant.

\begin{lemma}
\label{lemma-proper-finite-fibre-finite-in-neighbourhood}
(Pour une version plus générale, voir
Compléments sur les morphismes d'espaces, lemme
\ref{spaces-more-morphisms-lemma-proper-finite-fibre-finite-in-neighbourhood}).
Soit $S$ un schéma. Soit $f : X \to Y$ un morphisme d'espaces algébriques
au-dessus de $S$. Soit $\overline{y}$ un point géométrique de $Y$.
Supposons que
\begin{enumerate}
\item $Y$ soit localement noethérien,
\item $f$ soit propre, et
\item $|X_{\overline{y}}|$ soit fini.
\end{enumerate}
Alors il existe un voisinage ouvert $V \subset Y$ de $\overline{y}$
tel que $f|_{f^{-1}(V)} : f^{-1}(V) \to V$ soit fini.
\end{lemma}

\begin{proof}
Le morphisme $f$ est quasi-fini en tous les points géométriques de $X$
au-dessus de $\overline{y}$ d'après
Morphismes d'espaces, lemme \ref{spaces-morphisms-lemma-quasi-finite-at-point}.
D'après Morphismes d'espaces, lemme
\ref{spaces-morphisms-lemma-locally-finite-type-quasi-finite-part}, l'ensemble
des points où $f$ est quasi-fini est un sous-espace ouvert $U \subset X$.
Posons $Z = X \setminus U$. Alors $\overline{y} \not \in f(Z)$. Puisque $f$
est propre, l'ensemble $f(Z) \subset Y$ est fermé. Choisissons un voisinage ouvert
$V \subset Y$ de $\overline{y}$ tel que $Z \cap V = \emptyset$. Alors
$f^{-1}(V) \to V$ est localement quasi-fini et propre.
Ainsi, $f^{-1}(V) \to V$ a des fibres discrètes $X_k$
(Morphismes d'espaces, lemme
\ref{spaces-morphisms-lemma-locally-quasi-finite})
qui sont quasi-compactes, donc finies.
Par conséquent, $f^{-1}(V) \to V$
est fini d'après le lemme \ref{lemma-characterize-finite}.
\end{proof}




\begin{multicols}{2}[\section{Autres chapitres}]
\noindent
Préliminaires
\begin{enumerate}
\item \hyperref[introduction-section-phantom]{Introduction}
\item \hyperref[conventions-section-phantom]{Conventions}
\item \hyperref[sets-section-phantom]{Théorie des ensembles}
\item \hyperref[categories-section-phantom]{Catégories}
\item \hyperref[topology-section-phantom]{Topologie}
\item \hyperref[sheaves-section-phantom]{Faisceaux sur les espaces}
\item \hyperref[sites-section-phantom]{Sites et faisceaux}
\item \hyperref[stacks-section-phantom]{Champs}
\item \hyperref[fields-section-phantom]{Corps}
\item \hyperref[algebra-section-phantom]{Algèbre commutative}
\item \hyperref[brauer-section-phantom]{Groupes de Brauer}
\item \hyperref[homology-section-phantom]{Algèbre homologique}
\item \hyperref[derived-section-phantom]{Catégories dérivées}
\item \hyperref[simplicial-section-phantom]{Méthodes simpliciales}
\item \hyperref[more-algebra-section-phantom]{Compléments d'algèbre}
\item \hyperref[smoothing-section-phantom]{Lissage des morphismes d'anneaux}
\item \hyperref[modules-section-phantom]{Faisceaux de modules}
\item \hyperref[sites-modules-section-phantom]{Modules sur les sites}
\item \hyperref[injectives-section-phantom]{Objets injectifs}
\item \hyperref[cohomology-section-phantom]{Cohomologie des faisceaux}
\item \hyperref[sites-cohomology-section-phantom]{Cohomologie sur les sites}
\item \hyperref[dga-section-phantom]{Algèbre différentielle graduée}
\item \hyperref[dpa-section-phantom]{Algèbre à puissances divisées}
\item \hyperref[sdga-section-phantom]{Faisceaux différentiels gradués}
\item \hyperref[hypercovering-section-phantom]{Hyperrecouvrements}
\end{enumerate}
Schémas
\begin{enumerate}
\setcounter{enumi}{25}
\item \hyperref[schemes-section-phantom]{Schémas}
\item \hyperref[constructions-section-phantom]{Constructions de schémas}
\item \hyperref[properties-section-phantom]{Propriétés des schémas}
\item \hyperref[morphisms-section-phantom]{Morphismes de schémas}
\item \hyperref[coherent-section-phantom]{Cohomologie des schémas}
\item \hyperref[divisors-section-phantom]{Diviseurs}
\item \hyperref[limits-section-phantom]{Limites de schémas}
\item \hyperref[varieties-section-phantom]{Variétés}
\item \hyperref[topologies-section-phantom]{Topologies sur les schémas}
\item \hyperref[descent-section-phantom]{Descente}
\item \hyperref[perfect-section-phantom]{Catégories dérivées de schémas}
\item \hyperref[more-morphisms-section-phantom]{Compléments sur les morphismes}
\item \hyperref[flat-section-phantom]{Compléments sur la platitude}
\item \hyperref[groupoids-section-phantom]{Schémas en groupoïdes}
\item \hyperref[more-groupoids-section-phantom]{Compléments sur les schémas en groupoïdes}
\item \hyperref[etale-section-phantom]{Morphismes étales de schémas}
\end{enumerate}
Thèmes de théorie des schémas
\begin{enumerate}
\setcounter{enumi}{41}
\item \hyperref[chow-section-phantom]{Homologie de Chow}
\item \hyperref[intersection-section-phantom]{Théorie de l'intersection}
\item \hyperref[pic-section-phantom]{Schémas de Picard des courbes}
\item \hyperref[weil-section-phantom]{Théories cohomologiques de Weil}
\item \hyperref[adequate-section-phantom]{Modules adéquats}
\item \hyperref[dualizing-section-phantom]{Complexes dualisants}
\item \hyperref[duality-section-phantom]{Dualité pour les schémas}
\item \hyperref[discriminant-section-phantom]{Discriminants et différentes}
\item \hyperref[derham-section-phantom]{Cohomologie de de Rham}
\item \hyperref[local-cohomology-section-phantom]{Cohomologie locale}
\item \hyperref[algebraization-section-phantom]{Géométrie algébrique et formelle}
\item \hyperref[curves-section-phantom]{Courbes algébriques}
\item \hyperref[resolve-section-phantom]{Résolution des surfaces}
\item \hyperref[models-section-phantom]{Réduction semi-stable}
\item \hyperref[functors-section-phantom]{Foncteurs et morphismes}
\item \hyperref[equiv-section-phantom]{Catégories dérivées de variétés}
\item \hyperref[pione-section-phantom]{Groupes fondamentaux de schémas}
\item \hyperref[etale-cohomology-section-phantom]{Cohomologie étale}
\item \hyperref[crystalline-section-phantom]{Cohomologie cristalline}
\item \hyperref[proetale-section-phantom]{Cohomologie pro-étale}
\item \hyperref[relative-cycles-section-phantom]{Cycles relatifs}
\item \hyperref[more-etale-section-phantom]{Compléments de cohomologie étale}
\item \hyperref[trace-section-phantom]{La formule des traces}
\end{enumerate}
Espaces algébriques
\begin{enumerate}
\setcounter{enumi}{64}
\item \hyperref[spaces-section-phantom]{Espaces algébriques}
\item \hyperref[spaces-properties-section-phantom]{Propriétés des espaces algébriques}
\item \hyperref[spaces-morphisms-section-phantom]{Morphismes d'espaces algébriques}
\item \hyperref[decent-spaces-section-phantom]{Espaces algébriques décents}
\item \hyperref[spaces-cohomology-section-phantom]{Cohomologie des espaces algébriques}
\item \hyperref[spaces-limits-section-phantom]{Limites d'espaces algébriques}
\item \hyperref[spaces-divisors-section-phantom]{Diviseurs sur les espaces algébriques}
\item \hyperref[spaces-over-fields-section-phantom]{Espaces algébriques sur des corps}
\item \hyperref[spaces-topologies-section-phantom]{Topologies sur les espaces algébriques}
\item \hyperref[spaces-descent-section-phantom]{Descente et espaces algébriques}
\item \hyperref[spaces-perfect-section-phantom]{Catégories dérivées d'espaces}
\item \hyperref[spaces-more-morphisms-section-phantom]{Compléments sur les morphismes d'espaces}
\item \hyperref[spaces-flat-section-phantom]{Platitude sur les espaces algébriques}
\item \hyperref[spaces-groupoids-section-phantom]{Groupoïdes dans les espaces algébriques}
\item \hyperref[spaces-more-groupoids-section-phantom]{Compléments sur les groupoïdes dans les espaces}
\item \hyperref[bootstrap-section-phantom]{Amorçage}
\item \hyperref[spaces-pushouts-section-phantom]{Sommes amalgamées d'espaces algébriques}
\end{enumerate}
Thèmes de géométrie
\begin{enumerate}
\setcounter{enumi}{81}
\item \hyperref[spaces-chow-section-phantom]{Groupes de Chow des espaces}
\item \hyperref[groupoids-quotients-section-phantom]{Quotients de groupoïdes}
\item \hyperref[spaces-more-cohomology-section-phantom]{Compléments sur la cohomologie des espaces}
\item \hyperref[spaces-simplicial-section-phantom]{Espaces simpliciaux}
\item \hyperref[spaces-duality-section-phantom]{Dualité pour les espaces}
\item \hyperref[formal-spaces-section-phantom]{Espaces algébriques formels}
\item \hyperref[restricted-section-phantom]{Algébrisation des espaces formels}
\item \hyperref[spaces-resolve-section-phantom]{Résolution des surfaces revisitée}
\end{enumerate}
Théorie des déformations
\begin{enumerate}
\setcounter{enumi}{89}
\item \hyperref[formal-defos-section-phantom]{Théorie formelle des déformations}
\item \hyperref[defos-section-phantom]{Théorie des déformations}
\item \hyperref[cotangent-section-phantom]{Le complexe cotangent}
\item \hyperref[examples-defos-section-phantom]{Problèmes de déformation}
\end{enumerate}
Champs algébriques
\begin{enumerate}
\setcounter{enumi}{93}
\item \hyperref[algebraic-section-phantom]{Champs algébriques}
\item \hyperref[examples-stacks-section-phantom]{Exemples de champs}
\item \hyperref[stacks-sheaves-section-phantom]{Faisceaux sur les champs algébriques}
\item \hyperref[criteria-section-phantom]{Critères de représentabilité}
\item \hyperref[artin-section-phantom]{Axiomes d'Artin}
\item \hyperref[quot-section-phantom]{Espaces de Quot et de Hilbert}
\item \hyperref[stacks-properties-section-phantom]{Propriétés des champs algébriques}
\item \hyperref[stacks-morphisms-section-phantom]{Morphismes de champs algébriques}
\item \hyperref[stacks-limits-section-phantom]{Limites de champs algébriques}
\item \hyperref[stacks-cohomology-section-phantom]{Cohomologie des champs algébriques}
\item \hyperref[stacks-perfect-section-phantom]{Catégories dérivées de champs}
\item \hyperref[stacks-introduction-section-phantom]{Introduction aux champs algébriques}
\item \hyperref[stacks-more-morphisms-section-phantom]{Compléments sur les morphismes de champs}
\item \hyperref[stacks-geometry-section-phantom]{La géométrie des champs}
\end{enumerate}
Thèmes de théorie des espaces de modules
\begin{enumerate}
\setcounter{enumi}{107}
\item \hyperref[moduli-section-phantom]{Champs de modules}
\item \hyperref[moduli-curves-section-phantom]{Espaces de modules de courbes}
\end{enumerate}
Divers
\begin{enumerate}
\setcounter{enumi}{109}
\item \hyperref[examples-section-phantom]{Exemples}
\item \hyperref[exercises-section-phantom]{Exercices}
\item \hyperref[guide-section-phantom]{Guide bibliographique}
\item \hyperref[desirables-section-phantom]{Desiderata}
\item \hyperref[coding-section-phantom]{Style de programmation}
\item \hyperref[obsolete-section-phantom]{Obsolète}
\item \hyperref[fdl-section-phantom]{Licence de documentation libre GNU}
\item \hyperref[index-section-phantom]{Index généré automatiquement}
\end{enumerate}
\end{multicols}


\bibliography{my}
\bibliographystyle{amsalpha}

\end{document}
